\documentclass[11pt]{article}
\usepackage[margin=2.4cm]{geometry}
\usepackage{amsmath,amssymb,amsthm,booktabs,xcolor}
\usepackage{tikz}
\usepackage{quantikz}
\usepackage{listings}
\usepackage[colorlinks,linkcolor=blue!55!black,citecolor=blue!55!black,urlcolor=blue!55!black]{hyperref}

\renewcommand{\ket}[1]{\lvert #1 \rangle}
\renewcommand{\bra}[1]{\langle #1 \rvert}
\newcommand{\MAJ}{\mathrm{MAJ}}
\newcommand{\UMA}{\mathrm{UMA}}
\newcommand{\ADD}{\mathrm{ADD}}
\newcommand{\SUB}{\mathrm{SUB}}
\newcommand{\CMULT}{\mathrm{CMULT}}
\newcommand{\TOF}{\mathrm{TOF}}
\newcommand{\CNOT}{\mathrm{CNOT}}
\newcommand{\id}{\mathbb{1}}
\newcommand{\code}[1]{\texttt{#1}}

\theoremstyle{plain}
\newtheorem{theorem}{Theorem}[section]
\newtheorem{lemma}[theorem]{Lemma}
\newtheorem{proposition}[theorem]{Proposition}
\theoremstyle{remark}
\newtheorem*{remark}{Remark}

\definecolor{boxblue}{RGB}{22,63,110}
\definecolor{boxred}{RGB}{140,38,28}
\usepackage[most]{tcolorbox}
\newtcolorbox{keybox}[1][]{colback=boxblue!4,colframe=boxblue,coltitle=white,
  fonttitle=\bfseries,boxrule=0.7pt,left=6pt,right=6pt,top=5pt,bottom=5pt,#1}
\newtcolorbox{warnbox}[1][]{colback=boxred!4,colframe=boxred,coltitle=white,
  fonttitle=\bfseries,boxrule=0.7pt,left=6pt,right=6pt,top=5pt,bottom=5pt,#1}

\definecolor{codebg}{RGB}{248,248,251}
\definecolor{codekw}{RGB}{0,80,150}
\definecolor{codestr}{RGB}{146,58,22}
\definecolor{codecom}{RGB}{96,110,96}
\lstdefinestyle{py}{
  language=Python,
  basicstyle=\small\ttfamily,
  keywordstyle=\color{codekw}\bfseries,
  commentstyle=\color{codecom}\itshape,
  stringstyle=\color{codestr},
  showstringspaces=false,
  backgroundcolor=\color{codebg},
  frame=single, framerule=0.3pt, rulecolor=\color{black!25},
  aboveskip=8pt, belowskip=6pt,
  columns=fullflexible, keepspaces=true, breaklines=true, tabsize=4,
}
\lstset{style=py}

\usepackage{listings}
\definecolor{codebg}{RGB}{248,248,251}
\definecolor{codekw}{RGB}{0,80,150}
\definecolor{codestr}{RGB}{146,58,22}
\definecolor{codecom}{RGB}{96,110,96}
\lstdefinestyle{py}{
  language=Python, basicstyle=\small\ttfamily,
  keywordstyle=\color{codekw}\bfseries, commentstyle=\color{codecom}\itshape,
  stringstyle=\color{codestr}, showstringspaces=false,
  backgroundcolor=\color{codebg},
  frame=single, framerule=0.3pt, rulecolor=\color{black!25},
  aboveskip=7pt, belowskip=5pt,
  columns=fullflexible, keepspaces=true, breaklines=true, tabsize=4,
}
\lstset{style=py}
\title{\bfseries Reversible modular arithmetic in Shor's algorithm\\[4pt]
\large What it must do, how it is built, and how it runs --- with verified Qiskit code}
\author{}
\date{}

\begin{document}
\maketitle
\vspace{-2.2em}

\begin{keybox}[title={What this is, and how to read it}]
Shor's circuit is three layers: Hadamards, a ladder of controlled modular multiplications, and an
inverse QFT. The first and third are drawn gate by gate everywhere. The middle one is almost always
a box labelled ``$\times a^{2^j}\bmod N$'' --- and it is over $99\%$ of the circuit. This note
opens that box twice: once as mathematics, once as running code.

\medskip
\begin{tabular}{@{}ll@{}}
\textbf{Parts I--III} & why the arithmetic is what it is: what phase estimation demands, why\\
                      & garbage is fatal, the tower from Toffoli up, and how a pile of\\
                      & permutations turns into a measurable phase.\\[4pt]
\textbf{Part IV}      & a complete, minimal, \emph{verified} implementation --- about $150$ lines\\
                      & of Qiskit, with the tests that make it believable.\\[4pt]
\textbf{Part V}       & making it fast: ripple-carry arithmetic, windowing, unary-iteration\\
                      & lookups, measurement-based uncomputation, nested windows.\\[4pt]
\textbf{Part VI}      & reading somebody else's implementation, and what its bugs teach.\\
\end{tabular}

\medskip
Every listing in Parts~IV--V is copied from a file that was executed, and every number quoted as
measured was measured. The implementation factors $15$, $21$ and $33$ end to end; $21$ and $33$
are the honest tests, because their orders do not divide $2^t$ and so force the
continued-fraction step to do real work.

\medskip
Notation: $N$ is the modulus, $n=\lceil\log_2N\rceil$, $a$ the base with $\gcd(a,N)=1$, $r$ its
order, $t=2n$ the counting register. Qubit $k$ of a register holds bit $k$ of its value, LSB first
(Qiskit's convention).
\end{keybox}

\vspace{4pt}
{\small\setlength{\parskip}{0pt}\tableofcontents}
\part*{Part I --- Orientation}
\addcontentsline{toc}{part}{Part I --- Orientation}

\section{Where the arithmetic sits}

\subsection{What phase estimation asks for, and what it costs}

Phase estimation \cite{kitaev,cemm} is a general procedure. Give it a unitary $U$, an eigenvector
$\ket u$ with $U\ket u=e^{2\pi i\varphi}\ket u$, and $t$ blank qubits, and it writes $t$ bits of
$\varphi$ into those qubits. It is entirely indifferent to what $U$ is.

\begin{center}
\begin{quantikz}[column sep=10pt,row sep=10pt]
\lstick{$\ket0$} & \gate{H} & \ctrl{4} & \qw & \qw & \ \ldots\ & \qw & \gate[4]{F^\dagger_{2^t}} & \meter{}\\
\lstick{$\ket0$} & \gate{H} & \qw & \ctrl{3} & \qw & & \qw & & \meter{}\\
\lstick{$\vdots$} & & & & & & & & \\
\lstick{$\ket0$} & \gate{H} & \qw & \qw & \qw & \ \ldots\ & \ctrl{1} & & \meter{}\\
\lstick{$\ket u$} & \qw & \gate{U^{2^0}} & \gate{U^{2^1}} & \qw & \ \ldots\ & \gate{U^{2^{t-1}}} & \qw & \qw
\end{quantikz}
\end{center}

But look at what it charges you. The rungs are $U^{1},U^{2},U^{4},\dots,U^{2^{t-1}}$, so building
them by repetition costs
\[
1+2+4+\cdots+2^{t-1}=2^t-1
\]
applications of $U$ --- exponentially many. \textbf{Phase estimation is efficient only for unitaries
whose powers are cheap.} For a generic $U$ --- a black box, a Hamiltonian evolution $e^{iHt}$ ---
there is no shortcut, and the procedure is useless. This is worth saying plainly because the usual
presentation, which introduces QPE abstractly and then ``applies'' it to $U_a$, makes the
application look routine. It is not routine; it is the whole trick.

Modular multiplication is special twice over:
\begin{description}
\item[Gift 1: the powers collapse.] $U_a^{2^j}=U_{a^{2^j}\bmod N}$. A high power of $U_a$ is not a
long circuit --- it is \emph{the same circuit with a different classical constant}, and that
constant is computed classically by repeated squaring before the machine is switched on. This
turns $2^t-1$ applications into $t$.
\item[Gift 2: one rung is cheap.] $U_a$ itself has a circuit of size $O(n^2)$. This is Part~II.
\end{description}

\noindent
Both are needed, and neither implies the other. Gift~1 without Gift~2 would be a short ladder of
unaffordable rungs. Gift~2 without Gift~1 would be an affordable rung repeated an unaffordable
number of times. Shor's algorithm \cite{shor} exists in the narrow space where arithmetic hands you
both.

\subsection{The four demands on $U_a$}

Phase estimation does not accept just any circuit that computes $ay \bmod N$. Four properties are
load-bearing, and each one is discharged at a specific level of the tower.

\begin{center}
\begin{tabular}{@{}p{0.30\textwidth}p{0.40\textwidth}l@{}}
\toprule
Demand & Why & Discharged at\\
\midrule
\textbf{Unitary} & QPE's whole framework is spectral; a non-unitary map has no eigenphase to read.
& Level 4 (\S\ref{sec:l4})\\[3pt]
\textbf{Cheap powers} & otherwise the ladder is exponentially long (\S1.1). & Level 5 (\S\ref{sec:l5})\\[3pt]
\textbf{Controllable} & an \emph{uncontrolled} $U$ contributes a global phase, which no measurement
can see; the control is what makes the phase relative. & Level 4 (\S\ref{sec:l4})\\[3pt]
\textbf{Garbage-free} & retained scratch entangles with the counting register and destroys the
interference the inverse QFT depends on. & Levels 1--2 (\S\ref{sec:l1}, \S\ref{sec:l2})\\
\bottomrule
\end{tabular}
\end{center}

\medskip\noindent
\textbf{Unitarity is a number-theoretic condition.} The map $y\mapsto ay\bmod N$ is injective on
$\{0,\dots,N{-}1\}$ precisely when $a$ is invertible mod $N$, i.e.\ when $\gcd(a,N)=1$. If that
fails, the map collapses distinct inputs together, is not a permutation, and is not unitary --- so
there is no operator, no spectrum, and nothing for QPE to estimate. (In the algorithm this is
checked classically first, and if it fails you have already factored $N$ by computing the gcd.) The
same condition returns at Level~4 in a second role: $\bar a=a^{-1}\bmod N$ is what makes the
multiplier realisable \emph{in place}. One hypothesis, two jobs.

\subsection{One circuit, two descriptions}\label{sec:dual}

Here is the conceptual hinge of the whole subject, and the reason a note about \emph{classical}
arithmetic belongs inside a \emph{quantum} algorithm.

The ladder is a single unitary. Written in the computational basis it says
\begin{equation}\label{eq:compbasis}
\ket x\ket y\ \longmapsto\ \ket x\,\ket{a^xy\bmod N}
\qquad\text{--- it tabulates the function }f(x)=a^x\bmod N,
\end{equation}
which is how Shor described it in 1994 \cite{shor}: build a periodic function, Fourier-sample its
period. Written in the eigenbasis of $U_a$ the very same unitary says
\begin{equation}\label{eq:eigbasis}
\ket x\ket{u_s}\ \longmapsto\ e^{2\pi i xs/r}\ket x\ket{u_s}
\qquad\text{--- it kicks an eigenphase onto the control register,}
\end{equation}
which is Kitaev's \cite{kitaev} and CEMM's \cite{cemm} reading: phase estimation on $U_a$. The two
are related by the change of basis $\ket1=r^{-1/2}\sum_s\ket{u_s}$, and they are descriptions of
one circuit, not two algorithms.

\begin{keybox}[title={Why this matters for the arithmetic}]
Reading \eqref{eq:compbasis}, the ladder does nothing but arithmetic --- and indeed everything in
Part~II is a permutation matrix with no complex number anywhere in it. Reading
\eqref{eq:eigbasis}, the same gates manufacture the phases that the algorithm ultimately measures.

Nothing was added in between. The phases are not written by any gate; they appear because a
classical permutation was applied to a \emph{superposition} rather than to a number, and then
interfered. That is the single thing the classical circuit --- from which every gate in Part~II is
borrowed --- could not do.
\end{keybox}

\subsection{Why garbage is fatal, not untidy}\label{sec:garbage}

This is the demand that most shapes the construction, so it deserves the argument rather than the
slogan. Suppose some ancilla ends in a state $\ket{g_x}$ that depends on which branch of the
superposition it is in. Then instead of the clean state the algorithm wants, you have
\[
\frac1{\sqrt{2^t}}\sum_x\ket x\,\ket{a^x\bmod N}\,\ket{g_x}.
\]
The inverse QFT works by making branches $x$ and $x'$ interfere. Their interference term carries a
factor $\braket{g_{x'}}{g_x}$. If the garbage records anything about $x$ --- and for a useful
computation it always does --- those states are orthogonal, the factor is $0$, and there is no
interference at all. The measurement returns uniform noise. The circuit is still perfectly unitary;
it simply computes nothing.

\begin{warnbox}[title={The consequence for every block below}]
You cannot erase scratch by overwriting it (that is not reversible) and you cannot discard it (it
stays entangled). The only way out is to \emph{run a computation backwards} and unwrite it. Hence:
\begin{itemize}\setlength\itemsep{1pt}
\item Level~1's adder ends with a reversed chain of $\UMA$ blocks that erase the carries;
\item Level~2's modular adder spends three of its seven blocks recomputing and clearing one flag
qubit;
\item Level~4's multiplier runs a second multiplication \emph{in reverse} purely to restore a
scratch register to $\ket0$.
\end{itemize}
Roughly half of every circuit in Part~II exists to undo the other half. That is not waste --- it is
the price of the interference, and \S\ref{sec:l1} shows the undoing can be made to write the answer
on its way out.
\end{warnbox}

\subsection{The zoom}\label{sec:zoom}

Every box at one level is a whole circuit at the level below. Reading downwards is the plan of the
document; the right-hand column is where the $n^3$ comes from. The last column names the method of
the \code{qiskit-shor} package \cite{repo} that realises the level --- Part~IV reads those methods
in full.

\begin{center}\small\setlength{\tabcolsep}{4pt}
\begin{tabular}{@{}rllrl@{}}
\toprule
Level & Object & Made of & Multiplier & In code (\S\ref{sec:pkg})\\
\midrule
5 & the ladder & $t$ rungs, one per counting qubit & $\times\,2n$ & \code{exponentiate\_modulo}\\
4 & $c\text{-}U_{a^{2^j}}$ & $2\times\CMULT$ and $n$ Fredkins & & \code{c\_multiply\_modulo}\\
3 & $\CMULT(k)$ & doubly-controlled modular adders & $\times\,n$ & \code{c\_add\_quantum\_modulo}\\
2 & $\ADD^{\bmod N}(k)$ & constant additions and sign tests & $\times\,5$ & \code{c\_add\_classical\_modulo}\\
1 & $\ADD(k)$ & $\MAJ$ chain, CNOT, $\UMA$ chain & $\times\,2n$ & \code{add\_classical}\\
0 & Toffoli & Clifford$+T$ & $7T$; $4T$ if uncomputed & \code{MAJGate}, \code{UMAGate}\\
\bottomrule
\end{tabular}
\end{center}

\noindent
Three of those steps each multiply by roughly $n$ --- rungs, bits of the multiplicand, bits per
addition --- and that, and nothing subtler, is why the algorithm is cubic.

\section{How to read these circuits}

The single most misleading thing about circuit diagrams: \textbf{the horizontal lines are not wires
and nothing flows along them.} There is no signal travelling rightwards. Each line is \emph{one
qubit}, and moving right means moving \emph{forward in time}. The diagram is closer to sheet music
than to a plumbing schematic: one staff line per qubit, read left to right.

So a diagram with eight lines has eight qubits and always will. Lines never split, merge, or get
created. Whatever is written at the far left is the input; whatever is at the far right is the
output. \emph{Read those two edges first} --- the middle is only the recipe for getting from one to
the other.

\begin{center}
\begin{tabular}{@{}ll@{}}
\toprule
Symbol & Meaning\\
\midrule
filled dot $\bullet$ & a \emph{control}: look at this qubit, change nothing\\
$\oplus$ & a \emph{target}: flip this qubit, if every dot on its vertical line is $1$\\
box across several lines & one operation on those qubits jointly (its innards are drawn elsewhere)\\
vertical line & not a wire --- visual glue saying ``these symbols are simultaneous''\\
line labelled $\ket{x}_{n+1}$ & a \emph{bundle}: $n{+}1$ qubits drawn as one line to keep the picture small\\
\bottomrule
\end{tabular}
\end{center}

\noindent
Everything in one vertical column happens at the same instant. When a control sits on a bundle, the
caption must tell you \emph{which} qubit of the bundle is the control --- in the modular adder of
\S\ref{sec:l2} it is always the top one, the sign qubit.

\part*{Part II --- The tower}
\addcontentsline{toc}{part}{Part II --- The tower}

\section{Level 0: the gate set}

\textbf{The conceptual point of this level.} Every gate in Part~II is a \emph{permutation matrix}:
one $1$ per row and per column, everything else $0$. Such a matrix is real and orthogonal, hence
unitary, and it maps computational basis states to computational basis states. In other words the
entire tower is a classical reversible circuit that happens to act linearly on superpositions.
\textbf{No complex number appears anywhere below Level~5.} Keep \S\ref{sec:dual} in mind: this is
exactly why the phases of Shor's algorithm cannot be attributed to any gate.

\begin{equation*}
X=\begin{pmatrix}0&1\\1&0\end{pmatrix},
\qquad
\CNOT=\begin{pmatrix}1&0&0&0\\0&1&0&0\\0&0&0&1\\0&0&1&0\end{pmatrix},
\end{equation*}
\begin{equation*}
\TOF=\left(\begin{smallmatrix}
1&0&0&0&0&0&0&0\\ 0&1&0&0&0&0&0&0\\ 0&0&1&0&0&0&0&0\\ 0&0&0&1&0&0&0&0\\
0&0&0&0&1&0&0&0\\ 0&0&0&0&0&1&0&0\\ 0&0&0&0&0&0&0&1\\ 0&0&0&0&0&0&1&0
\end{smallmatrix}\right),
\qquad
\mathrm{FRED}=\left(\begin{smallmatrix}
1&0&0&0&0&0&0&0\\ 0&1&0&0&0&0&0&0\\ 0&0&1&0&0&0&0&0\\ 0&0&0&1&0&0&0&0\\
0&0&0&0&1&0&0&0\\ 0&0&0&0&0&0&1&0\\ 0&0&0&0&0&1&0&0\\ 0&0&0&0&0&0&0&1
\end{smallmatrix}\right)
\end{equation*}

\[
X\ket a=\ket{a\oplus1},\qquad
\CNOT\ket{c,t}=\ket{c,\,t\oplus c},\qquad
\TOF\ket{c_1,c_2,t}=\ket{c_1,c_2,\,t\oplus c_1c_2},
\]
\[
\mathrm{FRED}\ket{c,x,y}=\ket{c,\,x\oplus c(x\oplus y),\,y\oplus c(x\oplus y)}
\quad\text{(swap $x,y$ iff $c=1$).}
\]

\medskip\noindent
\begin{minipage}[t]{0.52\textwidth}
\begin{center}
\begin{quantikz}[column sep=13pt,row sep=13pt]
\lstick{$\ket{c_1}$} & \ctrl{1} & \qw & \rstick{$\ket{c_1}$}\\
\lstick{$\ket{c_2}$} & \ctrl{1} & \qw & \rstick{$\ket{c_2}$}\\
\lstick{$\ket{t}$}   & \targ{}  & \qw & \rstick{$\ket{t\oplus c_1c_2}$}
\end{quantikz}
\end{center}
\end{minipage}
\hfill
\begin{minipage}[t]{0.44\textwidth}
\begin{center}
\begin{quantikz}[column sep=13pt,row sep=13pt]
\lstick{$\ket{c}$} & \ctrl{1} & \qw & \rstick{$\ket{c}$}\\
\lstick{$\ket{x}$} & \swap{1} & \qw & \rstick{$\ket{x'}$}\\
\lstick{$\ket{y}$} & \targX{} & \qw & \rstick{$\ket{y'}$}
\end{quantikz}
\end{center}
\end{minipage}

\medskip
The Toffoli is the \emph{only} nonlinear gate. CNOT implements $\mathbb F_2$-linear maps and is
therefore free in the fault-tolerant cost model; the Toffoli must be synthesised, and its standard
Clifford$+T$ decomposition costs $7$ $T$ gates and $6$ CNOTs:

\begin{center}
\begin{quantikz}[column sep=5.5pt,row sep=11pt]
\lstick{$\ket{c_1}$} & \qw & \qw & \qw & \ctrl{2} & \qw & \qw & \qw & \ctrl{2} & \qw & \qw & \qw & \ctrl{1} & \gate{T} & \ctrl{1} & \qw\\
\lstick{$\ket{c_2}$} & \qw & \ctrl{1} & \qw & \qw & \qw & \ctrl{1} & \qw & \qw & \gate{T} & \qw & \qw & \targ{} & \gate{T^\dagger} & \targ{} & \qw\\
\lstick{$\ket{t}$}   & \gate{H} & \targ{} & \gate{T^\dagger} & \targ{} & \gate{T} & \targ{} & \gate{T^\dagger} & \targ{} & \gate{T} & \gate{H} & \qw & \qw & \qw & \qw & \qw
\end{quantikz}
\end{center}

\noindent
This is why gate counts in this subject are quoted in \emph{Toffolis} rather than in gates: the
Toffoli count is a proxy for the $T$ count, and the $T$ count is what magic-state distillation
actually bills you for.

\begin{remark}[the halving that matters]
A Toffoli that is \emph{later uncomputed} --- a \emph{temporary logical-AND} \cite{gidney-and} ---
costs only $4$ $T$ for the compute/uncompute \emph{pair}, because the uncomputation can be done by
an $X$-basis measurement of the ancilla and a classically conditioned $CZ$, which is Clifford. By
\S\ref{sec:garbage} nearly every Toffoli below is of exactly this kind, so this is not a corner
case --- it is the common case.
\end{remark}

\section{Level 1: addition}\label{sec:l1}

\textbf{The problem this level solves.} Adding two $n$-bit numbers is column-by-column addition:
add the two digits and whatever carried in, write a digit, carry out. The carries are the
difficulty, because column $i{+}1$ cannot start until column $i$ has finished. Everything expensive
about Shor's algorithm traces back to that chain.

\subsection{The recurrences}
With $c_0=0$,
\begin{equation}\label{eq:schoolbook}
s_i=a_i\oplus b_i\oplus c_i,
\qquad
c_{i+1}=\MAJ(a_i,b_i,c_i)=a_ib_i\oplus a_ic_i\oplus b_ic_i,
\qquad
a+b=\sum_i s_i2^i+c_n2^n .
\end{equation}
Read this as a cost statement. The sum bits are $\mathbb F_2$-\emph{linear} given the carries, so
they are CNOTs and therefore free. The carries are \emph{quadratic}, so each one costs a Toffoli.
And they form a chain, so the depth is $\Theta(n)$. Linear parts free, quadratic parts expensive,
chain gives depth: that one line prices the algorithm.

\subsection{The two blocks}
On an ordered wire triple $(c,b,a)$ define, following Cuccaro, Draper, Kutin and Moulton
\cite{cuccaro},
\begin{align*}
\MAJ&:\ \CNOT(a\!\to\!b),\ \CNOT(a\!\to\!c),\ \TOF(c,b\!\to\!a);\\
\UMA&:\ \TOF(c,b\!\to\!a),\ \CNOT(a\!\to\!c),\ \CNOT(c\!\to\!b).
\end{align*}

\begin{center}
\begin{quantikz}[column sep=11pt,row sep=11pt]
\lstick{$\ket c$} & \qw & \targ{} & \ctrl{2} & \qw & \rstick{$\ket{c\oplus a}$}\\
\lstick{$\ket b$} & \targ{} & \qw & \ctrl{1} & \qw & \rstick{$\ket{b\oplus a}$}\\
\lstick{$\ket a$} & \ctrl{-1} & \ctrl{-2} & \targ{} & \qw & \rstick{$\ket{\MAJ(a,b,c)}$}
\end{quantikz}
\hspace{2.2em}
\begin{quantikz}[column sep=11pt,row sep=11pt]
\lstick{$\ket{c\oplus a}$} & \ctrl{2} & \targ{} & \ctrl{1} & \qw & \rstick{$\ket{c}$}\\
\lstick{$\ket{b\oplus a}$} & \ctrl{1} & \qw & \targ{} & \qw & \rstick{$\ket{a\oplus b\oplus c}$}\\
\lstick{$\ket{\MAJ}$} & \targ{} & \ctrl{-2} & \qw & \qw & \rstick{$\ket{a}$}
\end{quantikz}
\end{center}

\begin{lemma}[MAJ/UMA are partial inverses]\label{lem:majuma}
\[
\MAJ:(c,b,a)\mapsto\big(c\oplus a,\ b\oplus a,\ \MAJ(a,b,c)\big),
\qquad
\UMA:\big(c\oplus a,\ b\oplus a,\ \MAJ(a,b,c)\big)\mapsto\big(c,\ a\oplus b\oplus c,\ a\big).
\]
\end{lemma}
\begin{proof}
For $\MAJ$ the only nontrivial step is the Toffoli, which sends
$a\mapsto a\oplus(c\oplus a)(b\oplus a)=a\oplus\big(cb\oplus ca\oplus ab\oplus a\big)=ab\oplus ac\oplus bc$,
using $a^2=a$ and $a\oplus a=0$ over $\mathbb F_2$. For $\UMA$ the same computation read backwards
gives $\MAJ(a,b,c)\oplus(c\oplus a)(b\oplus a)=a$; then $c\oplus a\mapsto(c\oplus a)\oplus a=c$, and
$b\oplus a\mapsto(b\oplus a)\oplus c$, which is the sum bit $s=a\oplus b\oplus c$.
\end{proof}

\begin{keybox}[title={The idea, in one sentence}]
Look at what the second half of Lemma~\ref{lem:majuma} says: \textbf{$\UMA$ erases the carry and
writes the sum bit in the same three gates.}

By \S\ref{sec:garbage} the carries \emph{must} be erased, so a reverse pass is compulsory. The
design makes that compulsory pass do the useful work. The paper analogy is exact: when you add by
hand you scribble little carry digits above the columns, and you are left with a page full of
scribbles. Here you go back over the columns rubbing out each scribble, and the act of rubbing out
deposits the answer.
\end{keybox}

\subsection{The matrices}
Ordering the basis $\ket{c\,b\,a}$ as $0,\dots,7$, the blocks are the permutations
$\MAJ=(0)(1\,6\,7)(2)(3\,5)(4)$ and $\UMA=(0)(1\,7\,4\,6\,3\,5)(2)$:
\[
\MAJ=\begin{pmatrix}
1&0&0&0&0&0&0&0\\ 0&0&0&0&0&0&0&1\\ 0&0&1&0&0&0&0&0\\ 0&0&0&0&0&1&0&0\\
0&0&0&0&1&0&0&0\\ 0&0&0&1&0&0&0&0\\ 0&1&0&0&0&0&0&0\\ 0&0&0&0&0&0&1&0
\end{pmatrix},
\qquad
\UMA=\begin{pmatrix}
1&0&0&0&0&0&0&0\\ 0&0&0&0&0&1&0&0\\ 0&0&1&0&0&0&0&0\\ 0&0&0&0&0&0&1&0\\
0&0&0&0&0&0&0&1\\ 0&0&0&1&0&0&0&0\\ 0&0&0&0&1&0&0&0\\ 0&1&0&0&0&0&0&0
\end{pmatrix}
\]
(rows indexed by output, columns by input). Both are orthogonal, and
$\UMA\cdot\MAJ:(c,b,a)\mapsto(c,\,a\oplus b\oplus c,\,a)$ --- the permutation $0,3,2,1,6,5,4,7$,
which is Lemma~\ref{lem:majuma}.

\subsection{The CDKM ripple-carry adder}

\begin{keybox}[title={Why the wires look scrambled}]
The wire order is $c_0,b_0,a_0,b_1,a_1,\dots,z$ --- alternating, not ``all the $a$'s then all the
$b$'s''. This is the feature that makes the diagram look impenetrable, and it means \emph{nothing
mathematically}. It is a seating plan. Each block must touch three particular qubits; interleave
them this way and those three are always \emph{adjacent}, so the blocks are compact rather than a
tangle of long-range connections. Re-order the wires and the circuit is identical, only uglier.
\end{keybox}

\begin{center}
\begin{quantikz}[column sep=7pt,row sep=8pt]
\lstick{$\ket{c_0{=}0}$} & \gate[3]{\MAJ} & \qw & \qw & \qw & \qw & \qw & \gate[3]{\UMA} & \rstick{$\ket 0$}\\
\lstick{$\ket{b_0}$} & & \qw & \qw & \qw & \qw & \qw & & \rstick{$\ket{s_0}$}\\
\lstick{$\ket{a_0}$} & & \gate[3]{\MAJ} & \qw & \qw & \qw & \gate[3]{\UMA} & & \rstick{$\ket{a_0}$}\\
\lstick{$\ket{b_1}$} & \qw & & \qw & \qw & \qw & & \qw & \rstick{$\ket{s_1}$}\\
\lstick{$\ket{a_1}$} & \qw & & \gate[3]{\MAJ} & \qw & \gate[3]{\UMA} & & \qw & \rstick{$\ket{a_1}$}\\
\lstick{$\ket{b_2}$} & \qw & \qw & & \qw & & \qw & \qw & \rstick{$\ket{s_2}$}\\
\lstick{$\ket{a_2}$} & \qw & \qw & & \ctrl{1} & & \qw & \qw & \rstick{$\ket{a_2}$}\\
\lstick{$\ket{z{=}0}$} & \qw & \qw & \qw & \targ{} & \qw & \qw & \qw & \rstick{$\ket{s_3}$}
\end{quantikz}
\end{center}

\begin{theorem}[CDKM adder \cite{cuccaro}]\label{thm:cdkm}
The circuit above implements
\[
\ket{c_0{=}0}\,\ket b\,\ket a\,\ket{z}\ \longmapsto\ \ket{0}\,\ket{(a+b)\bmod 2^n}\,\ket a\,\ket{z\oplus c_n}
\]
using $2n$ Toffoli gates, $4n{+}1$ CNOTs and $2n{+}2$ qubits, with the operand $a$ preserved,
addition in place on $b$, and \textbf{zero garbage}. Fault-tolerantly it costs $4n$ $T$ gates, each
carry being a temporary AND \cite{gidney-and}.
\end{theorem}

\subsection*{A run, in full: $a=5$, $b=3$}

\begin{center}\small
\begin{tabular}{@{}lcccccccc@{}}
\toprule
after & $c_0$ & $b_0$ & $a_0$ & $b_1$ & $a_1$ & $b_2$ & $a_2$ & $z$\\
\midrule
start           & 0 & \textbf{1} & 1 & \textbf{1} & 0 & \textbf{0} & 1 & 0\\
$\MAJ(c_0,b_0,a_0)$ & 1 & 0 & \textbf{1} & 1 & 0 & 0 & 1 & 0\\
$\MAJ(a_0,b_1,a_1)$ & 1 & 0 & 1 & 1 & \textbf{1} & 0 & 1 & 0\\
$\MAJ(a_1,b_2,a_2)$ & 1 & 0 & 1 & 1 & 0 & 1 & \textbf{1} & 0\\
CNOT $a_2\!\to\!z$  & 1 & 0 & 1 & 1 & 0 & 1 & 1 & \textbf{1}\\
$\UMA(a_1,b_2,a_2)$ & 1 & 0 & 1 & 1 & 1 & \textbf{0} & 1 & 1\\
$\UMA(a_0,b_1,a_1)$ & 1 & 0 & 1 & \textbf{0} & 0 & 0 & 1 & 1\\
$\UMA(c_0,b_0,a_0)$ & \textbf{0} & \textbf{0} & 1 & 0 & 0 & 0 & 1 & 1\\
\bottomrule
\end{tabular}
\end{center}

\noindent
Read the $b$ column downwards: it starts as $011=3$ and ends as $000$, with $z=1$ --- that is
$1000_2=8$. The $a$ wires end at $1,0,1=5$, exactly as they started, and $c_0$ is back to $0$.
Nothing is left dirty, which is what Theorem~\ref{thm:cdkm} promises and \S\ref{sec:garbage}
demands.

\smallskip\noindent
The bold entries on the way down are the carries being deposited. Notice that $a_1$ shows $1$ after
its own block and then drops to $0$ after the \emph{next} one --- because each $\MAJ$ also XORs into
the wire below it. \textbf{The only wire holding a clean, unmodified carry at any moment is the
topmost one, $a_2$, immediately after the last $\MAJ$.} That is exactly where the single CNOT sits.
The placement is not arbitrary: it is the one instant at which the final carry is sitting in the
clear.

\begin{keybox}[title={Verified}]
The circuit was built as a literal $\CNOT/\TOF$ list and applied to all $2^{2n+2}$ basis states for
$n=2,3,4,5$: the map is a permutation (reversible), the gate counts are exactly $2n$ Toffoli and
$4n{+}1$ CNOT, and $\big(b,a,c_0,z\big)\mapsto\big((a{+}b)\bmod 2^n,\,a,\,0,\,z\oplus c_n\big)$
holds for every input pair. The trace above is machine-generated.
\end{keybox}

\begin{remark}[where the ancillas went]
CDKM was a genuine economy. The first reversible adders --- Vedral--Barenco--Ekert's \cite{vbe},
and the CARRY/SUM networks Draper \cite{draper} reviews before departing from them --- kept every
carry in its own ancilla: $3n$ qubits to add two $n$-bit numbers. The $\MAJ$ trick stores each
carry \emph{in place of} an operand bit it has already consumed, and $3n$ drops to $2n+2$.
Draper's own answer to the same problem is more radical still --- no carries at all --- and is the
Fourier route of \S\ref{sec:code-qft}.
\end{remark}

\subsection{Subtraction and comparison}
Running Theorem~\ref{thm:cdkm} \emph{backwards} gives $\SUB$: $\ket b\ket a\mapsto\ket{b-a\bmod2^n}\ket a$.

For comparison, the trick is to make the answer free rather than to build a comparator. Run the
$\MAJ$ chain on $a$ and $\bar b$ (bitwise complement) with carry-in $1$: the register then holds
$a+\bar b+1=a-b+2^n$, whose carry-out satisfies
\begin{equation}\label{eq:cmp}
c_n=[\,a\ge b\,].
\end{equation}
Copy it with one CNOT and reverse the chain: $n$ compute/uncompute Toffoli pairs, $4n$ $T$.
Equivalently --- and this is the form Level~2 uses --- work one bit wider than you need, in two's
complement, and the \emph{top bit} of $x-y$ \emph{is} the predicate $[\,x<y\,]$. \textbf{Comparison
therefore costs nothing beyond the subtraction you were doing anyway}; you simply look at a qubit
that is already there. That single sign bit is the workhorse of the whole level above.

\begin{remark}[compile-time constants]
Every addition inside Shor's algorithm has one operand \emph{classically known}: the constants
$2^ia\bmod N$ and the modulus $N$. Specialising the circuit to a fixed operand deletes or fixes
roughly half the gates before emission --- controls on known-$0$ bits vanish, controls on known-$1$
bits become unconditional. See the box in \S\ref{sec:l3} for what this does and does not buy.
\end{remark}

\section{Level 2: modular addition by a classical constant}\label{sec:l2}

\textbf{The problem this level solves, in one sentence:} you want $(x+a)\bmod N$, which classically
is ``add them, and \emph{if} the result is at least $N$, subtract $N$'' --- and a quantum computer
cannot do \emph{if}.

It cannot branch, because in general it is holding a superposition in which some branches
overshot and some did not; there is no single answer to ``did it overshoot?''. So the conditional
must be turned into arithmetic: compute the condition into a qubit, then use that qubit as a
control. Blocks (1)--(4) below do exactly that, and if you could stop there the arithmetic would be
finished. \textbf{Blocks (5)--(7) exist entirely to destroy the evidence}, for the reason in
\S\ref{sec:garbage}.

\begin{theorem}[Beauregard's modular adder \cite{beauregard}]\label{thm:modadd}
For $0\le a<N$ the seven blocks below implement $\ket x\mapsto\ket{(x+a)\bmod N}$ for all
$0\le x<N$, using five constant additions and one ancilla returned to $\ket 0$. The value register
is $n{+}1$ qubits wide so that intermediates up to $2N$ fit.
\end{theorem}

\begin{center}
\begin{quantikz}[column sep=10pt,row sep=15pt]
\lstick{$\ket{x}_{n+1}$} & \gate{+a} & \gate{-N} & \ctrl{1} & \gate{+N} & \gate{-a} & \qw & \ctrl{1} & \gate{+a} & \qw & \rstick{$\ket{(x{+}a)\bmod N}$}\\
\lstick{$\ket{t{=}0}$}   & \qw       & \qw       & \targ{}  & \ctrl{-1} & \qw       & \gate{X} & \targ{} & \qw & \qw & \rstick{$\ket{0}$}
\end{quantikz}
\end{center}
{\small\noindent Left to right, the columns are blocks (1)--(7) of the table below (block (6) is the
$X$ together with the CNOT that follows it). The two vertical lines into $t$ are CNOTs \emph{from
the sign qubit} --- the top bit of the $(n{+}1)$-qubit value register --- and block (4) is the
addition of $N$ controlled on $t$.}

\medskip
\begin{center}
\begin{tabular}{@{}rlll@{}}
\toprule
& block & value register & $t$\\
\midrule
(1) & $\ADD(a)$ & $x+a\ (<2N$, fits$)$ & $0$\\
(2) & $\SUB(N)$ & $x+a-N$, sign bit $=[\,x+a<N\,]$ & $0$\\
(3) & CNOT sign $\to t$ & $x+a-N$ & $[\,x+a<N\,]$\\
(4) & $\ADD(N)$ controlled on $t$ & $R:=(x+a)\bmod N$ & unchanged\\
(5) & $\SUB(a)$ & $R-a$, sign bit $=[\,R<a\,]$ & unchanged\\
(6) & $X$ on $t$; CNOT sign $\to t$ & $R-a$ & $\mathbf 0$\\
(7) & $\ADD(a)$ & $R$ & $0$\\
\bottomrule
\end{tabular}
\end{center}

\begin{proof}
Blocks (1)--(2) form $x+a-N$, which is negative exactly when $x+a<N$; in two's complement the sign
qubit \emph{is} that predicate, so (3) records it and (4) restores $R\in[0,N)$. The remaining
question is how to clear the ancilla, and the answer is that the flag is recomputable from the
output alone. Since $x<N$,
\[
t=1\iff x+a<N\iff R=x+a\ge a,
\qquad
t=0\iff R=x+a-N<a,
\]
so $[\,R<a\,]=\neg t$. Block (5) forms $R-a$, whose sign qubit is therefore $\neg t$; the $X$ in (6)
maps $t\mapsto\neg t$ and the following CNOT gives $\neg t\oplus\neg t=0$. Block (7) restores $R$.
\end{proof}

\begin{keybox}[title={The idea: deduce the flag, don't remember it}]
After block (4) you hold the right answer but you are stuck with a qubit remembering whether you
wrapped around. You may not erase it and you may not discard it.

The way out is that \textbf{you do not need to remember, because you can work it out again from the
answer}. Since $x$ started below $N$, wrapping happened exactly when $R<a$. So block (5) recomputes
the flag from the output and block (6) XORs it against itself.

It is like this: you cannot remember whether you borrowed during a subtraction, but you can look at
your answer and deduce that you must have --- and then honestly forget. \textbf{Every block in this
tower ends by recomputing its own flag from its own output.} That single discipline is what keeps
the whole construction garbage-free.
\end{keybox}

\begin{warnbox}[title={The precondition is load-bearing}]
The proof uses $x<N$ twice. Fed a value $\ge N$ the circuit silently leaves $t$ dirty and the value
register wrong by $N$ --- and by \S\ref{sec:garbage} a dirty $t$ does not produce a wrong answer,
it produces \emph{no} answer, since the interference downstream dies.

The invariant holds because \emph{every} adder in the tower reduces mod $N$ (\S\ref{sec:l3}). Note
that this is the same hypothesis that licenses the $3$-of-$7$ control saving in \S\ref{sec:l4}: the
uncontrolled blocks compose to the identity only when $b<N$. One precondition, two uses.
\end{warnbox}

\begin{remark}[traced, $N=15$]
$x=3,\ a=7$: $3{+}7{=}10\to10{-}15{=}-5$ (sign $1$) $\to t{=}1\to+15\to R{=}10\to10{-}7{=}3$
(sign $0$) $\to t=1\oplus1\oplus0=0$. \quad
$x=9,\ a=7$: $9{+}7{=}16\to16{-}15{=}1$ (sign $0$) $\to t{=}0\to R{=}1\to1{-}7{=}-6$ (sign $1$)
$\to t=0\oplus1\oplus1=0$. Both branches: correct output, ancilla clean.
\end{remark}

\section{Level 3: modular multiplication}\label{sec:l3}

\textbf{The problem this level solves.} You have an adder; you need a multiplier. The bridge is the
schoolbook one --- reversible shift-and-add multiplication goes back to Vedral, Barenco and Ekert
\cite{vbe} --- and each shifted multiple is a \emph{classically known constant}. For
$x=\sum_i x_i2^i$,
\begin{equation}\label{eq:shiftadd}
a\,x\bmod N=\Big(\sum_{i=0}^{n-1}x_i\,\big(2^ia\bmod N\big)\Big)\bmod N ,
\end{equation}
i.e.\ binary long multiplication written reversibly. The bit $x_i$ decides whether its multiple
joins the sum, so it becomes a \emph{control}: the data has turned into conditioning.

\begin{theorem}\label{thm:cmult}
\[
\CMULT(a):\ \ket c\ket x\ket b\ \mapsto\ \ket c\ket x\ket{\,b+c\cdot ax\bmod N\,}
\qquad\text{is}\qquad
\prod_{i=0}^{n-1}\ADD^{\bmod N}_{c,\,x_i}\!\big(2^ia\bmod N\big),
\]
$n$ doubly-controlled modular additions of compile-time constants.
\end{theorem}

\begin{center}
\begin{quantikz}[column sep=9pt,row sep=9pt]
\lstick{$\ket c$}    & \ctrl{5} & \ctrl{5} & \ctrl{5} & \ \ldots\  & \ctrl{5} & \qw & \rstick{$\ket c$}\\
\lstick{$\ket{x_0}$} & \ctrl{4} & \qw      & \qw      &            & \qw      & \qw & \rstick{$\ket{x_0}$}\\
\lstick{$\ket{x_1}$} & \qw      & \ctrl{3} & \qw      &            & \qw      & \qw & \rstick{$\ket{x_1}$}\\
\lstick{$\ket{x_2}$} & \qw      & \qw      & \ctrl{2} &            & \qw      & \qw & \rstick{$\ket{x_2}$}\\
\lstick{$\vdots$}    &          &          &          & \raisebox{1pt}{$\ddots$} & \ctrl{1} & & \rstick{$\vdots$}\\
\lstick{$\ket{b}$}   & \gate{{+}a} & \gate{{+}2a} & \gate{{+}4a} & \ \ldots\ & \gate{{+}2^{n-1}a} & \qw & \rstick{$\ket{b{+}ax}$}
\end{quantikz}
\end{center}
{\small\noindent Each box is one Level-2 modular adder for the constant $2^ia\bmod N$, doubly
controlled by the outer control $c$ and the bit $x_i$. Every partial sum is reduced mod $N$, which
is what maintains the invariant $b<N$ that Theorem~\ref{thm:modadd} requires.}

\begin{proof}
Immediate from \eqref{eq:shiftadd}: the $i$-th adder fires exactly when $x_i=1$, and each adder
reduces, so all partial sums stay below $N$.
\end{proof}

\begin{keybox}[title={The constant is known; the operand is not}]
It is tempting to think that because $a$ is chosen in advance, the multiplication is ``already
done''. But a multiplication has \emph{two} operands, and precomputation gives you one:
\[
\underbrace{a}_{\text{known at compile time}}\;\times\;\underbrace{x}_{\text{unknown --- it is the register}}\bmod N .
\]
The register holds $a^x\bmod N$ for a counting register in superposition over all $2^t$ values of
$x$. If you knew which $x$ mattered you would compute $a^x\bmod N$ on a laptop and never build the
machine. The machine exists precisely because it evaluates at every input at once --- so every
multiplication in the ladder is fed a \emph{superposition}, and must be a circuit that maps all
inputs correctly and simultaneously. For $N=15$, $a=7$ the first rung must realise
\[
1\!\to\!7,\quad 2\!\to\!14,\quad 4\!\to\!13,\quad 7\!\to\!4,\quad
8\!\to\!11,\quad 11\!\to\!2,\quad 13\!\to\!1,\quad 14\!\to\!8
\]
in one shot. Precomputing the whole input/output table instead would need $2^n$ rows: about
$10^{616}$ for RSA-2048, against $10^{80}$ atoms in the observable universe. The constant, by
contrast, is $2048$ bits --- $256$ bytes. That gap is exactly the gap between what precomputation
gives you and what still has to be built.

What precomputation \emph{does} buy is real and is used everywhere: the rung constants
(\S\ref{sec:l5}), the shifted multiples $2^ia\bmod N$ above, the inverse $\bar a$
(\S\ref{sec:l4}), and the roughly half of all gates that can be deleted because one adder operand
is classical.
\end{keybox}

\begin{remark}[$N=15$, $a=7$]
The four constants are $2^i\cdot7\bmod15=7,14,13,11$. For $x=3=0011_2$ the first two adders fire:
$7+14=21\equiv6\pmod{15}$, and indeed $7\cdot3=21\equiv6$.
\end{remark}

\section{Level 4: the controlled multiplier, in place}\label{sec:l4}

\textbf{The problem this level solves.} $\CMULT$ writes its answer into a \emph{fresh} register,
but the ladder must apply many multiplications one after another to the \emph{same} register. A
circuit that consumed a clean register per rung would need thousands of registers. So the
out-of-place multiplier has to be converted into an in-place map $\ket x\mapsto\ket{ax\bmod N}$,
and the leftover input $x$ has to disappear --- which, by \S\ref{sec:garbage}, means it must be
\emph{computed away}, not discarded.

The escape is that $a$ is invertible. You can uncompute $x$ from the product. The construction is
Beauregard's \cite{beauregard}.

\begin{theorem}\label{thm:cu}
Let $\gcd(a,N)=1$ and $\bar a:=a^{-1}\bmod N$ (classical, extended Euclid). Then
\[
c\text{-}U_a\;=\;\CMULT(\bar a)^\dagger\cdot c\text{-}\mathrm{SWAP}\cdot\CMULT(a)
\]
implements $\ket c\ket x\ket 0\mapsto\ket c\ket{a^cx\bmod N}\ket0$ for all $0\le x<N$.
\end{theorem}

\begin{center}
\begin{quantikz}[column sep=9pt,row sep=11pt]
\lstick{$\ket c$}    & \ctrl{1} & \ctrl{1} & \ctrl{1} & \qw & \rstick{$\ket c$}\\
\lstick{$\ket x$}    & \gate[2]{\CMULT(a)} & \swap{1} & \gate[2]{\CMULT(\bar a)^{\dagger}} & \qw & \rstick{$\ket{a^cx \bmod N}$}\\
\lstick{$\ket{0}_b$} & & \targX{} & & \qw & \rstick{$\ket 0$}
\end{quantikz}
\end{center}

\begin{proof}
Take $c=1$ and follow the register pair:
\[
\ket x\ket0
\ \xrightarrow{\ \CMULT(a)\ }\ \ket x\ket{ax}
\ \xrightarrow{\ \mathrm{SWAP}\ }\ \ket{ax}\ket x
\ \xrightarrow{\ \CMULT(\bar a)^\dagger\ }\ \ket{ax}\ket{x-\bar a(ax)}=\ket{ax}\ket0,
\]
since $\bar aa\equiv1\pmod N$. For $c=0$ every factor is the identity. Both branches leave the
scratch in $\ket0$, unentangled, so $c\text{-}U_a$ is realised exactly.
\end{proof}

\noindent
The controlled SWAP is $n$ Fredkin gates sharing one control. Note that no controlled QFT appears
and the constant $\bar a$ carries no control --- its inverse is classical. For $N=15,a=7$:
$\bar a=13$, since $7\cdot13=91=6\cdot15+1$.

\begin{keybox}[title={Where the outer control goes, and why not everywhere}]
Let $V=G_L\cdots G_1$ and let $V_S$ replace $G_j$ by $c\text{-}G_j$ for $j\in S$. Since
$V_S=\ket0\!\bra0\otimes R_S+\ket1\!\bra1\otimes V$ with $R_S=\prod_{j\notin S}G_j$, we have
$V_S=c\text{-}V$ on a subspace $W$ \emph{iff} $R_S=\id$ on $W$. In words: \textbf{a gate needs the
control only if leaving it uncontrolled would not cancel out}. Two patterns exhaust the uses:
\begin{itemize}\setlength\itemsep{1pt}
\item \emph{Conjugation.} If $V=ABA^\dagger$ then controlling $B$ alone suffices, since
$AA^\dagger=\id$. A conjugating layer never needs the control.
\item \emph{Compensating residue.} In Theorem~\ref{thm:modadd} only blocks (1), (5), (7) take the
outer control: with the control off, blocks (2), (3), (4), (6) conspire to do nothing --- $b-N$
goes negative, so $t$ is set, block (4) adds $N$ back, and block (6) clears $t$. This works only
because $b<N$, the same precondition as before.
\end{itemize}
Measured on $N=15$, $a=7$, transpiled to $\{\CNOT,U\}$: the modular adder costs $2\,634$ CNOTs at
depth $4\,836$ if every gate is controlled, against $206$ / $354$ for the $3$-of-$7$ construction
($12.8\times$); the full $c\text{-}U_a$, $20\,892$ / $37\,727$ against $1\,748$ / $2\,957$
($12.0\times$). An order of magnitude, from one lemma. Both patterns reappear, letter for letter,
in the code of Part~IV (\S\ref{sec:code-l1}, \S\ref{sec:code-l2}).
\end{keybox}

\section{Level 5: the ladder}\label{sec:l5}

\textbf{What this level supplies.} Gift~1 of \S1.1. The rungs use classically precomputed constants
$c_j=a^{2^j}\bmod N$, obtained by repeated squaring $c_0=a$, $c_{j+1}=c_j^2\bmod N$ --- this is the
identity $U_a^{2^j}=U_{a^{2^j}\bmod N}$, which is why the ladder has $t$ rungs and not $2^t$.

\begin{center}
\begin{quantikz}[column sep=10pt,row sep=10pt]
\lstick{$\ket0$} & \gate{H} & \ctrl{4} & \qw & \qw & \ \ldots\ & \qw & \gate[4]{F^\dagger_{2^t}} & \meter{}\\
\lstick{$\ket0$} & \gate{H} & \qw & \ctrl{3} & \qw & & \qw & & \meter{}\\
\lstick{$\vdots$} & & & & & & & & \\
\lstick{$\ket0$} & \gate{H} & \qw & \qw & \qw & \ \ldots\ & \ctrl{1} & & \meter{}\\
\lstick{$\ket1_n$} & \qw & \gate{\times a} & \gate{\times a^{2}} & \qw & \ \ldots\ & \gate{\times a^{2^{t-1}}} & \qw & \qw
\end{quantikz}
\end{center}
{\small\noindent Each $\times a^{2^j}$ box is one $c\text{-}U_a$ of Theorem~\ref{thm:cu}, hence
$2n$ modular adders (Theorem~\ref{thm:cmult}) and $n$ Fredkins; each modular adder is five constant
additions (Theorem~\ref{thm:modadd}); each addition is the CDKM circuit (Theorem~\ref{thm:cdkm}) at
$2n$ Toffoli.}

\begin{keybox}[title={The ladder is not one circuit}]
Every rung uses a \emph{different} constant, so the rungs are $t$ structurally different circuits
that happen to sit side by side. For $N=15,a=7$ the constants are $7,4,1,1,1,\dots$ --- so most of
that ladder is literally the identity, which is why $15$ is everyone's demo integer.

A consequence worth stating, because it surprises people: \textbf{the compiled circuit is specific
to $a$}, and Shor's algorithm may need to retry with a fresh base. That is genuine, and it is
cheap. Recompiling means repeated squaring plus one extended Euclid --- $O(n^3)$ classical bit
operations, measured at $50$\,ms for a $2048$-bit modulus, against days of quantum runtime. And it
is rarely needed: a bad measurement or a $\gcd(s,r)>1$ failure is re-run on the \emph{same} circuit,
and only a bad base forces recompilation, which happens with probability at most $1/2$ for $N=pq$.
Expected number of distinct circuits: about two.
\end{keybox}

\part*{Part III --- Closing the loop}
\addcontentsline{toc}{part}{Part III --- Closing the loop}

\section{From permutations to phases}

Now follow the state through, and watch classical arithmetic turn into an eigenphase. After the
Hadamard layer,
\[
\frac1{\sqrt{2^t}}\sum_{x=0}^{2^t-1}\ket x\ket1 .
\]
The ladder is the only place any arithmetic happens, and by \S\ref{sec:l5} it computes
\[
\ket x\ket y\longmapsto\ket x\,\big|\,a^{\sum_jx_j2^j}y\bmod N\,\big\rangle=\ket x\ket{a^xy\bmod N},
\qquad\text{so}\qquad
\frac1{\sqrt{2^t}}\sum_x\ket x\ket{a^x\bmod N}.
\]
Every gate that produced this was a permutation matrix with entries $0$ and $1$. Now rewrite the
\emph{same state} in the eigenbasis of $U_a$, using $\ket1=r^{-1/2}\sum_s\ket{u_s}$:
\[
\frac1{\sqrt{2^t}}\sum_x\ket x\ket{a^x\bmod N}
\;=\;
\frac1{\sqrt r}\sum_{s=0}^{r-1}
\underbrace{\Big[\frac1{\sqrt{2^t}}\sum_x e^{2\pi ixs/r}\ket x\Big]}_{\text{a phase ramp of frequency }s/r}
\otimes\ket{u_s}.
\]
The phases were not written by any gate. They are what the permutation \emph{looks like} in the
other basis --- and they are visible only because the register was in a superposition, so that a
permutation of basis labels became a relative phase between branches. The inverse QFT then reads the
frequency of the ramp, and the measurement returns an approximation to $s/r$ for a uniformly random
$s$; continued fractions recover $r$ \cite{shor}.

\begin{keybox}[title={The whole story in four lines}]
\begin{enumerate}\setlength\itemsep{1pt}
\item Phase estimation is cheap only when $U$ has cheap powers; arithmetic makes $U_a$'s powers
cheap (Gift 1) \emph{and} makes one power cheap to build (Gift 2).
\item Building it means reversible classical arithmetic --- adders, a modular reduction, a
shift-and-add multiplier, an in-place trick that needs $a^{-1}$.
\item Every scratch bit must be uncomputed, or the interference that reads the phase does not
happen. Half of every circuit in Part~II is that uncomputation.
\item Nothing in the arithmetic is complex-valued. The phases appear on the change of basis, and
that change of basis is the algorithm.
\end{enumerate}
\end{keybox}

\section{The bill}

\[
\underbrace{2n}_{\text{rungs}}\times\underbrace{2n}_{\substack{\text{modular adds}\\ \text{per }c\text{-}U_a}}
\times\underbrace{5}_{\substack{\text{additions per}\\ \text{modular add}}}\times\underbrace{2n}_{\substack{\text{Toffoli per}\\ \text{addition}}}
\;\approx\;40\,n^3\ \text{Toffoli},
\]
which windowed arithmetic \cite{gidney-window} plus the $3$-of-$7$ control trick reduces to
$\approx0.3\,n^3$. Windowing is a \emph{different} construction from the one built here --- it
replaces the controlled blocks by table lookups addressed by the exponent bits --- and it is built
and measured in Part~V (\S\ref{im:sec:window}).

At $n=2048$ that is $\approx2.6\times10^9$ Toffoli for the modular exponentiation against
$\approx10^5$ gates for the approximate QFT: \textbf{the Fourier layer, which carries all the
conceptual content, is $0.003\%$ of the circuit.} Everything that determines whether RSA falls is
the cost of multiplying two $2048$-bit numbers, reversibly, in superposition.

\begin{keybox}[title={What was machine-checked}]
$\MAJ$/$\UMA$ (Lemma~\ref{lem:majuma}) on all $8$ inputs and as $8\times8$ orthogonal matrices;
the Clifford$+T$ Toffoli against the $8\times8$ permutation ($7$ $T$ gates);
the CDKM circuit (Theorem~\ref{thm:cdkm}) as a literal gate list on all $2^{2n+2}$ basis states for
$n=2,3,4,5$ --- reversibility, gate counts, and arithmetic --- and the $a=5,b=3$ trace;
the comparator identity \eqref{eq:cmp} for all pairs at $n=4,6$;
the seven-block modular adder (Theorem~\ref{thm:modadd}) for all $x,a<N$ at $N=15,21,33,55$,
checking both the output and that $t$ returns to $0$;
the shift-and-add identity \eqref{eq:shiftadd} and the conjugation of Theorem~\ref{thm:cu} for all
$x<N$ and all $a$ coprime to $N$ at $N=15,21,55$;
and the ladder identity $\ket x\ket1\mapsto\ket x\ket{a^x\bmod N}$.
\end{keybox}


\part*{Part IV --- Building it: a verified implementation}
\addcontentsline{toc}{part}{Part IV --- Building it: a verified implementation}

\noindent
Parts~I--III said what the arithmetic must do and why. This part does it: a complete order-finding
implementation in about $150$ lines of Qiskit, bottom up, with the test at each level that makes it
believable. It is self-contained --- no circuit-library arithmetic, and the QFT built from scratch
--- and it takes the Fourier route, because that is the shortest correct path from nothing to a
factored integer. Part~V rebuilds the slow parts three different ways.

\section{Level 0: the QFT}\label{im:sec:l0}

Everything in Parts~IV and~V needs only these imports --- no Qiskit circuit-library
primitives, no prebuilt arithmetic:

\begin{lstlisting}
import math
from fractions import Fraction
import numpy as np
from qiskit.circuit import ClassicalRegister, QuantumCircuit, QuantumRegister
\end{lstlisting}

The one primitive everything rests on is the QFT, so build it rather than import it. It is
\[
\ket j\ \longmapsto\ \frac1{\sqrt{2^n}}\sum_{k=0}^{2^n-1}e^{2\pi i\,jk/2^n}\ket k,
\]
and the standard circuit follows from the product form of that sum: for each qubit, one Hadamard,
then a controlled phase from every less significant qubit, at halving angles. A final layer of
swaps reverses the qubit order, which the product form leaves reversed.

\begin{lstlisting}
def qft(n, approx=0):
    """QFT on n qubits, Qiskit convention (qubit 0 = LSB):
        |j> -> 2^(-n/2) sum_k exp(2*pi*i*j*k/2^n) |k>.
    approx > 0 drops controlled rotations by less than pi/2^approx."""
    qc = QuantumCircuit(n, name="qft")
    for j in reversed(range(n)):
        qc.h(j)
        for k in range(j):
            if approx and (j - k) > approx:
                continue                     # angle too small to matter
            qc.cp(np.pi / 2 ** (j - k), k, j)
    for i in range(n // 2):                  # reverse the qubit order
        qc.swap(i, n - 1 - i)
    return qc.to_gate()
\end{lstlisting}

That is $n(n{+}1)/2$ gates plus $\lfloor n/2\rfloor$ swaps: $O(n^2)$. The \code{approx} parameter
is the approximate QFT of \S\ref{im:sec:beyond} --- rotations by less than $\pi/2^{d}$ are below any
hardware's resolution, and dropping them costs negligible error while cutting the count to
$O(n\log n)$. It is one \code{continue}, so there is no reason not to have it.

\begin{keybox}[title={Verify this before anything else}]
Do not test the QFT by eye. Build its unitary and compare it with the definition:
\begin{lstlisting}
ref = np.array([[np.exp(2j*np.pi*j*k/2**n) for j in range(2**n)]
                for k in range(2**n)]) / np.sqrt(2**n)
assert np.allclose(Operator(qft(n)).data, ref)        # qiskit.quantum_info.Operator
\end{lstlisting}
Checked for $n=1..6$ against both the definition and Qiskit's own \code{QFTGate}: they agree
exactly, and \code{qft(n).inverse()} undoes it. If you drop the swap layer, or reverse the loop
order, you get a QFT of the bit-reversed integer --- which is still unitary, still looks fine, and
will make every addition below wrong. This is the single most common place to lose an afternoon.
\end{keybox}

\section{Level 1: addition by a classical constant}\label{im:sec:l1}

Classically, addition means carries, and carries mean a chain of Toffolis and ancilla qubits.
Draper's insight \cite{draper} is to avoid them entirely by working in the Fourier basis. Write the
QFT of an $n$-qubit register as a product state:
\[
F\ket y=\bigotimes_{i=0}^{n-1}\frac{\ket0+e^{2\pi i\,y\,2^{\,i-n}}\ket1}{\sqrt2}.
\]
Every qubit carries the whole value $y$ --- in the \emph{phase}, at a different scale. Adding $X$
means advancing each phase by the corresponding amount, so
\[
\text{qubit } i \ \text{gets}\ P\big(2\pi X\,2^{\,i-n}\big),
\qquad
e^{2\pi iy2^{i-n}}\longmapsto e^{2\pi i(y+X)2^{i-n}},
\]
and an inverse QFT lands back in the computational basis holding $\ket{(y+X)\bmod 2^n}$. No
carries, no ancillas, and since $X$ is classical every qubit needs exactly \emph{one} phase gate,
so the addition layer has depth~$1$.

\begin{lstlisting}
def add_const(qc, X, y):
    """|y> -> |y + X mod 2^n> on the qubit list y (LSB first)."""
    n = len(y)
    qc.append(qft(n), y)
    for i in range(n):
        qc.p(2 * np.pi * X * 2.0 ** (i - n), y[i])
    qc.append(qft(n).inverse(), y)

def c_add_const(qc, ctrls, X, y):
    """add_const controlled on every qubit in ctrls."""
    n = len(y)
    qc.append(qft(n), y)
    for i in range(n):
        qc.mcp(2 * np.pi * X * 2.0 ** (i - n), ctrls, y[i])
    qc.append(qft(n).inverse(), y)
\end{lstlisting}

\begin{keybox}[title={Three things to notice}]
\textbf{Negative constants just work.} $X=-N$ needs no separate subtractor: the phases are
periodic, so $-N$ and $2^{n}-N$ give the same gates. Subtraction is addition of a negative Python
integer, and \S\ref{im:sec:l2} relies on this.

\textbf{Only the phases take the control.} In \code{c\_add\_const} the two QFT layers are
\emph{uncontrolled}. They form a conjugation $F^\dagger\cdots F$, so with the control off they
cancel: controlling them would be pure waste. Keep this in mind --- it is the same argument that
saves an order of magnitude in \S\ref{im:sec:l2}.

\textbf{The composite is a permutation.} A single $P(\theta)$ is not, and intermediate states carry
irrational phases, but $F^\dagger\cdot(\text{phases})\cdot F$ is exactly the $0/1$ permutation
matrix $\ket y\mapsto\ket{y+X}$. Nothing above Level~1 can tell this apart from a ripple-carry
adder.
\end{keybox}

\section{Level 2: modular addition --- the hard one}\label{im:sec:l2}

You want $(y+X)\bmod N$, which classically is ``add, and \emph{if} the result is $\ge N$, subtract
$N$''. A quantum computer cannot branch: it holds a superposition in which some branches overshot
and some did not, so there is no single answer to ``did it overshoot?''.

So compute the condition into a qubit and use that qubit as a control. Give the value register one
extra qubit (the \emph{sign} qubit) so intermediates up to $2N$ fit and two's-complement negatives
are readable: after subtracting, the top bit \emph{is} the predicate. Then:

\begin{center}\small
\begin{tabular}{@{}rlll@{}}
\toprule
& block & value register & flag $t$\\
\midrule
(1) & $+X$ (controlled) & $y+X$ \ ($<2N$, fits) & $0$\\
(2) & $-N$ & $y+X-N$, sign bit $=[\,y+X<N\,]$ & $0$\\
(3) & CNOT sign $\to t$ & unchanged & $[\,y+X<N\,]$\\
(4) & $+N$ if $t$ & $R:=(y+X)\bmod N$ & unchanged\\
\midrule
(5) & $-X$ & $R-X$, sign bit $=\neg t$ & unchanged\\
(6) & clear $t$ from that sign bit & $R-X$ & $\mathbf 0$\\
(7) & $+X$ & $R$ & $0$\\
\bottomrule
\end{tabular}
\end{center}

Blocks (1)--(4) do the arithmetic. \textbf{Blocks (5)--(7) exist entirely to destroy the evidence},
and they are the interesting half. After (4) you hold the right answer but one qubit remembers
whether you wrapped --- which you may neither erase nor discard.

The escape: \emph{you do not need to remember, because you can deduce it from the answer.} Since
$y<N$, wrapping happened exactly when $R\ge X$. So block (5) recomputes the flag from the output
and block (6) XORs it against itself. Concretely, $t=1\iff y+X<N\iff R=y+X\ge X$, so
$[\,R<X\,]=\neg t$, which is the sign bit after (5); XOR that into $t$ and apply $X$ to finish.

\begin{lstlisting}
def c_add_mod(qc, ctrls, X, y, sign, flag, N):
    """|y> -> |(y + X) mod N> when all ctrls are 1; identity otherwise.
    Needs 0 <= X < N, 0 <= y < N, sign and flag in |0>."""
    yw = list(y) + [sign]                     # (n+1)-wide value register
    c_add_const(qc, ctrls, X, yw)             # (1) + X          (controlled)
    add_const(qc, -N, yw)                     # (2) - N
    qc.cx(sign, flag)                         # (3) sign -> flag
    c_add_const(qc, [flag], N, yw)            # (4) + N  if flag
    add_const(qc, -X, yw)                     # (5) - X
    qc.mcx(list(ctrls) + [sign], flag)        # (6) clear flag from the answer
    qc.x(flag)
    add_const(qc, X, yw)                      # (7) + X
\end{lstlisting}

\begin{keybox}[title={Where the control goes --- the biggest constant factor in the circuit}]
Notice what is \emph{not} controlled. Blocks (2)--(5) and (7) are plain, and inside every
\code{add\_const} the QFT layers are plain too. Only block (1), block (4)'s $+N$, and block (6)'s
X carry a control. The governing lemma: replacing $G_j$ by $c\text{-}G_j$ only for $j\in S$ gives
a correct $c\text{-}V$ \emph{iff} the leftover product $\prod_{j\notin S}G_j$ is the identity on
the states you care about. Two patterns satisfy it, and both appear here.

\textbf{Conjugation.} If $V=ABA^\dagger$, controlling $B$ alone suffices, since $AA^\dagger$ is the identity.
This is the QFT sandwich in \S\ref{im:sec:l1}: never control a conjugating layer.

\textbf{Compensating residue.} Ask what the uncontrolled blocks do when the control is
\textbf{off}: $y-N$ goes negative $\Rightarrow$ sign set $\Rightarrow$ (3) sets $t$ $\Rightarrow$
(4) adds $N$ back. Then (5) subtracts $X$, the Toffoli in (6) does nothing (its outer control is
off), the \code{x} flips $t$ back to $0$, and (7) restores $X$. Everything cancels.
\textbf{This works only because $y<N$} --- the same precondition the correctness proof needs.
With the control \textbf{on}: after (5) the sign qubit holds $\neg t$, so the Toffoli sets
$t\mapsto t\oplus\neg t=1$ and the \code{x} clears it.

\textbf{Measured} ($N=15$, transpiled to $\{\CNOT,U\}$), against the naive circuit built by
wrapping an uncontrolled modular adder in \code{.control(k)} --- i.e.\ controlling every gate:

\begin{center}\small
\begin{tabular}{@{}lrrrr@{}}
\toprule
& \multicolumn{2}{c}{every gate controlled} & \multicolumn{2}{c}{as written here}\\
\cmidrule(lr){2-3}\cmidrule(lr){4-5}
& CNOT & depth & CNOT & depth\\
\midrule
1 control & $2\,353$ & $4\,351$ & $287$ & $372$\\
2 controls (the real usage, \S\ref{im:sec:l34}) & $5\,916$ & $10\,689$ & $315$ & $418$\\
\bottomrule
\end{tabular}
\end{center}

\noindent
$8\times$ and $19\times$. Almost all of it is the conjugation pattern --- leaving the QFT layers
alone. Isolating the residue pattern (blocks controlled or not, QFTs plain either way) is worth a
further $1.6\times$ at two controls. Do both; the first is nearly free to get right.

If you prefer, Beauregard's original choice is to control blocks (1), (5), (7) and leave (6) a
plain CNOT. Both are correct, and only the residue bookkeeping changes.
\end{keybox}

\begin{warnbox}[title={The precondition is load-bearing}]
The proof uses $y<N$ twice. Fed $y\ge N$, the circuit leaves $t$ dirty and the value wrong --- and
a dirty $t$ does not give a wrong answer, it gives \emph{no} answer. The invariant holds because
every adder above reduces mod $N$; if you ever add without reducing, you have broken it.
\end{warnbox}

\section{Levels 3--4: multiplication, in place}\label{im:sec:l34}

\textbf{Level 3, accumulate.} Multiplication is shifted addition, and each shifted multiple
$2^iA\bmod N$ is classically known. Bit $x_i$ of the register decides whether its multiple joins
the sum --- so it becomes a \emph{control}, and the data turns into conditioning:
\[
A\,x \bmod N=\Big(\sum_{i=0}^{n-1}x_i\,\big(2^iA\bmod N\big)\Big)\bmod N .
\]
Every partial sum is reduced, which is what maintains the $y<N$ invariant Level~2 demands.

\begin{lstlisting}
def c_mult_acc(qc, c, A, x, y, sign, flag, N):
    """|c>|x>|y> -> |c>|x>|(y + A*x) mod N> : shift-and-add accumulate."""
    for i in range(len(x)):
        c_add_mod(qc, [c, x[i]], (A << i) % N, y, sign, flag, N)
\end{lstlisting}

\textbf{Level 4, in place.} \code{c\_mult\_acc} writes into a \emph{fresh} register, but the ladder
must apply rung after rung to the \emph{same} register. Swapping is easy; the problem is that the
old input $x$ is still sitting in the scratch, and by the no-garbage rule it must be computed away,
not discarded.

The escape is that $A$ is invertible. Multiply, swap, then subtract $A^{-1}$ times the product:
\[
\ket x\ket0
\xrightarrow{\ \times A\ }\ket x\ket{Ax}
\xrightarrow{\ \mathrm{SWAP}\ }\ket{Ax}\ket x
\xrightarrow{\ -A^{-1}\cdot\ }\ket{Ax}\ket{x-A^{-1}Ax}=\ket{Ax}\ket0 .
\]

\begin{lstlisting}
def c_ua(qc, c, A, x, y, sign, flag, N):
    """|c>|x>|0> -> |c>|A^c x mod N>|0> : multiply, swap, uncompute."""
    c_mult_acc(qc, c, A % N, x, y, sign, flag, N)
    for i in range(len(x)):
        qc.cswap(c, x[i], y[i])
    B = pow(A, -1, N)                        # classical inverse, gcd(A,N)=1
    c_mult_acc(qc, c, (-B) % N, x, y, sign, flag, N)
\end{lstlisting}

Note what is \emph{not} here: no circuit is reversed. The adjoint of ``add $A^{-1}x$ mod $N$'' is
``add $-A^{-1}x$ mod $N$'', so the dagger becomes a minus sign on a classical constant. Note also
that $\gcd(a,N)=1$ is doing two jobs at once: it makes $y\mapsto ay$ a permutation (hence unitary,
hence something phase estimation can act on \emph{at all}), and it makes $A^{-1}$ exist so the
in-place trick works. One hypothesis, two roles.

\section{Level 5: the ladder and the full circuit}\label{im:sec:l5}

\begin{lstlisting}
def order_circuit(A, N, t=None):
    """Counting register of t qubits (default 2n), target, n+2 scratch."""
    n = math.ceil(math.log2(N))
    t = t if t is not None else 2 * n
    ctr, tgt = QuantumRegister(t, "ctr"), QuantumRegister(n, "tgt")
    anc = QuantumRegister(n, "anc")
    sf = QuantumRegister(2, "sf")            # sign qubit, flag qubit
    out = ClassicalRegister(t, "out")
    qc = QuantumCircuit(ctr, tgt, anc, sf, out)
    qc.h(ctr)                                # Hadamard layer
    qc.x(tgt[0])                             # target = |1>
    for i in range(t):                       # the ladder: t rungs
        c_ua(qc, ctr[i], pow(A, 2**i, N), tgt, anc, sf[0], sf[1], N)
    qc.append(qft(t).inverse(), ctr)     # inverse QFT
    qc.measure(ctr, out)
    return qc
\end{lstlisting}

\textbf{Qubit budget:} $t$ counting $+\ n$ target $+\ n$ scratch product $+\ 1$ sign $+\ 1$ flag
$=4n+2$ with $t=2n$. For $N=15$: $8+4+4+2=18$.

\textbf{Why it works.} After the Hadamards and the ladder the state is
$\frac1{\sqrt{2^t}}\sum_x\ket x\ket{a^x\bmod N}$ --- and every gate that produced it was a
permutation matrix of $0$s and $1$s. Rewriting the \emph{same state} in the eigenbasis of $U_a$,
using $\ket1=r^{-1/2}\sum_s\ket{u_s}$:
\[
\frac1{\sqrt{2^t}}\sum_x\ket x\ket{a^x\bmod N}
=\frac1{\sqrt r}\sum_{s=0}^{r-1}
\underbrace{\Big[\tfrac1{\sqrt{2^t}}\sum_x e^{2\pi ixs/r}\ket x\Big]}_{\text{phase ramp of frequency } s/r}
\otimes\ket{u_s}.
\]
No gate wrote those phases. They are what a permutation of basis labels \emph{looks like} in
another basis, visible only because the register was in superposition. The inverse QFT reads the
frequency, and the measurement returns $\approx 2^t s/r$ for a uniformly random $s$.

\begin{keybox}[title={Why $t=2n$ --- and what happens if you skimp}]
The measured integer $y$ satisfies $y/2^t\approx s/r$. To pin $r<N$ down uniquely from a rational
approximation, continued fractions need the error below $1/2r^2$, and since $r$ can be nearly $N$
that means $2^t>N^2$, i.e.\ $t\ge2n$.

Beware the $N=15$ trap: there $r=4$ divides $2^t$ exactly, so $y/2^t$ is exact for \emph{any}
$t\ge2$ and a too-short register still works. $N=21$ with $a=2$ ($r=6$, $n=5$) is the honest test,
and running it at three precisions shows the whole story:
\begin{center}\small
\begin{tabular}{@{}lllp{0.42\textwidth}@{}}
\toprule
$t$ & qubits & recovered $r$ & what the peaks look like\\
\midrule
$5=n$   & $17$ & $\mathbf{18}$ --- wrong & $y/2^t$ lands on $7/20$, $3/19$: junk fractions\\
$7$     & $19$ & $6$ --- correct & $21/128\to1/6$, $107/128\to5/6$\\
$10=2n$ & $22$ & $6$ --- correct & clean peaks at $0,171,341,512,683,853$\\
\bottomrule
\end{tabular}
\end{center}
The $t=n$ failure is the instructive one. It did not return nonsense --- it returned $18$, which
\emph{passes} \code{pow(A, r, N) == 1} because $18$ is a multiple of $6$. A too-short register
gives you a plausible, self-consistent, wrong answer. ($t=7$ sufficing here is luck: the bound
$t\ge2n$ covers the worst case $r\approx N$, and you do not know $r$ in advance.)
\end{keybox}

\section{The classical half}\label{im:sec:classical}

\begin{lstlisting}
def order_from_counts(counts, A, N, t):
    for bits in sorted(counts, key=counts.get, reverse=True):
        y = int(bits, 2)
        if y == 0:
            continue
        r = Fraction(y, 2**t).limit_denominator(N - 1).denominator
        if pow(A, r, N) == 1:
            return r
    return 0

def find_factor(N, run):                     # run: circuit -> counts dict
    if N % 2 == 0:
        return 2
    for k in range(2, N.bit_length() + 1):   # prime-power check
        d = round(N ** (1 / k))
        if d > 1 and d**k == N:
            return d
    import random
    while True:
        a = random.randrange(2, N)
        if (d := math.gcd(a, N)) > 1:
            return d                         # lucky gcd
        t = 2 * math.ceil(math.log2(N))
        r = order_from_counts(run(order_circuit(a, N)), a, N, t)
        if r and r % 2 == 0:
            d = math.gcd(pow(a, r // 2, N) - 1, N)
            if 1 < d < N:
                return d
\end{lstlisting}

\code{limit\_denominator} \emph{is} the continued-fraction expansion --- Python ships the
number theory. Candidates are checked with \code{pow(A, r, N) == 1} before being believed;
$y=0$ carries no information and is skipped.

\begin{keybox}[title={Most of the peaks are useless --- by design}]
The loop looks like defensive coding. It is not: \textbf{a majority of the peaks cannot give you
$r$}, and you must expect to discard them. The measurement returns $y\approx2^ts/r$ for uniformly
random $s$, and continued fractions recover $s/r$ \emph{in lowest terms} --- so what you get back
is $r/\gcd(s,r)$. Only $s$ coprime to $r$ yields $r$ itself.

At $N=21$, $a=2$ ($r=6$, $t=10$), the six peaks land as:
\begin{center}\small
\begin{tabular}{@{}lcccccc@{}}
\toprule
$s$ & $0$ & $1$ & $2$ & $3$ & $4$ & $5$\\
\midrule
peak $y\approx$ & $0$ & $171$ & $341$ & $512$ & $683$ & $853$\\
$y/2^t\to$ & $0$ & $1/6$ & $1/3$ & $1/2$ & $2/3$ & $5/6$\\
recovered $r$ & $1$ & $\mathbf 6$ & $3$ & $2$ & $3$ & $\mathbf 6$\\
\code{pow} check & fails & \textbf{passes} & fails & fails & fails & \textbf{passes}\\
\bottomrule
\end{tabular}
\end{center}
Two of six --- that is $\varphi(6)/6$ --- and the other four return proper divisors of $r$ that the
\code{pow} check rejects. This is not hypothetical: in the measured run the most frequent
\emph{nonzero} outcome was $y=512$, which gives $r=2$, and $\gcd(2^1{-}1,21)=1$ factors nothing.
Trusting the top outcome would have failed; scanning all of them by descending frequency, as the
loop does, finds $y=171$ and $r=6$. The $N=33$ run repeats the lesson at larger scale: with
$r=10$ only $\varphi(10)/10=40\%$ of peaks work, the recovered fractions splitting cleanly into
$\{\tfrac1{10},\tfrac3{10},\tfrac7{10},\tfrac9{10}\}$ (accepted) and
$\{\tfrac15,\tfrac25,\tfrac35,\tfrac45\}$ (rejected), and again the top nonzero outcome was one of
the rejects.

One caveat the check cannot catch: \code{pow(A, r, N) == 1} accepts any \emph{multiple} of the true
order, so a stray $y$ can return $12$ instead of $6$. Harmless for factoring when the result is
even, but do not treat the returned value as certainly minimal.
\end{keybox}

Retries are cheap and expected: a bad measurement re-runs the \emph{same} circuit; only a bad base
$a$ forces rebuilding, and that happens with probability $\le1/2$. That bound is close to tight,
and easy to see for yourself --- at $N=33$, of the $19$ nontrivial bases coprime to $33$, nine
fail: four have odd order ($r=5$, for $a=4,16,25,31$) and five hit $a^{r/2}\equiv-1$, where
$\gcd(a^{r/2}{-}1,N)$ collapses to $1$ ($a=2,8,17,29,32$). Ten succeed. So $47\%$ of bases send you
back to the top of the loop with a \emph{correct} order in hand --- order finding worked, and the
reduction still failed. Budget for it: this is the common case, not an error.

\section{Verifying it --- level by level}\label{im:sec:test}

This is the part most write-ups omit, and it is where the time actually goes. Each level is testable
in isolation on a simulator, and \textbf{a bug at level $k$ is invisible at level $k+1$ except as
noise}, so test in order. Prepare a basis state, run one primitive, measure, compare with Python.

\begin{lstlisting}
def set_value(qc, val, qubits):              # prepare a computational basis state
    for i, q in enumerate(qubits):
        if (val >> i) & 1:
            qc.x(q)
\end{lstlisting}

The suite below was run against the listings in this note; the results are its actual output.

\begin{center}\small
\begin{tabular}{@{}lp{0.42\textwidth}p{0.42\textwidth}@{}}
\toprule
Level & Test & Result\\
\midrule
0 & \code{qft}: unitary vs.\ the definition and vs.\ Qiskit's \code{QFTGate}, $n=1..6$
  & exact agreement; \code{.inverse()} undoes it\\[2pt]
0 & \code{\_maj}/\code{\_uma} truth tables, all $8$ inputs (\S\ref{im:sec:rc})
  & $\mathrm{MAJ}=(c{\oplus}a,\,b{\oplus}a,\,\mathrm{MAJ})$ and
    $\mathrm{UMA}\!\cdot\!\mathrm{MAJ}=(c,\,a{\oplus}b{\oplus}c,\,a)$\\[2pt]
1 & \code{rc\_add}: all $(x,y)$ pairs, $n=2..5$, control on and off
  & sum correct, $x$ preserved, carry ancilla clean; exactly $2n$ Toffoli\\[2pt]
1 & \code{rc\_add\_const}: all $y$, negative constants, $n=3,4$
  & correct mod $2^n$; all $n{+}1$ ancillas clean\\[2pt]
2 & \code{rc\_c\_add\_mod}: all pairs $y,X<N$ for $N=9,15$, control on/off
  & matches the Fourier version exactly; sign, flag and scratch all clean\\[2pt]
4 & \code{rc\_c\_ua}: $N=15$, all coprime $A$, all $x$, control on/off
  & $x\mapsto Ax\bmod15$; all scratch clean --- the tower is adder-agnostic\\[2pt]
0 & \code{lookup}: every address, $m=1,2,3$ (\S\ref{im:sec:window})
  & $T[\text{addr}]$ correct; address preserved, unary ancillas clean\\[2pt]
2 & \code{add\_quantum\_mod}: all $(y,t)$ pairs, $N=9,15$
  & $(y{+}t)\bmod N$; addend preserved, sign and flag clean\\[2pt]
3 & \code{windowed\_c\_mult\_acc}: $N=9$, $w=1,2$, all $x$, control on/off
  & agrees with \code{c\_mult\_acc} on every input; all scratch clean\\[2pt]
3 & \code{run\_windowed\_acc\_mbu} (\S\ref{im:sec:mbu}): $N=15$, all $k$, all $x$, control on/off
  & every lookup uncomputed by measurement; $x$ preserved, $y=kx\bmod15$, scratch clean, $180$ mid-circuit measurements\\[2pt]
0 & lookup uncomputed by measurement on a \emph{random non-uniform} address superposition
  & address restored exactly; leaked norm $0$\\[2pt]
0 & \code{lookup\_ui} (unary iteration, \S\ref{im:sec:ui}): every address, both control values, $m=1..4$
  & correct; exactly $2(L{-}1)$ Toffoli; $m$ ancillas not $2^m$\\[2pt]
0 & \code{phase\_fixup} (\S\ref{im:sec:mbu}): every address, every split $l=0..m$, $m=2,3,4$
  & acts as exactly $\mathrm{diag}((-1)^{F[a]})$; ancillas clean\\[2pt]
0 & $X$-basis measurement identity, $L=4,8$
  & every outcome leaves exactly $(-1)^{\mathrm{parity}(m\wedge T[a])}$; output disentangled\\[2pt]
5 & \code{windowed\_order\_circuit(7,15,w=2,t=8)}, $24$ qubits
  & peaks at $\{0,64,128,192\}$, all $2048$ shots; $r=4$; $15=3\times5$ ($58$\,s)\\[2pt]
5 & \code{windowed\_order\_circuit}, $N{=}21$, $w{=}2$, $t{=}2n{=}10$, $29$ qubits
  & six peaks at multiples of $170.67$ (smeared, $6\nmid2^t$); $r=6$; $21=7\times3$ ($4\,062$\,s)\\[2pt]
5 & \code{windowed\_exponentiate} (nested, \S\ref{im:sec:nested}): every $e<16$, $N=15$, $k=7$
  & $\ket e\ket1\ket0\mapsto\ket e\ket{7^e\bmod15}\ket0$; scratch clean\\[2pt]
5 & \code{nested\_order\_circuit}, $N{=}15$, $w_e{=}w_m{=}2$, $23$ qubits
  & $2$ multiplications not $4$; $r=4$; $15=3\times5$ ($28$\,s)\\[2pt]
1 & \code{add\_const}: all $y$, several $X$ incl.\ negative, $n=3,4$
  & single outcome $=(y+X)\bmod2^n$ every time\\[2pt]
2 & \code{c\_add\_mod}: all $y,X<N$ for $N=9,15$, control on \emph{and} off; measure sign and flag too
  & $(y+X)\bmod N$ with control on, identity with it off; \textbf{sign and flag always $\ket{00}$}\\[2pt]
4 & \code{c\_ua}: all $x<N$, all $A\in\{2,4,7,8,11,13\}$, $N=15$, both control values
  & $x\mapsto Ax\bmod N$; \textbf{all $n+2$ scratch qubits back to $\ket0$}\\[2pt]
5 & \code{order\_circuit(7,15)}, $t=8$, $4096$ shots
  & peaks at exactly $\{0,64,128,192\}$, all $4096$ shots on peaks\\[2pt]
5 & post-processing & $r=4$ for both $a=7$ and $a=2$; \code{find\_factor(15)} returns $5$\\
\bottomrule
\end{tabular}
\end{center}

\medskip\noindent
$N=15$ is too friendly to prove much, so the same suite was run at \textbf{$N=21$} ($n=5$,
$4n+2=22$ qubits, orders $2$, $3$ and $6$ --- and $6\nmid2^t$, so continued fractions do real
work):

\begin{center}\small
\begin{tabular}{@{}lp{0.42\textwidth}p{0.42\textwidth}@{}}
\toprule
Level & Test & Result\\
\midrule
2 & \code{c\_add\_mod}: all $441$ pairs $y,X<21$, control on and off
  & correct, ancillas clean, every case ($7$\,s)\\[2pt]
4 & \code{c\_ua}: all $11$ bases coprime to $21$, all $x<21$, both control values
  & $x\mapsto Ax\bmod21$; scratch clean throughout ($88$\,s)\\[2pt]
5 & \code{order\_circuit(2,21)}, $t=10$, $22$ qubits, depth ${\sim}23\,000$
  & six peaks at $0,171,341,512,683,853$ as predicted; $r=6$; $\gcd(2^3{-}1,21)=7$ \ $\Rightarrow 21=7\times3$ ($63$\,s)\\[2pt]
5 & \code{order\_circuit(8,21)} ($r=2$, exact)
  & two peaks, $1049$ and $999$ of $2048$ shots; $r=2$; again $21=7\times3$\\
\midrule
\multicolumn{3}{@{}l}{\emph{and once more at $N=33$ ($n=6$, $26$ qubits, depth ${\sim}43\,700$)}}\\[2pt]
2 & \code{c\_add\_mod}: all $1089$ pairs $y,X<33$, control on/off & correct, ancillas clean ($24$\,s)\\[2pt]
4 & \code{c\_ua}: all $19$ bases coprime to $33$, all $x<33$ & correct, scratch clean ($575$\,s)\\[2pt]
5 & \code{order\_circuit(5,33)}, $t=12$, $2048$ shots
  & ten peaks at every multiple of $409.6$; $r=10$; $\gcd(5^5{-}1,33)=11$ \ $\Rightarrow 33=11\times3$ ($1373$\,s)\\[2pt]
5 & \code{order\_circuit(10,33)} ($r=2$)
  & $r=2$; $\gcd(10{-}1,33)=3$ \ $\Rightarrow 33=3\times11$ ($1511$\,s)\\
\bottomrule
\end{tabular}
\end{center}

\noindent
The $a=2$ run is the one worth staring at. The peaks sit at $2^{10}s/6=170.67s$, which is
\emph{not} an integer, so the distribution smears onto neighbours ($682$ and $854$ pick up $68$ and
$65$ shots beside the main $683$ and $853$) --- exactly the behaviour $N=15$ can never show you,
and exactly what the continued-fraction step exists to absorb.

\begin{keybox}[title={The two assertions that matter most}]
\textbf{Assert the ancillas are clean.} At Levels 2 and 4, measure the scratch qubits along with
the answer and assert they read $\ket{0\cdots0}$. A circuit can compute the right value and still
be broken; this is the only cheap way to catch it before it becomes noise.

\textbf{Assert the peaks, not the factors.} At $t=8$, $r=4$, the measurement must concentrate on
$y=2^ts/r\in\{0,64,128,192\}$. Do not assert ``\code{find\_factor} returned $3$ or $5$'': the
classical tail verifies its own candidates, so it can succeed by chance on a \emph{uniform}
distribution --- exactly what a garbage bug produces. A published from-scratch notebook has this
failure: flat histogram over all sixteen outcomes, retry loop still prints $3\times5$. The
algorithm's robustness masks the bug. Test the distribution.
\end{keybox}

\begin{keybox}[title={Pitfalls, in the order you will hit them}]
\begin{itemize}\setlength\itemsep{1pt}
\item \textbf{Endianness.} Qiskit is little-endian: \code{qr[0]} is the LSB, and printed bitstrings
read MSB-first. Half of all bugs here are a reversed register.
\item \textbf{Phase sign and scale.} \code{2*np.pi*X*2**(i-n)} --- get the exponent wrong and you
add the wrong constant, silently. Test \code{add\_const} exhaustively; it takes seconds.
\item \textbf{Negative exponents.} Plain Python is safe (\code{2 ** -3} is \code{0.125}), but if
$n$ or $i$ arrives as a \emph{numpy} integer --- easy, if a bit width came from an array ---
\code{2 ** (i-n)} raises \code{ValueError: Integers to negative integer powers are not allowed},
and \code{1 << (i-n)} raises on any type. Write \code{2.0 ** (i - n)} and stop thinking about it.
\item \textbf{The missing $+1$.} The value register must be $n+1$ qubits (sign qubit included)
everywhere in \S\ref{im:sec:l2}. Drop it and the comparison silently reads a data bit.
\item \textbf{Reducing the constant.} Pass $X$ with $0\le X<N$ into \code{c\_add\_mod} ---
hence \code{(A << i) \% N} and \code{(-B) \% N} at Level~3/4.
\item \textbf{Simulation size.} $4n+2$ qubits at $N=15$ is $18$ --- fine. $N=21$ is $22$ qubits at
depth ${\sim}23\,000$, about a minute per run on a laptop; $N=33$ is $26$ and roughly $16\times$
worse. This is a wall, not a slope. Test the arithmetic levels with
\code{Statevector.from\_instruction} instead of sampling --- they are permutations on basis states,
so one exact evolution replaces thousands of shots, and the $N=21$ sweeps above finish in seconds.
\item \textbf{Transpile once.} Run the pass manager on the circuit once and reuse it across shots.
\end{itemize}
\end{keybox}

\subsection*{Running it}

\begin{lstlisting}
from qiskit.transpiler.preset_passmanagers import generate_preset_pass_manager
from qiskit_aer import AerSimulator
from qiskit_aer.primitives import SamplerV2 as Sampler

SIM, SAMPLER = AerSimulator(), Sampler(seed=7)
PM = generate_preset_pass_manager(backend=SIM, optimization_level=1)

def run(qc, shots=4096):
    return SAMPLER.run([PM.run(qc)], shots=shots).result()[0].data.out.get_counts()

counts = run(order_circuit(7, 15))           # -> peaks at 0, 64, 128, 192
print(order_from_counts(counts, 7, 15, t=8)) # -> 4
print(find_factor(15, run))                  # -> 3 or 5
\end{lstlisting}


\part*{Part V --- Making it fast}
\addcontentsline{toc}{part}{Part V --- Making it fast}

\noindent
Part~IV's implementation is correct and about as simple as the algorithm gets. It is also slower
than it needs to be in three separate ways, each with its own literature and its own trade. This
part builds all three and measures them: ripple-carry arithmetic (a different gate set --- the one
error correction actually wants), windowing (fewer additions, paid for in qubits), and the two
refinements, unary-iteration lookups and measurement-based uncomputation, that make windowing worth
having.

\section{The other Level 1: ripple-carry arithmetic}\label{im:sec:rc}

Everything so far rests on phase gates, and that is a liability. The smallest rotation in an
$n$-qubit Fourier addition is $2\pi/2^{n+1}$: $0.20$ rad at $n=4$, $10^{-9}$ at $n=32$, and
$10^{-616}$ at $n=2048$ --- a number no apparatus will ever apply. Worse, arbitrary-angle rotations
are not fault-tolerant gates; each must be synthesised from Clifford$+T$ at a cost that grows with
the precision demanded.

The alternative is to do addition the way hardware wants it: carries, in the computational basis,
using only $X$, $\CNOT$ and Toffoli. Then the \emph{Toffoli count} is the cost, because it is a
proxy for the $T$ count, and $T$ gates are what magic-state distillation actually bills you for.
This section builds Level 1 a second time. \textbf{Nothing above Level 1 changes.}

\subsection{Two blocks that do everything}

Schoolbook addition is $s_i=a_i\oplus b_i\oplus c_i$ with carry
$c_{i+1}=\mathrm{MAJ}(a_i,b_i,c_i)=a_ib_i\oplus a_ic_i\oplus b_ic_i$. Read that as a cost statement:
the sum bits are $\mathbb F_2$-linear, so they are CNOTs and effectively free; the carries are
quadratic, so each one costs a Toffoli; and they form a chain, so depth is $\Theta(n)$.

Cuccaro, Draper, Kutin and Moulton \cite{cuccaro} express the whole adder with two three-qubit
blocks acting on a triple $(c,b,a)$ --- carry-in, addend, augend:
\begin{align*}
\mathrm{MAJ}&:\ \CNOT(a\!\to\!b),\ \CNOT(a\!\to\!c),\ \mathrm{TOF}(c,b\!\to\!a);\\
\mathrm{UMA}&:\ \mathrm{TOF}(c,b\!\to\!a),\ \CNOT(a\!\to\!c),\ \CNOT(c\!\to\!b).
\end{align*}
$\mathrm{MAJ}$ sends $(c,b,a)\mapsto(c\oplus a,\ b\oplus a,\ \mathrm{MAJ}(a,b,c))$: it computes the
carry \emph{into the wire that held $a$}, which is why no ancilla per carry is needed. Running the
chain forward computes all $n$ carries; running $\mathrm{UMA}$ back down undoes them. And here is
the point of the design:
\[
\mathrm{UMA}\cdot\mathrm{MAJ}:(c,b,a)\mapsto(c,\ a\oplus b\oplus c,\ a)
\]
--- \textbf{$\mathrm{UMA}$ erases the carry and writes the sum bit in the same three gates}. The
uncomputation that \S\ref{im:sec:l1}'s no-garbage rule makes compulsory is exactly the pass that
delivers the answer. You are not paying twice; you are paying once and collecting on the way out.

\begin{lstlisting}
def _maj(qc, c, b, a, ctrls=()):
    """MAJ: |c,b,a> -> |c^a, b^a, MAJ(a,b,c)>.  Only the b-target takes ctrls."""
    qc.mcx([*ctrls, a], b)               # plain CNOT when ctrls is empty
    qc.cx(a, c)
    qc.ccx(c, b, a)

def _uma(qc, c, b, a, ctrls=()):
    """UMA: undoes MAJ's carry and writes the sum bit into b."""
    qc.ccx(c, b, a)
    qc.cx(a, c)
    qc.mcx([*ctrls, c], b)               # plain CNOT when ctrls is empty

def rc_add(qc, x, y, carry, ctrls=()):
    """|x>|y> -> |x>|(y + x) mod 2^n>.  carry: one clean ancilla, returned |0>."""
    n = len(y)
    _maj(qc, carry, y[0], x[0], ctrls)
    for i in range(n - 1):
        _maj(qc, x[i], y[i + 1], x[i + 1], ctrls)
    for i in range(n - 1, 0, -1):
        _uma(qc, x[i - 1], y[i], x[i], ctrls)
    _uma(qc, carry, y[0], x[0], ctrls)
\end{lstlisting}

\noindent
That is $2n$ Toffolis, $4n$ CNOTs and one ancilla, with $x$ preserved and $y$ overwritten by the
sum --- verified exhaustively below. The interleaved wire order you see in the literature
($c_0,b_0,a_0,b_1,a_1,\dots$) is only a seating plan that keeps each block's three qubits adjacent;
it means nothing mathematically.

\subsection{Adding a classical constant}

CDKM adds two \emph{quantum} registers, but Level 2 wants ``add the known constant $X$''. The
cheapest honest fix: write $X$ into $n$ clean ancillas with $X$ gates, run the adder, and unwrite
it. That is where the ripple-carry route's extra qubits go.

\begin{lstlisting}
def rc_add_const(qc, X, y, anc, ctrls=()):
    """|y> -> |y + X mod 2^n>.  anc: n+1 clean ancillas, returned clean."""
    n = len(y)
    carry, xreg = anc[0], anc[1:n + 1]
    Xm = X % (2**n)                      # negative X just works, as in Fourier
    for i in range(n):                   # encode the constant -- uncontrolled
        if (Xm >> i) & 1:
            qc.x(xreg[i])
    rc_add(qc, xreg, y, carry, ctrls)
    for i in range(n):                   # unencode
        if (Xm >> i) & 1:
            qc.x(xreg[i])
\end{lstlisting}

\noindent
Negative constants work for the same reason as in \S\ref{im:sec:l1} --- \code{X \% 2**n} is the two's
complement --- so subtraction again needs no separate circuit. With this in place the modular adder
of \S\ref{im:sec:l2} is reproduced \emph{line for line}, only the adder calls changing:

\begin{lstlisting}
def rc_c_add_mod(qc, ctrls, X, y, sign, flag, N, anc):
    yw = list(y) + [sign]
    rc_add_const(qc, X, yw, anc, ctrls)  # (1) + X          (controlled)
    rc_add_const(qc, -N, yw, anc)        # (2) - N
    qc.cx(sign, flag)                     # (3) sign -> flag
    rc_add_const(qc, N, yw, anc, [flag])  # (4) + N  if flag
    rc_add_const(qc, -X, yw, anc)        # (5) - X
    qc.mcx(list(ctrls) + [sign], flag)    # (6) clear flag from the answer
    qc.x(flag)
    rc_add_const(qc, X, yw, anc)         # (7) + X
\end{lstlisting}

And Levels 3--5 are then genuinely unchanged --- same loops, same constants, same modular inverse,
with only the extra ancilla register threaded through:

\begin{lstlisting}
def rc_c_mult_acc(qc, c, A, x, y, sign, flag, N, anc):
    for i in range(len(x)):
        rc_c_add_mod(qc, [c, x[i]], (A << i) % N, y, sign, flag, N, anc)

def rc_c_ua(qc, c, A, x, y, sign, flag, N, anc):
    rc_c_mult_acc(qc, c, A % N, x, y, sign, flag, N, anc)
    for i in range(len(x)):
        qc.cswap(c, x[i], y[i])
    B = pow(A, -1, N)                        # classical inverse, unchanged
    rc_c_mult_acc(qc, c, (-B) % N, x, y, sign, flag, N, anc)
\end{lstlisting}

\noindent
Checked exhaustively at $N=15$ ($17$ qubits): every base coprime to $15$, every $x<15$, control on
and off, $x\mapsto Ax\bmod15$ correct and all $n{+}4$ scratch qubits returned to $\ket0$. The
abstraction really does hold --- swap the adder, keep the tower.

\begin{keybox}[title={The control trick, again --- and why it telescopes}]
Look at \code{\_maj} and \code{\_uma}: of their six gates, exactly \emph{two} accept the outer
control, and both are the CNOT targeting the $b$ wire. Everything else is unconditional. The
justification is the residue argument of \S\ref{im:sec:l2}, and here it is unusually clean.

With the control off, what survives of $\mathrm{MAJ}$ is $[\CNOT(a\!\to\!c),\
\mathrm{TOF}(c,b\!\to\!a)]$ and what survives of $\mathrm{UMA}$ is
$[\mathrm{TOF}(c,b\!\to\!a),\ \CNOT(a\!\to\!c)]$ --- the same two gates in the opposite order, i.e.\
each other's inverse. Since the $\mathrm{UMA}$ chain retraces the $\mathrm{MAJ}$ chain in reverse,
the residues cancel pairwise from the inside out and the whole adder telescopes to the identity.

So a controlled $n$-bit addition costs $2n$ controlled gates rather than $6n$ --- and the constant
encoding stays uncontrolled too, since encode/unencode brackets an identity. Verified exhaustively
for $n=2..5$ with the control both on and off.
\end{keybox}

\subsection{What it costs, measured}

Both backends, the same doubly-usable modular adder, same modulus, counted before and after
transpiling to $\{\CNOT,U\}$:

\begin{center}\small
\begin{tabular}{@{}llrrrrr@{}}
\toprule
$N$ & backend & qubits & rotations & Toffoli & CNOT & depth\\
\midrule
$15$ ($n{=}4$) & Fourier      & $7$  & $125$ & $0$   & $287$ & $492$\\
               & ripple-carry & $13$ & $\mathbf 0$ & $71$ & $437$ & $578$\\
\midrule
$63$ ($n{=}6$) & Fourier      & $9$  & $245$ & $0$   & $545$ & $948$\\
               & ripple-carry & $17$ & $\mathbf 0$ & $99$ & $637$ & $\mathbf{844}$\\
\bottomrule
\end{tabular}
\end{center}

\noindent
Read the trend, not the absolute numbers. Fourier rotations grow as $O(n^2)$ per addition
($125\to245$) because of the QFT sandwich; ripple-carry Toffolis grow as $O(n)$ ($71\to99$). By
$n=6$ the ripple-carry circuit is already \emph{shallower}, and the gap widens from there. What you
buy is a factor of $n$ in gate count and a gate set that survives error correction; what you pay is
$n+2$ extra qubits, taking the full order-finding circuit from $4n+2$ to $5n+4$.

\begin{keybox}[title={Which to use}]
For \emph{simulating} small $N$ on a laptop, Fourier wins: fewer qubits, and simulation cost is
exponential in qubit count, so $4n+2$ versus $5n+4$ is the whole ballgame.

For anything aimed at real fault-tolerant hardware, ripple-carry wins and it is not close. Zero
arbitrary-angle rotations, a gate count lower by $O(n)$, and a cost model ($T$ count) that actually
predicts the resource bill. The $10^{-616}$ rad rotation is not a rounding problem you can engineer
around --- it is a statement that Fourier arithmetic at cryptographic sizes must be approximated
before it can be attempted at all.
\end{keybox}

\section{Windowing: trading lookups for additions}\label{im:sec:window}

Level 3 spends one modular addition per bit of $x$ --- $n$ of them, each controlled twice. Gidney's
windowing \cite{gidney-window} removes most of them by the classical trick of merging iterations
into a table lookup. Unlike the control placement of \S\ref{im:sec:l2}, which is a fixed constant
factor, this one \emph{improves with $n$} --- it is the optimisation that changes the exponent's
logarithm, and the reason published Toffoli counts are what they are.

\subsection{The idea in three steps}

\textbf{Step 1: a controlled addition is a lookup.} Adding the constant $k$ into $y$ controlled on
one qubit $q$ is the same as adding $T[q]$ where $T=[0,k]$ --- a two-entry table addressed by $q$.
Nothing has changed yet.

\textbf{Step 2: address more than one qubit.} If a two-entry table works, so does a $2^w$-entry one.
Take $w$ bits of $x$ at a time; the value they contribute to $kx$ is
$k\cdot(\text{window value})\cdot2^i \bmod N$, which is classical --- so tabulate all $2^w$
possibilities, look up the one the register actually holds, add it \emph{once}, then uncompute the
lookup. One addition per \emph{window} instead of one per bit: $\lceil n/w\rceil$ instead of $n$.

\textbf{Step 3: fold the control into the address.} The outer control $c$ becomes address bit $0$,
with every $c=0$ entry set to $0$ (Gidney Fig.~5). Adding zero modulo $N$ is the identity, so the
modular addition needs \emph{no control at all}. This is why windowing composes so well with the
ladder: the expensive part stops being controlled.

The cost of a window is one lookup instead of $w-1$ extra additions, so the Toffoli count goes as
$O\!\big(\tfrac nw(n+2^w)\big)$. Setting $w=\lg n$ balances the two terms and gives
$O(n^2/\lg n)$ --- a logarithmic factor better than schoolbook.

\subsection{What has to be built}

Two new primitives. First a \textbf{QROM}: read a classical table at a quantum address. Encode the
address one-hot, and the data load becomes pure CNOTs --- one per set bit of the table.

\begin{lstlisting}
def _unary_gate(m):
    """binary(m qubits) -> one-hot(2^m qubits), using 2^m - 1 Fredkins."""
    addr, unary = QuantumRegister(m, "a"), QuantumRegister(1 << m, "u")
    qc = QuantumCircuit(addr, unary)
    qc.x(unary[0])
    for i in reversed(range(m)):                  # MSB down to LSB
        step = 1 << i
        for j in range(0, 1 << m, 2 * step):
            qc.cswap(addr[i], unary[j], unary[j + step])
    return qc.to_gate(label="unary")

def lookup(qc, addr, out, table, unary):
    """|addr>|out> -> |addr>|out ^ table[addr]>; unary ancillas returned clean."""
    g = _unary_gate(len(addr))
    qc.append(g, list(addr) + list(unary))
    for j in range(len(unary)):
        Tj = table[j] if j < len(table) else 0
        for b in range(len(out)):
            if (Tj >> b) & 1:
                qc.cx(unary[j], out[b])
    qc.append(g.inverse(), list(addr) + list(unary))
\end{lstlisting}

\noindent
Second, the lookup produces a \emph{quantum} value, so Level 2 now needs to add a register rather
than a constant. That is a smaller change than it sounds: the phase ramp of \S\ref{im:sec:l1} just
takes its rotations from the bits of $t$ instead of the bits of a Python integer, and Beauregard's
seven blocks are otherwise untouched.

\begin{lstlisting}
def add_quantum(qc, a, y, sign=+1):
    """y += sign*a (mod 2^len(y)); a and y both quantum.  The quantum twin of
    add_const: same phase ramp, each rotation controlled by a bit of a."""
    n = len(y)
    qc.append(qft(n), y)
    for i in range(len(a)):
        for j in range(n):
            if i + j < n:                          # higher terms are trivial
                qc.cp(sign * 2 * np.pi * 2.0 ** (i + j - n), a[i], y[j])
    qc.append(qft(n).inverse(), y)

def add_quantum_mod(qc, t, y, sign, flag, N):
    """|y>|t> -> |(y+t) mod N>|t>.  Blocks (1),(5),(7) go quantum; rest as-is."""
    yw = list(y) + [sign]
    add_quantum(qc, t, yw, +1)            # (1) + t
    add_const(qc, -N, yw)                 # (2) - N
    qc.cx(sign, flag)                     # (3) sign -> flag
    c_add_const(qc, [flag], N, yw)        # (4) + N  if flag
    add_quantum(qc, t, yw, -1)            # (5) - t
    qc.cx(sign, flag)                     # (6) [R < t] == not flag, so this
    qc.x(flag)                            #     pair returns flag to |0>
    add_quantum(qc, t, yw, +1)            # (7) + t
\end{lstlisting}

\noindent
With those, Level 3 becomes a loop over windows rather than over bits. It calls \code{lookup\_ui} rather than the \code{lookup} above: same operation, but built by unary iteration so it needs $w$ ancillas instead of $2^w$ --- see \S\ref{im:sec:ui}, which is worth reading before you build anything. The one-hot version above is the easier way to see what a QROM \emph{is}; it is not the one to ship.

\begin{lstlisting}
def windowed_c_mult_acc(qc, c, k, x, y, sign, flag, N, t, anc, w):
    """Drop-in for c_mult_acc, but ceil(len(x)/w) modular additions instead of
    len(x).  c is the top-level control of each lookup, so when c=0 no data
    loads at all and the modular addition -- which is uncontrolled -- adds 0.

    t:   len(y) clean qubits holding one looked-up value (returned clean)
    anc: w clean qubits for the unary-iteration QROM     (returned clean)
    w:   window size -- bits of x consumed per lookup
    """
    n = len(y)
    for i in range(0, len(x), w):
        win = list(x[i:i + w])
        # window value j contributes (k * j * 2^i) mod N
        table = [((k * j) << i) % N for j in range(2**len(win))]
        lookup_ui(qc, c, win, t[:n], table, anc[:len(win)])
        add_quantum_mod(qc, t[:n], y, sign, flag, N)
        lookup_ui(qc, c, win, t[:n], table, anc[:len(win)])   # uncompute
\end{lstlisting}

\noindent
Level 4 and the ladder then need no new ideas at all --- same multiply/swap/uncompute, same
classical inverse, with the lookup scratch threaded through:

\begin{lstlisting}
def windowed_c_ua(qc, c, A, x, y, sign, flag, N, t, anc, w):
    """|c>|x>|0> -> |c>|A^c x mod N>|0>.  Same multiply/swap/uncompute as
    c_ua; only the accumulate is windowed."""
    windowed_c_mult_acc(qc, c, A % N, x, y, sign, flag, N, t, anc, w)
    for i in range(len(x)):
        qc.cswap(c, x[i], y[i])
    B = pow(A, -1, N)                        # classical inverse, unchanged
    windowed_c_mult_acc(qc, c, (-B) % N, x, y, sign, flag, N, t, anc, w)
\end{lstlisting}

\begin{keybox}[title={Why the $c=0$ entries being zero is the whole trick}]
Set $c=0$ and every table entry is $0$, so the lookup writes $\ket0$ into $t$ and block (1) adds
nothing. Trace the seven blocks with $t=0$: (2) drives the register negative, so the sign bit sets
and (3) sets the flag; (4) adds $N$ back; (5) subtracts nothing, leaving $R=y\ge0$ so the sign bit
is now $0$; and (6) computes $\text{flag}\oplus\text{sign}=1$, which the $X$ clears. Identity,
ancilla clean --- \emph{without a single controlled gate in the adder}. The control survives only
as a classical fact about the table.
\end{keybox}

\subsection{What it costs, measured}

One accumulate $y \mathbin{+{=}} c\!\cdot\!kx \bmod N$, transpiled to $\{\CNOT,U\}$:

\begin{center}\small
\begin{tabular}{@{}lrrrrr@{}}
\toprule
method & mod.\ adds & qubits & CNOT & depth & vs.\ depth\\
\midrule
\multicolumn{6}{@{}l}{\emph{$N=15$, $n=4$}}\\
shift-and-add        & $4$ & $11$ & $1\,254$ & $2\,119$ & ---\\
window $w=2$         & $2$ & $23$ & $1\,184$ & $1\,128$ & $1.9\times$\\
\midrule
\multicolumn{6}{@{}l}{\emph{$N=63$, $n=6$}}\\
shift-and-add        & $6$ & $15$ & $3\,474$ & $6\,001$ & ---\\
window $w=3$         & $2$ & $37$ & $2\,412$ & $2\,082$ & $2.9\times$\\
\midrule
\multicolumn{6}{@{}l}{\emph{$N=255$, $n=8$}}\\
shift-and-add        & $8$ & $19$ & $7\,390$ & $13\,020$ & ---\\
window $w=3$         & $3$ & $43$ & $4\,646$ & $4\,873$ & $2.7\times$\\
window $w=4$         & $2$ & $59$ & $4\,432$ & $3\,428$ & $3.8\times$\\
\bottomrule
\end{tabular}
\end{center}

\noindent
The addition count falls exactly as $\lceil n/w\rceil$ predicts, and --- unlike many asymptotic
optimisations --- the win is already real at $n=4$ and grows: $1.06\times$, $1.44\times$,
$1.67\times$ in CNOTs at $n=4,6,8$, and far more in depth, because the lookups for one window are
shallow and the additions they replace were not. What you pay is qubits.

The same holds for the \emph{whole} order-finding circuit, not just one accumulate. At $N=15$,
with the unary-iteration QROM of \S\ref{im:sec:ui}, windowing halves the depth for six extra qubits:

\begin{center}\small
\begin{tabular}{@{}lrrr@{}}
\toprule
full circuit, $N=15$ & qubits & CNOT & depth\\
\midrule
plain, $t=4$                          & $14$ & $10\,212$ & $17\,239$\\
windowed $w=2$, $t=4$                 & $20$ & $7\,122$  & $\mathbf{8\,488}$\\
plain, $t=2n=8$                       & $18$ & $20\,462$ & $34\,537$\\
windowed $w=2$, $t=2n=8$              & $24$ & $14\,276$ & $\mathbf{17\,028}$\\
\bottomrule
\end{tabular}
\end{center}

\subsection{End to end}

Windowed arithmetic factors $15$ with nothing else changed. With the unary-iteration QROM of
\S\ref{im:sec:ui} the full-precision circuit --- $t=2n=8$, $w=2$ --- fits in $24$ qubits and gives a
textbook distribution:

\begin{center}\small
\begin{tabular}{@{}rrrr@{}}
\toprule
$y$ & shots & $y/2^t$ & continued fraction\\
\midrule
$0$   & $530$ & $0$   & ---\\
$64$  & $490$ & $1/4$ & $r=4$ \ \checkmark\\
$128$ & $508$ & $1/2$ & $r=2$, rejected\\
$192$ & $520$ & $3/4$ & $r=4$ \ \checkmark\\
\bottomrule
\end{tabular}
\end{center}

\noindent
All $2048$ shots land on the four ideal peaks, $r=4$ is recovered, and $\gcd(7^2{-}1,15)=3$ gives
$15=3\times5$; base $a=2$ likewise. Note that even in this clean run the $y=128$ peak returns
$r=2$ and is discarded --- \S\ref{im:sec:classical}'s scan is doing work here too.

It is worth seeing what \S\ref{im:sec:ui} bought at this level. Built on the one-hot QROM the same
circuit needed $30$ qubits at $t=2n$ --- $16$\,GB of statevector, out of reach --- so it had to be
run at $t=4$ instead, which is sound at $N=15$ only because $r=4$ divides $2^4$. Unary iteration
turned a compromise into the real thing, and made it faster: $58$\,s against $100$\,s.

$N=15$ is too forgiving to prove much, so the windowed circuit was also run on $N=21$, where the
order is $6$ and $6\nmid2^t$, so the continued-fraction step has to do real work. At the full
$t=2n=10$ ($29$ qubits, depth $14\,510$, $8$\,GB of statevector, $68$ minutes) it lands on all six
ideal peaks $2^{10}s/6=170.67\,s$:

\begin{center}\small
\begin{tabular}{@{}rrrrrrrrr@{}}
\toprule
$y$ & $0$ & $171$ & $341$ & $512$ & $683$ & $853$ & $682$ & $854$\\
shots & $326$ & $223$ & $247$ & $337$ & $221$ & $234$ & $74$ & $72$\\
$\to$ & --- & $1/6$ & $1/3$ & $1/2$ & $2/3$ & $5/6$ & $2/3$ & $5/6$\\
\bottomrule
\end{tabular}
\end{center}

\noindent
--- the last two columns being the smearing onto neighbours that a non-integer $2^t/r$ forces, and
which $N=15$ can never show. It recovers $r=6$, and $\gcd(2^3{-}1,21)=7$ gives $21=7\times3$. Only
$1/6$ and $5/6$ pass the \code{pow} check (\S\ref{im:sec:classical}); the other four peaks return
divisors of $6$ and are discarded. Windowed arithmetic is not doing anything special at $N=15$; it
is just doing arithmetic.

Nested windowing (\S\ref{im:sec:nested}) factors $15$ too, on $23$ qubits at $t=4$, using
\emph{two} modular multiplications where the plain ladder uses four: peaks at $\{0,4,8,12\}$,
$r=4$, $15=3\times5$, for both $a=7$ and $a=2$.

\subsection{Making the lookup cheap: unary iteration}\label{im:sec:ui}

The one-hot QROM above is correct but wasteful: it materialises all $2^m$ address states at once,
so it needs $2^m$ ancillas. Gidney's Fig.~2 walks the address \emph{tree} instead, holding only
the AND of the current prefix \cite{babbush}. Recursively: split on the top address bit, compute
$\text{node}=\text{ctrl}\wedge a$ to select the upper half, flip it with one CNOT to get
$\text{ctrl}\wedge\neg a$ for the lower half, then clear it.

\begin{lstlisting}
def _walk(qc, ctrl, addr, out, table, offset, anc):
    """out ^= table[offset + value(addr)], all conditioned on ctrl."""
    if not addr:
        T = table[offset] if offset < len(table) else 0
        for b in range(len(out)):
            if (T >> b) & 1:
                qc.cx(ctrl, out[b])           # leaf: data load is CNOTs
        return
    a, rest = addr[-1], addr[:-1]             # split on the most significant bit
    half, node = 1 << len(rest), anc[0]
    qc.ccx(ctrl, a, node)                     # node = ctrl AND a      (upper half)
    _walk(qc, node, rest, out, table, offset + half, anc[1:])
    qc.cx(ctrl, node)                         # node = ctrl AND NOT a  (lower half)
    _walk(qc, node, rest, out, table, offset, anc[1:])
    qc.cx(ctrl, node)                         # put node back to ctrl AND a ...
    qc.ccx(ctrl, a, node)                     # ... and clear it

def lookup_ui(qc, ctrl, addr, out, table, anc):
    """|addr>|out> -> |addr>|out ^ table[addr]> when ctrl is 1, else identity.
    anc: len(addr) clean ancillas, returned clean.  2(L-1) Toffolis."""
    _walk(qc, ctrl, list(addr), list(out), list(table), 0, list(anc))
\end{lstlisting}

\noindent
The recursion depth is $m$, so it needs $m$ ancillas, not $2^m$ --- and measured Toffoli counts
land exactly on $2(L-1)$: $2,6,14,30$ for $m=1,2,3,4$, against $2,4,8,16$ ancillas saved down to
$1,2,3,4$. (With the temporary-AND trick \cite{gidney-and} --- built in \S\ref{sec:tempand} --- the uncompute Toffolis become Clifford, giving
quoted $L-1$.) Note this also replaces the control-folding of \S\ref{im:sec:window}: $c$ is now simply
the top-level \code{ctrl}, so the tree halves as well.

Swapping it in shrinks the windowed circuit substantially --- and, unlike the one-hot version,
the qubit cost is now barely sensitive to $w$:

\begin{center}\small
\begin{tabular}{@{}lrrrr@{}}
\toprule
one accumulate, $N=255$ ($n=8$) & mod.\ adds & qubits & CNOT & depth\\
\midrule
shift-and-add               & $8$ & $19$ & $7\,390$ & $13\,020$\\
window $w=2$                & $4$ & $29$ & $4\,840$ & $6\,497$\\
window $w=3$                & $3$ & $30$ & $3\,902$ & $4\,873$\\
window $w=4$                & $2$ & $31$ & $\mathbf{3\,264}$ & $\mathbf{3\,866}$\\
\bottomrule
\end{tabular}
\end{center}

\noindent
With one-hot at $w=4$ the same circuit needed $59$ qubits; unary iteration brings it to $31$, and
the CNOT advantage over shift-and-add improves from $1.67\times$ to $2.26\times$ (depth
$3.4\times$). The earlier tables in \S\ref{im:sec:window} are the one-hot numbers; these are the ones
to build on.

\subsection{Uncomputing a lookup for free}\label{im:sec:mbu}

We still uncompute each lookup by running it again, at equal cost --- and that is avoidable,
because the output is about to be discarded. Measure every output qubit in the $X$ basis. On
$\sum_a\alpha_a\ket a\ket{T[a]}$ this yields a uniformly random outcome $m$ and leaves
\[
\sum_a (-1)^{\,\mathrm{parity}(m\,\wedge\,T[a])}\,\alpha_a\ket a\ \otimes\ \ket m ,
\]
so the output register is disentangled \emph{for free}; all that survives is a phase. Removing it
is a phase-only lookup, and \emph{that} can be done in $O(\sqrt L)$: split the address into $l$ low
bits and $h$ high bits, one-hot the low half ($2^l$ Fredkins), then unary-iterate the high half
applying $CZ$ into the one-hot register ($2(2^h{-}1)$ Toffolis). With $l=h=m/2$ the total is
$O(\sqrt L)$ instead of $O(L)$.

\begin{lstlisting}
def phase_fixup(qc, addr, F, one, unary_lo, anc_hi, l):
    """Apply the phase (-1)^F[a] to every basis state |a> of addr, in O(sqrt(L))."""
    lo, hi = list(addr[:l]), list(addr[l:])
    g = _unary_gate(l)
    qc.append(g, lo + list(unary_lo[:1 << l]))
    qc.x(one)                                        # always-true control
    _walk_cz(qc, one, hi, list(unary_lo[:1 << l]), F, 0, list(anc_hi))
    qc.x(one)
    qc.append(g.inverse(), lo + list(unary_lo[:1 << l]))
\end{lstlisting}

\noindent
Verified two ways: \code{phase\_fixup} acts as exactly $\mathrm{diag}((-1)^{F[a]})$ with all
ancillas clean, for every split $l=0\ldots m$ at $m=2,3,4$; and the measurement identity above
holds for \emph{every} outcome $m$ (checked amplitude by amplitude at $L=4,8$). At $m=4$ the
$\sqrt L$ split costs $8$ where a direct phase lookup costs $16$.

\medskip\noindent
One caveat on the measured costs elsewhere in this section: the CNOT and depth figures in
\S\ref{im:sec:window}--\ref{im:sec:nested} were taken from the \emph{static} circuits, which uncompute by
re-running the lookup. They are therefore a \textbf{lower bound} on what windowing delivers --- the
hybrid version below is cheaper still, by the factors in the table.

\subsection{Nesting: windowing the exponent too}\label{im:sec:nested}

Everything so far windows the \emph{multiplicand}. Gidney's headline construction (\S3.5 of
\cite{gidney-window}) windows the \emph{exponent} as well, and that is where the second logarithmic
factor comes from. Take $w_e$ counting qubits at a time and let the lookup select the factor
$k^{v\cdot 2^i}$ for the window value $v$; the address becomes the exponent-window bits followed by
the multiplicand-window bits, so one table serves both.

Three things fall out at once. There are $\lceil t/w_e\rceil$ multiplications instead of $t$. The
multiplications need \emph{no control at all} --- when the exponent window reads $0$ the factor is
$k^0=1$, and multiplying by $1$ is the identity. And the swap between accumulators becomes a
\emph{relabelling} in the classical loop, so the $n$ Fredkin gates of \S\ref{im:sec:l34} vanish.

\begin{lstlisting}
for i in range(0, len(e), we):
    ei = list(e[i:i + we])
    ne = len(ei)
    kes = [pow(k, (2**i) * v, N) for v in range(2**ne)]
    kes_inv = [pow(v, -1, N) for v in kes]

    # B += A * k_e   (mod N):  maps (x, 0) -> (x, x*k_e)
    for j in range(0, n, wm):
        mi = A[j:j + wm]
        addr = ei + list(mi)
        table = [((kes[ad & ((1 << ne) - 1)] * (ad >> ne)) << j) % N
                 for ad in range(1 << len(addr))]
        lookup_ui(qc, one, addr, tmp[:n], table, anc[:len(addr)])
        qc.append(madd, tmp[:n] + B + [sign, flag])
        lookup_ui(qc, one, addr, tmp[:n], table, anc[:len(addr)])

    # A -= B * k_e^{-1} (mod N): maps (x, x*k_e) -> (0, x*k_e)
    for j in range(0, n, wm):
        mi = B[j:j + wm]
        addr = ei + list(mi)
        table = [((kes_inv[ad & ((1 << ne) - 1)] * (ad >> ne)) << j) % N
                 for ad in range(1 << len(addr))]
        lookup_ui(qc, one, addr, tmp[:n], table, anc[:len(addr)])
        qc.append(msub, tmp[:n] + A + [sign, flag])
        lookup_ui(qc, one, addr, tmp[:n], table, anc[:len(addr)])

    A, B = B, A                               # relabelling swap -- free
\end{lstlisting}

\noindent
Verified on the ladder identity directly: $\ket e\ket1\ket0\mapsto\ket e\ket{k^e\bmod N}\ket0$ for
\emph{every} exponent $e<2^4$ at $N=15$, $k=7$, with all scratch returned to $\ket0$. Measured on
the full circuit at $t=2n=8$:

\begin{center}\small
\begin{tabular}{@{}lrrrr@{}}
\toprule
$N=15$, $t=2n=8$ & mults & qubits & CNOT & depth\\
\midrule
plain shift-and-add       & $8$ & $18$ & $20\,462$ & $34\,537$\\
windowed $w_m=2$          & $8$ & $24$ & $14\,276$ & $17\,028$\\
nested $w_e=2$, $w_m=2$   & $4$ & $27$ & $\mathbf{11\,908}$ & $\mathbf{14\,229}$\\
nested $w_e=4$, $w_m=2$   & $2$ & $29$ & $16\,292$ & $20\,095$\\
\bottomrule
\end{tabular}
\end{center}

\noindent
$1.72\times$ the CNOTs and $2.43\times$ the depth of plain shift-and-add, and a further $1.2\times$
on top of single-level windowing. The last row is the useful warning: at $n=4$, $w_e=4$ makes the
table $2^{6}$ entries and the lookup costs more than the two multiplications it saved. Windows
have an optimum ($w\approx\lg n$), and overshooting is worse than not windowing that level at all.

\subsection{What of the paper is actually here}\label{im:sec:coverage}

Windowing is a family of constructions, and it is worth being exact about which of them this note
builds, since partial implementations of it are easy to mistake for the whole thing.

\begin{center}\small
\begin{tabular}{@{}llp{0.40\textwidth}@{}}
\toprule
\cite{gidney-window} & status & where / why not\\
\midrule
\S2, Fig.~2 \ table lookup by unary iteration & \textbf{built} & \S\ref{im:sec:ui}, \code{lookup\_ui}\\
\S3.3 \ windowed modular product addition     & \textbf{built} & \S\ref{im:sec:window}, \code{windowed\_c\_mult\_acc}\\
\S3.4 \ windowed modular multiplication       & \textbf{built} & \S\ref{im:sec:window}, \code{windowed\_c\_ua}\\
\S3.5 \ nested exponent $+$ multiplication windows & \textbf{built} & \S\ref{im:sec:nested}, \code{windowed\_exponentiate}\\
\midrule
\S2, Fig.~3 \ measurement-based uncomputation & \textbf{built} & \S\ref{im:sec:mbu}, \code{run\_windowed\_acc\_mbu} --- hybrid, so it runs as circuit/measure/classical/circuit\\[3pt]
\S3.2 \ windowed in-place multiplication      & not built & not needed: we multiply out-of-place and swap (Fig.~6)\\
\bottomrule
\end{tabular}
\end{center}

\begin{keybox}[title={Wiring it in: the classical controller sits in the middle}]
The fixup table $F[a]=\mathrm{parity}(m\wedge T[a])$ depends on the measurement outcome, so this
cannot be one static circuit --- it is genuinely
\emph{circuit $\to$ measure $\to$ classical $\to$ circuit}, which is exactly how it runs on
hardware. \code{run\_windowed\_acc\_mbu} does that: it emits a stage, measures, builds $F$ in
Python, emits the fixup, and continues.

Verified by wiring it into the real windowed accumulate: every $k$, every $x<15$, control on and
off, across $180$ mid-circuit measurements --- $x$ preserved, $y=kx\bmod15$, all scratch back to
$\ket0$. Also checked on a lookup addressed by a \emph{non-uniform random} superposition, where any
residual phase would show: the address register is restored exactly, with leaked norm $0$.

\begin{center}\small
\begin{tabular}{@{}rrrrr@{}}
\toprule
$w$ & $L=2^w$ & recompute & measured & saving\\
\midrule
$2$  & $4$    & $6$     & $3$   & $2.0\times$\\
$4$  & $16$   & $30$    & $9$   & $3.3\times$\\
$8$  & $256$  & $510$   & $45$  & $11.3\times$\\
$11$ & $2048$ & $4\,094$ & $157$ & $26.1\times$\\
\bottomrule
\end{tabular}
\end{center}

\noindent
$2(L{-}1)$ against $2^{l}{-}1+2(2^{h}{-}1)$ with $l=h=w/2$. Uncomputation stops being half the
lookup cost and becomes a rounding error --- at the $w=\lg n\approx11$ of a cryptographic run,
$26\times$ cheaper, which turns the compute/uncompute pair from $2\times$ a lookup into
$1.04\times$.
\end{keybox}

\section{Emitting temporary ANDs}\label{sec:tempand}

Every Toffoli in the unary-iteration walk of \S\ref{im:sec:ui} is one half of a compute/uncompute
pair --- the tree computes $\mathrm{ctrl}\wedge a$, uses it, and clears it. Gidney's observation
\cite{gidney-and} is that such a pair need not cost two Toffolis. Computing an AND \emph{into a
qubit known to be $\ket0$} costs $4$ $T$ rather than $7$, and \emph{uncomputing} it costs none at
all: measure the ancilla in the $X$ basis and repair the phase with a Clifford $CZ$.

\begin{lstlisting}
def and_compute(qc, a, b, t):
    # t must start in |0>; leaves t = a AND b using 4 T gates, no phase
    qc.h(t); qc.t(t)
    qc.cx(b, t); qc.tdg(t)
    qc.cx(a, t); qc.t(t)
    qc.cx(b, t); qc.tdg(t)
    qc.h(t); qc.sdg(t)

def and_uncompute(qc, a, b, t, cbit):
    # erase t = a AND b for free: X-basis measurement + Clifford repair
    qc.h(t)
    qc.measure(t, cbit)
    with qc.if_test((cbit, 1)):
        qc.cz(a, b)
        qc.x(t)                                   # return the ancilla to |0>
\end{lstlisting}

\noindent
The final gate is $S^\dagger$, not $S$; three plausible variants of this circuit differ only in
that phase and are all wrong, which a truth table alone will not catch --- the check has to include
the phase. One classical bit suffices for the whole lookup, since each outcome is consumed by the
$CZ$ immediately after it is produced and may then be overwritten.

Substituting the pair into the walk gives \code{lookup\_ui\_ta}, verified to compute the same map
on every address for $m=1,2,3$ with both control values and all ancillas clean. Transpiled to
Clifford$+T$:

\begin{center}\small
\begin{tabular}{@{}rrrrr@{}}
\toprule
$m$ & $L$ & $T$ (plain) & $T$ (temporary AND) & saving\\
\midrule
$2$ & $4$  & $42$  & $12$  & $3.50\times$\\
$3$ & $8$  & $98$  & $28$  & $3.50\times$\\
$4$ & $16$ & $210$ & $60$  & $3.50\times$\\
$5$ & $32$ & $434$ & $124$ & $3.50\times$\\
\bottomrule
\end{tabular}
\end{center}

\noindent
Exactly $7/2$, and for the reason you would predict: half the Toffolis were uncomputes and are now
free, and the other half fell from $7$ $T$ to $4$. Since the $T$ count is what magic-state
distillation bills, this is the cheapest real saving in the note --- a few lines of code, no change
to any interface, and it applies to every lookup in Part~V.

\section{One counting qubit: $4n+2 \to 2n+3$}\label{sec:onectrl}

The counting register is the largest in the circuit and it is barely used: each of its qubits is
Hadamarded, controls one rung, and is measured. Worse, in the inverse QFT every controlled rotation
is controlled by a qubit that is \emph{about to be measured anyway}. So measure it first and make
the control classical --- the semiclassical Fourier transform of Griffiths and Niu
\cite{griffiths-niu}, applied to factoring by Mosca--Ekert \cite{mosca-ekert}, Parker--Plenio
\cite{parker-plenio} and Beauregard \cite{beauregard}. The counting register then never exists all
at once: one qubit is Hadamarded, drives its rung, receives phase corrections conditioned on every
previously measured bit, is measured, and is reset. The loop knows nothing about modular arithmetic
--- it takes the rungs as callbacks --- so it lives in \code{semiclassical.py} and the
elliptic-curve circuit of Part VII reuses it unchanged:

\begin{lstlisting}
def semiclassical_iqft(qc, ctrl, out, rungs):
    t = len(rungs)
    assert len(out) >= t, "one classical bit per rung"
    for i in range(t):
        qc.h(ctrl)
        rungs[t - 1 - i](qc, ctrl)                   # heaviest rung first
        for j in range(i):                           # the inverse QFT, classically
            with qc.if_test((out[j], 1)):
                qc.p(-math.pi / 2 ** (i - j), ctrl)
        qc.h(ctrl)
        qc.measure(ctrl, out[i])
        with qc.if_test((out[i], 1)):                # reset for the next rung
            qc.x(ctrl)


def order_circuit_1c(A, N, t=None):
    n = math.ceil(math.log2(N))
    t = t if t is not None else 2 * n
    ctr = QuantumRegister(1, "ctr")                  # the whole counting register
    tgt = QuantumRegister(n, "tgt")
    anc = QuantumRegister(n, "anc")
    sf = QuantumRegister(2, "sf")                    # sign, flag
    out = ClassicalRegister(t, "out")
    qc = QuantumCircuit(ctr, tgt, anc, sf, out)
    qc.x(tgt[0])                                     # target = |1>

    def rung(k):                                     # controlled U^(2^k)
        return lambda qc, c: c_ua(qc, c, pow(A, 2 ** k, N),
                                  list(tgt), list(anc), sf[0], sf[1], N)
    semiclassical_iqft(qc, ctr[0], out, [rung(k) for k in range(t)])
    return qc
\end{lstlisting}

\noindent
$1+n+n+2=2n+3$, the count in the title of Beauregard's paper. Measured:

\begin{center}\small
\begin{tabular}{@{}lrrrl@{}}
\toprule
$N$ & qubits & plain & time & result\\
\midrule
$15$ ($a=7$, $a=2$) & $\mathbf{11}$ & $18$ & $6$--$7$\,s & $r=4$, \ $15=3\times5$\\
$21$ ($a=2$)        & $\mathbf{13}$ & $22$ & $48$\,s & $r=6$, \ $21=7\times3$\\
\bottomrule
\end{tabular}
\end{center}

\begin{warnbox}[title={Three things that bite}]
\textbf{The bit order is not what you expect.} With the rungs run most-significant-exponent first,
the \emph{first} measured bit is the \emph{least} significant bit of the phase --- which happens to
be Qiskit's own bit order, so \code{order\_from\_counts} reads this circuit unchanged. Assuming the
reverse (as the author first did) silently yields $r=0$ at $N=15$; only testing both orderings
against a known $r$ settles it.

\textbf{Fewer qubits does not mean faster simulation.} Mid-circuit measurement forces Aer to
simulate shot by shot rather than evolve one statevector, so every shot is paid for individually.
On hardware this is the cheapest circuit in this note; in simulation part of the qubit saving is
given back. It is still a large net win --- $N=21$ in $48$\,s against roughly $1\,000$\,s for the
$22$-qubit version.

\textbf{It is incompatible with exponent windowing.} Windowing $w_e$ exponent bits
(\S\ref{im:sec:nested}) needs those $w_e$ qubits to coexist as lookup address bits; this circuit
keeps exactly one alive. The two are ends of one spectrum --- counting qubits $=w_e$ --- and you
cannot have both. Windowing \emph{costs} qubits ($24$ at $N=15$) where this \emph{saves} them
($11$).
\end{warnbox}

\section{Coset representation: deleting the modular adder}\label{sec:coset}

Everything so far pays five constant additions and a flag qubit for each modular addition
(\S\ref{im:sec:l2}), and measurement says that machinery costs $5.2$--$5.5\times$ the addition it
wraps --- flat in $n$, from $N=15$ to $N=1023$. It exists for one reason: to turn addition
mod $2^m$, which the register does natively, into addition mod $N$.

Zalka's coset representation \cite{zalka-coset}, as used by Gidney and Eker\aa\ \cite{gidney-ekera},
removes the need. Instead of storing $k$ as $\ket k$, store the periodic superposition
\[
\ket{k}_{\mathrm{coset}} \;=\; 2^{-c_{\mathrm{pad}}/2}\sum_{j=0}^{2^{c_{\mathrm{pad}}}-1}\ket{jN+k},
\]
in a register widened by $c_{\mathrm{pad}}$ padding qubits above the high bits. Every branch is
congruent to $k$ modulo $N$, so a \emph{plain, non-modular} addition of $a$ already produces a state
whose every branch is congruent to $k+a$. When $k+a\ge N$ the sum simply re-indexes,
\[
\sum_j\ket{jN+(k{+}a)}\;=\;\sum_{j'=1}^{2^{c_{\mathrm{pad}}}}\ket{j'N+(k{+}a{-}N)},
\]
which differs from the ideal coset state of $(k{+}a)\bmod N$ in only the two edge terms out of
$2^{c_{\mathrm{pad}}}$. \textbf{That is the entire approximation.} Level~2 collapses into Level~1,
and the sign and flag qubits go with it:

\begin{lstlisting}
def add(qc, X, reg, N=None):
    """Modular addition -- as a PLAIN addition.  This is the entire point:
    no compare, no subtract, no restore, no flag qubit, no sign qubit."""
    if N is not None:
        X %= N
    add_const(qc, X, list(reg))
\end{lstlisting}

\subsection{The deviation law, measured}

Gidney and Eker\aa\ bound the deviation per addition by $n/(c_{\mathrm{sep}}2^{c_{\mathrm{pad}}})$,
subadditive under composition, so a sequence of $A$ additions deviates by at most
$A\,2^{-c_{\mathrm{pad}}}$ when there are no carry runways. Measured, with $A=8$ random additions:

\begin{center}\small
\begin{tabular}{@{}rrrrr@{}}
\toprule
$N$ & $2^n-N$ & $c_{\mathrm{pad}}$ & measured error & bound $A/2^{c_{\mathrm{pad}}}$\\
\midrule
$15$ & $1$ & $2$ & $0.749$ & $2.000$\\
     &     & $3$ & $0.373$ & $1.000$\\
     &     & $4$ & $0.129$ & $0.500$\\
     &     & $5$ & $0.031$ & $0.250$\\
     &     & $6$ & $0.000$ & $0.125$\\
\midrule
$31$ & $1$ & $2$ & $0.507$ & $2.000$\\
     &     & $4$ & $0.061$ & $0.500$\\
     &     & $6$ & $0.000$ & $0.125$\\
\bottomrule
\end{tabular}
\end{center}

\noindent
The error halves with each padding qubit, exactly as advertised, and every measurement sits below
the bound --- which is loose by $2$--$4\times$, as a worst-case bound should be.

\begin{keybox}[title={Where the approximation actually bites}]
A first attempt to measure this at $N=21$ returned \emph{zero} error at every padding size, which
looks like a broken test and is not. The register holds $m=n+c_{\mathrm{pad}}$ bits, and the
headroom above the top branch is
$H=2^{c_{\mathrm{pad}}}(2^n-N)+N-1-k$. For $N=21$, $2^n-N=11$, so the headroom is enormous and no
branch ever wraps. The approximation only bites when the accumulated offset exceeds $H$, i.e.\ when
$N$ sits close to $2^n$ --- which is why the table above uses $N=15$ and $N=31$, where
$2^n-N=1$. \textbf{Choosing a friendly modulus hides the very effect you are trying to measure.}
\end{keybox}

\subsection{What it saves}

One addition, coset versus the seven-block modular adder:

\begin{center}\small
\begin{tabular}{@{}lrrrr@{}}
\toprule
adder & $N=15$ & $N=63$ & $N=255$ & $N=1023$\\
\midrule
Fourier (CNOT ratio)      & $1.7\times$ & $2.2\times$ & $2.5\times$ & $2.8\times$\\
ripple-carry (Toffoli ratio) & $3.0\times$ & $3.1\times$ & $3.5\times$ & $3.9\times$\\
\bottomrule
\end{tabular}
\end{center}

\noindent
Less than the naive $5.3\times$, because the coset register is $c_{\mathrm{pad}}$ qubits wider and
the adder must pay for that width; and better with ripple-carry, whose cost is linear in width,
than with Fourier, whose QFT sandwich is quadratic. The ripple-carry figure of $4n$ Toffoli for the
plain addition matches \cite{gidney-ekera} exactly. Both ratios improve with $n$, since
$c_{\mathrm{pad}}$ grows only logarithmically: $c_{\mathrm{pad}}=2\lg n+\lg n_e+\delta$ with
$\delta\in[2,10]$.

\subsection{Multiplying in the coset representation}

There is a subtlety that looks fatal and is not. In branch $j$ the register holds $v=jN+x$, so its
\emph{bits differ between branches} --- and shift-and-add multiplication uses those bits as
controls. It works anyway:
\[
\sum_i v_i\,(A2^i\bmod N)\;\equiv\;Av\;=\;A(jN+x)\;\equiv\;Ax \pmod N ,
\]
so every branch accumulates something congruent to $Ax$, whatever its own $j$. Verified at $N=7$
with both registers coset-encoded:

\begin{center}\small
\begin{tabular}{@{}rrrrr@{}}
\toprule
$c_{\mathrm{pad}}$ & $m$ & qubits & worst error & bound $A/2^{c_{\mathrm{pad}}}$\\
\midrule
$3$ & $6$ & $13$ & $0.0760$ & $0.7500$\\
$4$ & $7$ & $15$ & $0.0135$ & $0.4375$\\
\bottomrule
\end{tabular}
\end{center}

\begin{warnbox}[title={What is built, and what is not}]
Built and verified: the coset encoding, the plain-adder modular addition, the deviation law, and
the coset multiplier. \textbf{Not built: the full coset order-finding circuit.} Level~4's
multiply/swap/uncompute in coset form is the remaining piece --- the accumulator must start in
$\ket0_{\mathrm{coset}}=\sum_j\ket{jN}$ rather than $\ket0$, and the uncomputation must return it
there to within the deviation budget.

There is also a change of character worth stating plainly. Every other construction in this note is
\emph{exact}, and its tests assert exact equality with ancillas provably clean. The coset
representation is approximate by design, so its tests measure an error \emph{rate} against a
padding budget instead. That is not a weaker standard, but it is a different one, and it propagates:
once the arithmetic is approximate, every claim above it becomes probabilistic.
\end{warnbox}

\section{Where to go from here}\label{im:sec:beyond}

Two more upgrades, both easy and worth doing immediately.

\textbf{Approximate QFT (easy, big win).} Already built: pass
$d=\lceil\log_2 n\rceil+2$ as \code{approx} to \code{qft} (\S\ref{im:sec:l0}) and the count drops
from $O(n^2)$ to $O(n\log n)$ with negligible error \cite{barenco-aqft} --- measured, $144\to99$
gates at $n=16$, and the gap widens with $n$. Make it the default once your tests pass exactly.

\textbf{One counting qubit} --- built and measured in \S\ref{sec:onectrl}.

\textbf{Ripple-carry arithmetic} --- built and measured in \S\ref{im:sec:rc}. Do this one if you are
heading anywhere near fault tolerance.

A fuller Qiskit implementation of the plain (unwindowed) algorithm, with both adders and the one-control circuit, is \cite{repo}.\par\smallskip\textbf{Windowing} --- built and measured in \S\ref{im:sec:window}. Combined with the control savings
it is what takes published Toffoli counts from $\approx40n^3$ to $\approx0.3n^3$.

\begin{warnbox}[title={What this does not demonstrate}]
Simulation cost grows exponentially in qubit count, so noise-free demos stop at a few bits. On real
2020s hardware the binding constraint is not qubit count but noise: even the $N=15$ circuit's
two-qubit depth far exceeds current coherence. What you have built is the algorithm, correct and
testable --- not a factoring machine. The gap between them is fault tolerance, and at $n=2048$ it
is about $10^9$ Toffolis.
\end{warnbox}

\part*{Part VI --- Reading someone else's implementation}
\addcontentsline{toc}{part}{Part VI --- Reading someone else's implementation}

\section{The package at a glance}\label{sec:pkg}

Everything in Part~II exists as running code: the \code{qiskit-shor} package \cite{repo}
(see also the accompanying walkthrough \cite{blog}) implements the full tower on top of the Qiskit
SDK and factors integers end to end on a simulator. It is worth reading against the note for two
reasons. First, it realises the tower \emph{twice}, over two interchangeable Level-1 backends ---
the CDKM ripple-carry adder of Part~II \cite{cuccaro}, and a Fourier-space adder
\cite{draper,beauregard} that Part~II did not need but that dominates practical small-scale
demonstrations. Second, everything above Level~1 is written \emph{once}, generically over the
adder: the code is organised exactly along the levels of \S\ref{sec:zoom}, and the class hierarchy
is the proof of the note's claim that Levels 2--5 use nothing about addition except its interface.

\begin{center}
\begin{tabular}{@{}lll@{}}
\toprule
Module & Contents & Level(s)\\
\midrule
\code{gate\_utils.py} & \code{MAJGate}, \code{UMAGate}, controlled variants, QFT synthesis & 0\\
\code{rc\_adder.py}   & \code{RCAdderCircuit}: ripple-carry additions \cite{cuccaro} & 1\\
\code{qft\_adder.py}  & \code{QFTAdderCircuit}: Fourier-space additions \cite{draper,beauregard} & 1\\
\code{adder.py}       & \code{AdderCircuit} (abstract): modular add / multiply / exponentiate & 2--5\\
\code{shor.py}        & order-finding circuits, \code{find\_order}, \code{find\_factor} & 5 + classical\\
\bottomrule
\end{tabular}
\end{center}

\medskip\noindent
The class \code{AdderCircuit} is an abstract subclass of Qiskit's \code{QuantumCircuit}. It
declares two primitives, \code{add\_classical} and its controlled version --- ``add a known
constant to a register, optionally under a control''. The two backends supply those primitives;
the abstract class builds Levels 2--5 out of them, verbatim from
Theorems \ref{thm:modadd}--\ref{thm:cu}.

\begin{keybox}[title={Dictionary: the note's registers, the code's names}]
The code follows Qiskit's little-endian convention: qubit $k$ of a register carries bit $x_k$ of
$x=\sum_kx_k2^k$, so qubit $0$ is the least significant. The cast of \S\ref{sec:l2} maps as:
\begin{center}
\begin{tabular}{@{}lll@{}}
\toprule
In the note & In the code & Width\\
\midrule
value register $\ket x_{n+1}$ & \code{y\_reg} $+$ \code{overflow\_bit} & $n+1$\\
the sign qubit (top bit) & \code{overflow\_bit} & $1$\\
the flag $t$ & \code{ancilla\_bit} & $1$\\
multiplicand register $\ket x$ & \code{x\_reg} & $n$\\
counting register & \code{control\_register} & $2n$ (or $1$, \S\ref{sec:code-onectrl})\\
ripple-carry scratch & \code{a\_reg} (RC backend only) & $n+2$\\
\bottomrule
\end{tabular}
\end{center}
\end{keybox}

\section{Level 1 in code, twice}\label{sec:code-l1}

\subsection{The ripple-carry backend}

\code{gate\_utils.py} defines $\MAJ$ and $\UMA$ exactly as in \S\ref{sec:l1}, as reusable Qiskit
gates (listings are lightly abridged --- docstrings and boilerplate trimmed --- throughout Part~IV):

\begin{lstlisting}
class MAJGate(SingletonGate):
    """|c,b,a> -> |c+a, b+a, MAJ(a,b,c)>   (addition mod 2)"""
    def _define(self):
        q = QuantumRegister(3, "q")       # q0=c, q1=b, q2=a
        qc = QuantumCircuit(q, name="maj")
        qc.cx(q[2], q[1])                 # CNOT(a -> b)
        qc.cx(q[2], q[0])                 # CNOT(a -> c)
        qc.ccx(q[0], q[1], q[2])          # TOF(c, b -> a)
        self.definition = qc

class UMAGate(SingletonGate):
    """UMA . MAJ |c,b,a> = |c, a+b+c, a>"""
    def _define(self):
        q = QuantumRegister(3, "q")
        qc = QuantumCircuit(q, name="uma")
        qc.ccx(q[0], q[1], q[2])          # TOF(c, b -> a)
        qc.cx(q[2], q[0])                 # CNOT(a -> c)
        qc.cx(q[0], q[1])                 # CNOT(c -> b)
        self.definition = qc
\end{lstlisting}

\noindent
and \code{RCAdderCircuit.add\_quantum} chains them precisely as in the diagram of
Theorem~\ref{thm:cdkm} --- $\MAJ$ up the bits, one CNOT into the overflow qubit at the single
instant the final carry sits in the clear, $\UMA$ back down:

\begin{lstlisting}
def add_quantum(self, x_reg, y_reg, ancilla_bit, overflow_bit=None):
    """|x>_n |y>_n |0> -> |x>_n |y+x>_{n+1} |0>"""
    n = len(y_reg)
    maj, uma = MAJGate(), UMAGate(RC_version=True)
    self.append(maj, [ancilla_bit, y_reg[0], x_reg[0]])
    for i in range(n - 1):
        self.append(maj, [x_reg[i], y_reg[i + 1], x_reg[i + 1]])
    if overflow_bit:
        self.cx(x_reg[n - 1], overflow_bit)
    for i in range(n - 1, 0, -1):
        self.append(uma, [x_reg[i - 1], y_reg[i], x_reg[i]])
    self.append(uma, [ancilla_bit, y_reg[0], x_reg[0]])
\end{lstlisting}

\noindent
Two departures from the bare Theorem~\ref{thm:cdkm} are worth pausing on.

\begin{remark}[the $3$-CNOT $\UMA$]
The flag \code{RC\_version=True} selects Cuccaro et al.'s alternative six-gate $\UMA$
($X\,b$;\ $\CNOT\,a\!\to\!b$;\ $\TOF\,a,b\!\to\!c$;\ $X\,b$;\ $\CNOT\,c\!\to\!a$;\
$\CNOT\,c\!\to\!b$) \cite{cuccaro}. It computes the same permutation as the three-gate version ---
the package's unit tests check them against each other --- but its final two CNOTs share no qubit
with the \emph{next} $\UMA$'s first gates, so consecutive blocks overlap in time and the adder's
depth drops by a constant factor. Same matrix, better schedule: a first taste of the fact that gate
count and depth are different currencies.
\end{remark}

\begin{remark}[adding a constant with an ancilla register]
CDKM adds two \emph{quantum} integers, but Level~2 needs $y\mathbin{+{=}}X$ for a \emph{classical}
$X$. The RC backend does the simplest correct thing: it writes $X$ into $n$ clean ancillas with
$X$ gates (\code{bitwise\_AND\_classical}), runs \code{add\_quantum}, and unwrites $X$ with the
same $X$ gates --- which is why this backend carries the extra \code{a\_reg} of $n+2$ qubits
through every signature, and why its total budget is $5n+4$ qubits against the Fourier backend's
$4n+2$ (\S\ref{sec:code-counts}). Subtraction, needed for blocks (2) and (5) of
Theorem~\ref{thm:modadd}, is the same routine with the two's-complement of $X$: the code reduces
$X\bmod 2^{n+1}$ first, so \code{add\_classical(X - N, ...)} just works.
\end{remark}

\noindent
The controlled adder is where the code quietly reuses the control lemma of \S\ref{sec:l4}. Rather
than controlling all six gates of every block, \code{gate\_utils.py} defines \code{MCXMAJGate} and
\code{MCXUMAGate}, in which \emph{only the CNOT targeting the $b$ wire} acquires the outer control:
\begin{align*}
\MAJ_c&:\ c\text{-}\CNOT(a\!\to\!b),\ \CNOT(a\!\to\!c),\ \TOF(c,b\!\to\!a);\\
\UMA_c&:\ \TOF(c,b\!\to\!a),\ \CNOT(a\!\to\!c),\ c\text{-}\CNOT(c\!\to\!b).
\end{align*}
With the control off, what remains of $\MAJ$ is $\big[\CNOT(a\!\to\!c),\ \TOF(c,b\!\to\!a)\big]$
and what remains of $\UMA$ is its exact mirror $\big[\TOF(c,b\!\to\!a),\ \CNOT(a\!\to\!c)\big]$;
since the $\UMA$ chain retraces the $\MAJ$ chain in reverse order, the whole residue telescopes to
the identity --- the \emph{conjugation} pattern of \S\ref{sec:l4}, found in the wild. One controlled
gate per block, $2n{+}1$ in total (the overflow CNOT is controlled too), instead of $6$ per block.

\subsection{The Fourier backend}\label{sec:code-qft}

The default backend does no carry propagation at all. Draper's observation \cite{draper} is that
the QFT turns addition into phase accumulation. Write the full QFT on $n$ qubits, in the code's
little-endian convention, as
\[
F\ket y=\bigotimes_{i=0}^{n-1}\frac{\ket0+e^{2\pi i\,y\,2^{\,i-n}}\ket1}{\sqrt2}.
\]
Each qubit carries the value $y$ in the \emph{period} of a phase, at a different scale --- in
Draper's reading, qubit $i$ holds the binary fraction $e^{2\pi i\,(0.y_{n-i-1}\cdots y_1y_0)}$ of
the $n{-}i$ low bits of $y$: the value is spread across scales, not across carries. Adding a
classical constant $X$ is then $n$ \emph{parallel, uncontrolled} phase gates --- qubit $i$ gets
$P(2\pi X\,2^{\,i-n})$, turning $e^{2\pi iy2^{i-n}}$ into $e^{2\pi i(y+X)2^{i-n}}$ --- and an
inverse QFT returns to the computational basis holding $\ket{(y+X)\bmod2^n}$. That is the entire
implementation:

\begin{lstlisting}
def add_classical(self, X, y_reg, a_reg=None, include_QFT=True):
    """|y> -> |y + X mod 2^n>   (no ancillas needed)"""
    n = len(y_reg)
    qft = QFTFullGate(n, approximation_degree=self.qft_approx_degree(n))
    if include_QFT:
        self.compose(qft, y_reg, inplace=True)
    for i in range(n):
        self.p(2 * np.pi * X * 2 ** (i - n), y_reg[i])  # addition in Fourier space
    if include_QFT:
        self.compose(qft.inverse(), y_reg, inplace=True)
\end{lstlisting}

\noindent
No carries, no ancillas, and the controlled version (\code{c\_add\_classical}) simply upgrades each
phase gate to a controlled phase gate --- the QFT layers keep no control, being the conjugating
$A\cdots A^\dagger$ of the control lemma, \S\ref{sec:l4}. This is why the Fourier route needs two
registers fewer than the ripple-carry route, and why Beauregard \cite{beauregard} could squeeze the
whole algorithm into $2n+3$ qubits.

\begin{remark}[what the Fourier route buys in depth]
Two structural facts, one from each source paper, sharpen the economics. \emph{Beauregard
\cite{beauregard}:} because $X$ is classical, all the rotations that land on one qubit multiply
into a \emph{single} phase gate --- the loop above emits exactly one $P$ per qubit, so the
addition layer between the QFTs is $n$ parallel one-qubit gates at \textbf{depth 1}.
\emph{Draper \cite{draper}:} even in the quantum--quantum adder (\code{add\_quantum}, $O(n^2)$
controlled rotations) all rotations \emph{commute} --- unlike the QFT itself, whose Hadamards
impose an ordering --- so they schedule into $n{+}1$ time slices, or $O(\log n)$ slices under the
small-angle cutoff below. The depth bottleneck of the Fourier route is therefore the QFT sandwich,
at depth $O(n)$, not the addition it encloses. Finally, the same gate run backwards subtracts:
$\phi\mathrm{ADD}(a)^\dagger$ sends $\phi(b)$ to $\phi(b-a)$ if $b\ge a$ and to
$\phi(2^{n+1}-(a-b))$ otherwise, so the top qubit is the comparator $[\,b<a\,]$ ---
\S\ref{sec:l1}'s ``comparison is free'', relocated to Fourier space, and the workhorse of the sign
tests in \S\ref{sec:code-l2}.
\end{remark}

\begin{keybox}[title={Doesn't this break ``no complex number below Level 5''?}]
Inside the sandwich, yes: the intermediate states have irrational phases, and a single $P(\theta)$
is certainly not a permutation matrix. But the \emph{composite}
$F^\dagger\cdot(\text{phases})\cdot F$ is exactly the same $0/1$ permutation matrix
$\ket y\mapsto\ket{y+X\bmod2^n}$ as the CDKM circuit. The slogan of \S\ref{sec:dual} survives in a
sharper form: the two backends are two factorisations of one permutation --- one through
$\mathbb F_2$ carry algebra, one through the Fourier dual --- and everything above Level~1 cannot
tell them apart. The trade is not conceptual but economic: the Fourier route pays in
\emph{precision} (doubly-controlled rotations by angles as small as $2\pi/2^{n+1}$) instead of in
\emph{qubits}, which is why the package also offers the RC backend, whose gate set
$\{X,\CNOT,\TOF\}$ is the one fault-tolerant hardware actually wants.
\end{keybox}

\begin{remark}[the approximate QFT]
Rotations by $2\pi/2^k$ for large $k$ are far below any hardware's resolution, and dropping them
barely disturbs the result: keeping only angles at least $\pi/2^{d}$ with
$d=\lceil\log_2n\rceil+2$ costs an error that stays negligible while cutting the QFT from $O(n^2)$
to $O(n\log n)$ gates \cite{barenco-aqft,cheung}. The package exposes exactly this cutoff via the
flag \code{approx\_QFT}; it is the ``approximate QFT'' priced in Part~III's bill.
\end{remark}

\section{Levels 2--4: written once, above the interface}\label{sec:code-l2}

\subsection{The modular adder}

\code{AdderCircuit.add\_classical\_modulo} is Theorem~\ref{thm:modadd} with one refinement. Since
$a$ \emph{and} $N$ are both classical, blocks (1) and (2) --- add $a$, subtract $N$ --- merge into
a single addition of the constant $a-N$:

\begin{lstlisting}
def add_classical_modulo(self, X, y_reg, ancilla_bit, overflow_bit, N,
                         a_reg=None, reset_ancilla=True):
    """|y>|0> -> |(X + y) mod N>|0>,  assuming 0 <= X < N and 0 <= y < N."""
    y_bits = y_reg[:] + [overflow_bit]        # the (n+1)-wide value register
    n = math.ceil(math.log2(N))

    self.add_classical(X - N, y_bits, a_reg)  # blocks (1)+(2): y + X - N
    self.cx(y_bits[n], ancilla_bit)           # block (3): sign -> t
    self.c_add_classical(ancilla_bit, N, y_bits, a_reg)   # block (4)

    if reset_ancilla:                         # blocks (5)-(7): clear the flag
        self.add_classical(-X, y_bits, a_reg)
        self.cx(y_bits[n], ancilla_bit)
        self.x(ancilla_bit)
        self.add_classical(X, y_bits, a_reg)
\end{lstlisting}

\noindent
Read it against the table in \S\ref{sec:l2} line by line: \code{y\_bits[n]} is the sign qubit,
\code{ancilla\_bit} is $t$, and the last four lines are blocks (5), (6), (7) --- recompute the flag
from the answer, XOR it against itself, restore. (The code performs the CNOT before the $X$ where
the note performs the $X$ before the CNOT; both compute $t\mapsto\neg(t\oplus s)$.) The merge
means the uncontrolled modular adder costs \emph{four} constant additions rather than five ---
compile-time constants at work, as promised in the remark of \S\ref{sec:l1}.

The controlled version cannot merge blocks (1) and (2), because block (1) must switch with the
control while block (2) must always fire. It is also where the code makes a different --- and
instructive --- choice of which blocks carry the control:

\begin{lstlisting}
def c_add_classical_modulo(self, control_reg, X, y_reg, ancilla_bit,
                           overflow_bit, N, a_reg=None, reset_ancilla=True):
    y_bits = y_reg[:] + [overflow_bit]
    n = math.ceil(math.log2(N))

    self.c_add_classical(control_reg, X, y_bits, a_reg)  # (1), controlled
    self.add_classical(-N, y_bits, a_reg)                # (2), uncontrolled
    self.cx(y_bits[n], ancilla_bit)                      # (3)
    self.c_add_classical(ancilla_bit, N, y_bits, a_reg)  # (4), controlled on t

    if reset_ancilla:
        self.add_classical(-X, y_bits, a_reg)            # (5), uncontrolled!
        self.mcx(control_reg + [y_bits[n]], ancilla_bit) # (6), CNOT -> Toffoli
        self.x(ancilla_bit)
        self.add_classical(X, y_bits, a_reg)             # (7), uncontrolled!
\end{lstlisting}

\begin{keybox}[title={Two of seven, not three of seven}]
Where \S\ref{sec:l4} controlled blocks (1), (5), (7), the code controls block (1) and, instead of
controlling the additions of blocks (5) and (7), upgrades block (6)'s flag-clearing CNOT to a
Toffoli. Both choices are correct, and both are licensed by the same lemma; only the residue
argument changes. Control off, code version: block (2) drives the register to $y-N<0$, sign set,
$t$ flips on, block (4) adds $N$ back --- the familiar conspiracy. Then block (5) subtracts $X$
\emph{unconditionally}, the Toffoli does nothing (its outer control is off), the $X$ gate undoes
the earlier flip of $t$, and block (7) adds $X$ back: the residue $(-X)$ then $(+X)$ cancels as
pure conjugation. Control on: after block (5) the sign qubit holds $\neg t$
(Theorem~\ref{thm:modadd}'s proof), so the Toffoli sets $t\mapsto t\oplus\neg t=1$ and the $X$
clears it. The bargain: two additions run uncontrolled --- saving $O(n)$ controlled gates each ---
at the cost of upgrading one CNOT to a Toffoli, which is $O(1)$.
\end{keybox}

\begin{remark}[the hidden QFT sandwiches]
Beauregard's drawn circuit (his Fig.~5 \cite{beauregard}) keeps the value register in Fourier
space for the whole modular adder and drops out only twice --- a $\mathrm{QFT}^{-1}\cdots
\mathrm{QFT}$ pair around each read of the sign qubit --- because a two's-complement sign is a
computational-basis fact, invisible to a register that stores its value in phases. The code
contains the same structure without ever writing it: every \code{add\_classical} of the Fourier
backend closes with an inverse QFT and the next opens with a QFT, so wherever two additions abut,
the pair is redundant (Beauregard draws nothing there, and a transpiler cancels it), and wherever
a CNOT to the flag intervenes --- blocks (3) and (6) --- the pair survives as exactly his
sandwich. The RC backend has no such subtlety: its register never leaves the computational basis,
which is one more reason its gate count is a factor of $n$ better and its qubit count $n{+}2$
worse.
\end{remark}

\subsection{The multiplier, out of place and then in place}

Level~3 is a four-line loop --- Theorem~\ref{thm:cmult} transcribed. Note the constant
\code{((A \% N) * 2**i) \% N} is the shifted multiple $2^iA\bmod N$, computed in Python at circuit
construction time --- ``compile time'' in the note's vocabulary:

\begin{lstlisting}
def c_add_quantum_modulo(self, control_reg, x_reg, y_reg, ancilla_bit,
                         overflow_bit, N, a_reg=None, A=1):
    """|c>|x>|y>|0> -> |c>|x>|(y + c*A*x) mod N>|0>"""
    for i in range(len(x_reg)):
        self.c_add_classical_modulo(
            control_reg=control_reg + [x_reg[i]],   # doubly controlled
            X=((A % N) * 2**i) % N,                 # 2^i A mod N, classical
            y_reg=y_reg, ancilla_bit=ancilla_bit,
            overflow_bit=overflow_bit, N=N, a_reg=a_reg)
\end{lstlisting}

\noindent
Level~4 --- Theorem~\ref{thm:cu} --- is the multiply/swap/uncompute sandwich, with the modular
inverse handed to Python's built-in \code{pow}:

\begin{lstlisting}
def c_multiply_modulo(self, control_reg, A, x_reg, y_reg, overflow_bit,
                      ancilla_bit, N, a_reg=None, ...):
    """|c>|x>|0> -> |c>|A^c x mod N>|0>"""
    # CMULT(A):  |x>|0> -> |x>|Ax mod N>
    self.c_add_quantum_modulo(control_reg, x_reg, y_reg,
                              ancilla_bit, overflow_bit, N, a_reg, A=A)
    # controlled SWAP: n Fredkin gates sharing the control
    for i in range(n):
        self.append(SwapGate().control(k), control_reg + [x_reg[i], y_reg[i]])
    # CMULT(A^{-1})^dagger: uncompute x from the product
    B = pow(A, -1, N)                       # extended Euclid, classical
    self.c_add_quantum_modulo(control_reg, y_reg, x_reg,
                              ancilla_bit, overflow_bit, N, a_reg, A=-B)
\end{lstlisting}

\begin{remark}[the dagger, implemented as a minus sign]
Theorem~\ref{thm:cu} calls for $\CMULT(\bar a)^\dagger$ --- run the $\bar a$-multiplier backwards.
The code never reverses a circuit: it passes $A=-B$ and lets the constant reduction fold
$-\bar a\bmod N=N-\bar a$ into an ordinary forward multiplier. The two are the same operator,
because the inverse of ``add $\bar ax$ to this register mod $N$'' is exactly ``subtract
$\bar ax$ from this register mod $N$''. Adjoints of permutations are permutations, and for adders
they are again adders --- one more classical fact doing quantum work.
\end{remark}

\section{Level 5 and the two order-finding circuits}\label{sec:code-l5}

The ladder is the last, shortest loop --- rung constants by repeated squaring, one controlled
in-place multiplier per counting qubit:

\begin{lstlisting}
def exponentiate_modulo(self, A, x_reg, y_reg, ancilla_reg, N, a_reg=None):
    """|x>_m |y>_n |0>_{n+2} -> |x>_m |(A^x y) mod N>_n |0>_{n+2}"""
    n = math.ceil(math.log2(N))
    for i in range(len(x_reg)):
        self.c_multiply_modulo(
            control_reg=x_reg[i],
            A=pow(A, 2**i, N),              # the rung constant a^{2^i} mod N
            x_reg=y_reg, y_reg=ancilla_reg[:n],
            overflow_bit=ancilla_reg[n], ancilla_bit=ancilla_reg[n + 1],
            N=N, a_reg=a_reg)
\end{lstlisting}

\noindent
and the full order-finding circuit of \S\ref{sec:l5} --- Hadamards, ladder, inverse QFT,
measurement --- assembles in a dozen lines (\code{shor.py}):

\begin{lstlisting}
def order_finding_circuit(A, N, precision=None, adder="qft_adder"):
    n = math.ceil(math.log2(N))
    m = precision if precision is not None else 2 * n   # t = 2n counting qubits

    control_register = QuantumRegister(m)
    target_register  = QuantumRegister(n)
    ancilla_register = QuantumRegister(n + 2)
    output_register  = ClassicalRegister(m, name="output_bits")
    qc = QFTAdderCircuit(control_register, target_register,
                         ancilla_register, output_register)

    for i in range(m):
        qc.h(control_register[i])           # Hadamard layer
    qc.x(target_register[0])                # target = |1>
    qc.exponentiate_modulo(A=A, x_reg=control_register,
                           y_reg=target_register,
                           ancilla_reg=ancilla_register, N=N)
    qc.compose(QFTGate(m).inverse(), qubits=control_register, inplace=True)
    qc.measure(control_register, output_register)
    return qc
\end{lstlisting}

\noindent
Count the qubits: $2n$ counting $+\ n$ target $+\ (n{+}2)$ scratch ($n$ for the out-of-place
product, the sign qubit, the flag $t$) $=4n+2$ --- with the RC backend, $n{+}2$ more for
\code{a\_reg}, hence $5n+4$.

\subsection{One control qubit: the $2n+3$ circuit}\label{sec:code-onectrl}

Beauregard's headline saving \cite{beauregard} is that the $2n$ counting qubits can be replaced by
\emph{one}, recycled $2n$ times. The idea, due to Mosca--Ekert \cite{mosca-ekert} and
Parker--Plenio \cite{parker-plenio}, rests on the semiclassical Fourier transform of
Griffiths--Niu \cite{griffiths-niu}: in the inverse QFT, every two-qubit controlled rotation is
controlled by a qubit that is \emph{about to be measured anyway} --- so measure it first, and
replace the quantum control by a classical one conditioned on the outcome. The counting register
then never needs to exist all at once. One qubit is Hadamarded, drives the controlled
multiplication for its rung, receives phase corrections determined by all \emph{previously
measured} bits, is measured, and is reset for the next rung:

\begin{lstlisting}
    c_bit = control_register[0]             # the one and only counting qubit
    for i in range(m):
        qc.h(c_bit)
        qc.c_multiply_modulo(               # rung i, MSB first
            control_reg=c_bit, A=pow(A, 2 ** (m - i - 1), N),
            x_reg=target_register, y_reg=ancilla_register[:n],
            overflow_bit=ancilla_register[n],
            ancilla_bit=ancilla_register[n + 1], N=N)
        for j in range(i):                  # semiclassical inverse QFT:
            with qc.if_test((output_register[j], 1)):     # classically
                qc.p(-math.pi / 2 ** (i - j), c_bit)      # conditioned phases
        qc.h(c_bit)
        qc.measure(c_bit, output_register[i])
        with qc.if_test((output_register[i], 1)):
            qc.x(c_bit)                     # reset for the next rung
\end{lstlisting}

\noindent
The quantum inverse QFT has disappeared entirely, traded for mid-circuit measurement and
feed-forward (Qiskit's \code{if\_test} control flow --- on hardware, ``dynamic circuits''). Total:
$1+n+(n{+}2)=2n+3$ qubits, the count in the title of Beauregard's paper. The two variants are
toggled by one argument, and they sample from the same distribution: nothing about the
\emph{arithmetic} changed, only about where the counting register lives in time.

\subsection{The bill, as the package pays it}\label{sec:code-counts}

With $n=\lceil\log_2N\rceil$ and the package's own accounting \cite{repo} (gates counted in the
Qiskit basis before transpilation --- see the note in the README on why such counts are only
comparative):

\begin{center}
\begin{tabular}{@{}lcccc@{}}
\toprule
Backend & qubits & qubits (one-control) & gates & depth\\
\midrule
Fourier \cite{draper,beauregard} & $4n+2$ & $2n+3$ & $O(n^4)$ & $O(n^3)$\\
\quad with approximate QFT \cite{barenco-aqft,cheung} & $4n+2$ & $2n+3$ & $O(n^3\log n)$ & $O(n^2\log n)$\\
Ripple-carry \cite{cuccaro} & $5n+4$ & $3n+5$ & $O(n^3)$ & $O(n^3)$\\
\bottomrule
\end{tabular}
\end{center}

\medskip\noindent
The Fourier backend's extra factor of $n$ is the QFT sandwich around every addition; the
ripple-carry backend's extra $n{+}2$ qubits are the scratch register for constant encoding. The
package also implements, as a standalone method (\code{c\_add\_quantum\_optimized\_depth}), the
depth-$O(n)$ controlled Fourier addition of Pavlidis and Gizopoulos \cite{pavlidis}, which moves
the control off the $O(n^2)$ rotation block --- the rotations run uncontrolled, conjugated by
controlled bit flips of the operand --- but this trick does not compose with the modular reduction
of \S\ref{sec:code-l2}, so the order-finding circuit does not use it; their full architecture
reworks the modular level too, and is beyond both the package and this note.

Beauregard's own ledger \cite{beauregard} prices each level of the Fourier route, with
$k_{\max}$ the approximate-QFT cutoff: $\phi\mathrm{ADD}(a)$ is $n{+}1$ qubits at depth $O(1)$;
the doubly controlled modular adder $n{+}4$ qubits, $O(n\,k_{\max})$ gates, depth $O(n)$;
$\CMULT$ and $c\text{-}U_a$ are $2n{+}3$ qubits at depth $O(n^2)$; the whole circuit $2n{+}3$
qubits, $O(n^3\lg n)$ gates, depth $O(n^3)$. His $2n+3$ was, at the time, the last word in a
decade-long qubit race:

\begin{center}\small
\begin{tabular}{@{}llll@{}}
\toprule
Year & Construction & Qubits & Notes\\
\midrule
1996 & Vedral--Barenco--Ekert \cite{vbe} & $7n+1$ & $5n+2$ with basic optimisation\\
1996 & Beckman--Chari--Devabhaktuni--Preskill \cite{bcdp} & $5n+1$ & $4n+1$ with unbounded Toffolis\\
1998 & Zalka \cite{zalka} & $3n+O(\lg n)$ & elementary gates only\\
2003 & Beauregard \cite{beauregard} & $2n+3$ & the one-control circuit of \S\ref{sec:code-onectrl}\\
\bottomrule
\end{tabular}
\end{center}

\medskip\noindent
At $N=15$ the $2n+3$ circuit is eleven qubits. The 2001 NMR experiment that factored $15$
\cite{vandersypen} used seven --- by compiling for that one modulus, whose rung constants
$7,4,1,1,\dots$ (\S\ref{sec:l5}) make most of the ladder the identity. General-purpose and
special-cased circuits are different objects, and only the former is priced by the table above.

\section{Cross-check: factoring 15 from scratch}\label{sec:code-15}

The same two papers admit more than one faithful implementation, and it is instructive to read a
second. \emph{Building Shor's algorithm from scratch: factoring 15} \cite{notebook15} is an
independent Qiskit notebook that assembles the identical Draper--Beauregard stack --- its QPE
preamble follows Nielsen--Chuang \cite{nielsen-chuang} --- one class per level of the tower, each
verified on deterministic test vectors before the next is built ($5$ through the QFT and back;
$(5+3)\bmod 9=8$ through the modular adder; $(7+3\cdot3)\bmod 7=2$ through the multiplier). The
correspondence is exact:

\begin{center}\small
\begin{tabular}{@{}lll@{}}
\toprule
Object & \code{qiskit-shor} & notebook \cite{notebook15}\\
\midrule
QFT (Level 0/1) & \code{QFTFullGate} & \code{QFTBlock}\\
$\phi\mathrm{ADD}(a)$ (\S\ref{sec:code-qft}) & \code{add\_classical} & \code{QAdderClassicalBlock}\\
modular adder (Thm.~\ref{thm:modadd}) & \code{add\_classical\_modulo} & \code{add\_mod\_N\_classic}\\
$\CMULT$ (Thm.~\ref{thm:cmult}) & \code{c\_add\_quantum\_modulo} & \code{CMULTaMODN\_c}\\
$c\text{-}U_a$ (Thm.~\ref{thm:cu}) & \code{c\_multiply\_modulo} & \code{Controlled\_U\_a}\\
ladder $+$ QPE (\S\ref{sec:l5}) & \code{order\_finding\_circuit} & final assembly cell\\
\bottomrule
\end{tabular}
\end{center}

\medskip\noindent
Three divergences are worth reading closely, because each lands on a point the note has already
made.

\begin{remark}[three of seven, again]
The notebook builds its doubly controlled additions by calling Qiskit's \code{.control(2)} on the
$\phi\mathrm{ADD}$ gates of blocks (1), (5) and (7) --- Beauregard's original prescription, and
precisely the $3$-of-$7$ choice of \S\ref{sec:l4}. So the two codebases realise the \emph{two}
correct control placements side by side: blocks (1), (5), (7) in the notebook; block (1) plus the
Toffoli'd flag-clear in \code{qiskit-shor} (\S\ref{sec:code-l2}). Diff them and the residue
argument of the control lemma is the only invariant.
\end{remark}

\begin{warnbox}[title={Gift 1, forgone}]
The notebook's ladder does \emph{not} precompute rung constants. It applies the same
$c\text{-}U_a$ gate $2^i$ times for counting qubit $i$:
\begin{center}
\code{for j in range(2**i): qc.append(Ctrl\_U\_a\_gate, ...)}
\end{center}
--- $1+2+4+8=2^t-1=15$ multiplier applications where \code{qiskit-shor} uses $t=4$ rungs with
\code{pow(A, 2**i, N)}. This is the exponential sum of \S1.1, live in running code. At $N=15$ it
costs a factor of four and nobody notices; at any cryptographic size it is the difference between
an algorithm and a demonstration. Gift~1 is two lines of Python --- and it is still a gift you
must actually accept.
\end{warnbox}

\begin{remark}[an undersized counting register]
The notebook sets $t=n=4$ rather than $t=2n=8$, noting the shortcut. For this instance it is
sound: the order $r=4$ divides $2^t$, so the phase $s/r$ is \emph{exactly} representable in $4$
bits and no rounding occurs. For general $N$ the order does not divide $2^t$, the measured value
only approximates $s/r$, and the $2n$ bits are what guarantees the continued-fraction step
converges to the right denominator \cite{shor,cemm}. Another instance of $15$'s exceptional
friendliness (\S\ref{sec:l5}) --- and of why it proves so little.
\end{remark}

\begin{warnbox}[title={Read the histogram, not the factors}]
With $a=7$, $N=15$, $r=4$, $t=4$, theory owes four sharp peaks: the measurement should return
$y=2^t s/r\in\{0,4,8,12\}$, a quarter of the shots each. The run recorded in the notebook shows
instead a \emph{flat} distribution --- all sixteen outcomes near $512$ of $8\,192$ shots, and
measured phases like $3/16$ and $5/16$ that no $s/4$ can produce. That flatness is a diagnosis,
not statistics: it is the uniform-noise signature \S\ref{sec:garbage} derives for a counting
register left entangled with something --- interference killed by $\braket{g_{x'}}{g_x}=0$ ---
somewhere in the final assembly, even though every gate-level unit test passed deterministically.
The notebook's retry loop still ends by printing $3\times5$, because its classical tail ---
continued fractions, then \emph{verify} $\gcd(a^{r/2}\pm1,N)$ --- filters random candidates until
one passes, which uniform noise eventually supplies. Two lessons, both cheap. End-to-end
correctness lives in the \emph{distribution}, so test the distribution (the seeded
\code{find\_order} test of \S\ref{sec:code-run} is exactly such a test). And a pipeline whose
classical tail can succeed on white noise will not tell you when its quantum head has failed ---
the algorithm's own robustness is a bug-mask.
\end{warnbox}

\section{Running it}\label{sec:code-run}

What remains of Shor's algorithm \cite{shor} outside the circuit is classical, and \code{shor.py}
contains all of it. \code{find\_order} builds either circuit, transpiles it, samples it, and reads
the order out of the measurement distribution --- the continued-fraction step is Python's
\code{Fraction}:

\begin{lstlisting}
x = int(best_output, 2)                       # measured phase, t bits
r = Fraction(x / 2**precision).limit_denominator(N - 1).denominator
if pow(A, r, N) == 1:                          # verify classically
    return r
\end{lstlisting}

\noindent
--- exactly the ``continued fractions recover $r$'' of Part~III, with the candidate checked by one
classical exponentiation. \code{find\_factor} wraps it in the standard loop: strip factors of $2$
and prime powers, pick a random base, gcd-check it (the lucky case has already factored $N$),
find the order, and try $\gcd(a^{r/2}-1,N)$ when $r$ is even. End to end, on a noise-free
simulator:

\begin{lstlisting}
from qiskit_aer import AerSimulator
from qiskit_aer.primitives import SamplerV2 as AerSampler
from qiskit.transpiler.preset_passmanagers import generate_preset_pass_manager
from qiskit_shor.shor import find_factor, find_order

sim = AerSimulator()
pm = generate_preset_pass_manager(backend=sim, optimization_level=1)

order, dist = find_order(7, 15, AerSampler(), pm, num_shots=100,
                         one_control_circuit=True)   # -> 4
factor = find_factor(15, AerSampler(), pm, num_tries=3,
                     num_shots_per_trial=100)        # -> 3 or 5
\end{lstlisting}

\begin{keybox}[title={What the package machine-checks}]
The note's ``Verified'' boxes have counterparts in the repository's \code{pytest} suite: unit
tests drive every arithmetic method on all (or sampled) basis states through Qiskit's simulator ---
additions, modular additions, multipliers, for both backends, including that every ancilla returns
to $\ket0$ --- and the end-to-end tests recover the order $r=4$ of both $A=2$ and $A=7$ modulo
$15$, with both circuit variants, on a seeded \code{AerSimulator}. The same discipline at two
scales: trust nothing you have not enumerated.
\end{keybox}

\begin{warnbox}[title={What this does not demonstrate}]
Simulation cost grows exponentially in qubit count, so noise-free demonstrations stop around
$N$ of a few bits --- and on 2025-era hardware, noise, not qubit count, is the binding constraint:
a few hundred physical qubits exist, but the $\sim\!2\,000$-gate two-qubit-depth of even the
$N=15$ one-control circuit far exceeds current coherence \cite{repo}. The package is the tower of
Part~II made executable and testable, not a factoring machine; the gap between the two is the
fault-tolerance overhead priced in Part~III's bill.
\end{warnbox}















\clearpage
\part*{Part VII --- The other Shor: elliptic-curve discrete logarithms}
\addcontentsline{toc}{part}{Part VII --- The other Shor: elliptic-curve discrete logarithms}

\begin{keybox}[title={Why a second half}]
Everything so far factors integers. The same three layers --- Hadamards, an oracle ladder, an
inverse QFT --- also break elliptic-curve cryptography, and the outer two layers are unchanged. What
changes is the middle, and it changes completely.

For factoring, the group is $(\mathbb{Z}/N)^\times$ and the oracle multiplies by a
\emph{classical} constant. Multiplication by a constant is cheap: Part~II built it out of shifted
additions. For the discrete logarithm on a curve, the group is $E(\mathbb{F}_p)$ and the oracle adds
a classical \emph{point} --- and the chord-and-tangent law needs a division, i.e.\ a modular
\emph{inverse}, of a quantum value. That single fact is the whole story of this part: an inversion
costs an order of magnitude more than the multiplication it replaces, and every optimisation below
is an attempt to make it cheaper, do it less often, or avoid it altogether.
\end{keybox}

\section{What changes, and what does not}\label{ec:sec:orient}

Let $E$ be an elliptic curve over $\mathbb{F}_p$, $P$ a point of prime order $r$, and
$Q=[k]P$ for an unknown $k$. Recovering $k$ is the \emph{elliptic-curve discrete logarithm problem}
(ECDLP), and it is what ECDSA, ECDH and every deployed curve-based scheme rest on.

Shor's algorithm for it \cite{ec:proos-zalka} uses two counting registers instead of one:
\[
\frac{1}{2^{m}}\sum_{u,v}\ket u\ket v\ket{\mathcal O}
\;\xrightarrow{\ \text{oracle}\ }\;
\frac{1}{2^{m}}\sum_{u,v}\ket u\ket v\ket{[u]P+[v]Q}
\;\xrightarrow{\ F^\dagger\otimes F^\dagger\ }\;
\text{measure } (j_1,j_2).
\]
The third register depends on $(u,v)$ only through $u+kv \bmod r$, so its level sets are cosets of
the lattice generated by $(-k,1)$ and $(r,0)$. The Fourier transform concentrates on the dual
lattice, and a measured pair obeys, with $a_i=\mathrm{round}(j_i r/2^{m})$,
\begin{equation}\label{ec:eq:post}
a_2 \equiv k\,a_1 \pmod r,
\qquad\text{so}\qquad
k = a_2\,a_1^{-1} \bmod r
\end{equation}
whenever $a_1$ is invertible --- which discards $j_1=0$ automatically. One measurement suffices with
constant probability; \S\ref{ec:sec:verify} shows the distribution this produces on a toy curve, and
\S\ref{ec:sec:post} the search that \cite{ec:ekera} adds when $2^m$ is not a multiple of $r$.

\begin{center}
\begin{quantikz}[column sep=8pt,row sep=8pt]
\lstick{$\ket{u_0}$}     & \gate{H} & \ctrl{6} & \qw          & \ \ldots\ & \qw                & \qw       & \ \ldots\ & \qw                & \gate[3]{F^\dagger} & \meter{}\\
\lstick{$\vdots$}        &          &          &              &           &                    &           &           &                    &                     & \\
\lstick{$\ket{u_{m-1}}$} & \gate{H} & \qw      & \qw          & \ \ldots\ & \ctrl{4}           & \qw       & \ \ldots\ & \qw                &                     & \meter{}\\
\lstick{$\ket{v_0}$}     & \gate{H} & \qw      & \qw          & \ \ldots\ & \qw                & \ctrl{3}  & \ \ldots\ & \qw                & \gate[3]{F^\dagger} & \meter{}\\
\lstick{$\vdots$}        &          &          &              &           &                    &           &           &                    &                     & \\
\lstick{$\ket{v_{m-1}}$} & \gate{H} & \qw      & \qw          & \ \ldots\ & \qw                & \qw       & \ \ldots\ & \ctrl{1}           &                     & \meter{}\\
\lstick{$\ket{S}$}       & \qw      & \gate{+P}& \gate{+[2]P} & \ \ldots\ & \gate{+[2^{m-1}]P} & \gate{+Q} & \ \ldots\ & \gate{+[2^{m-1}]Q} & \qw                 & \qw
\end{quantikz}
\end{center}

\noindent
Each rung adds a \emph{classical} point --- the multiples $[2^i]P$ and $[2^i]Q$ are precomputed on a
laptop --- controlled by one counting qubit. There are $2m\approx 2n$ of them, against the $2n$
rungs of Part~II's ladder, so at the level of \emph{counting} nothing has changed. The cost is
entirely inside the box marked $+P$.

\begin{warnbox}[title={The accumulator never starts at infinity}]
The diagram starts the third register at a fixed point $S$, not at the identity $\mathcal O$. Affine
addition formulas cannot represent the point at infinity at all: there is no $(x,y)$ for it, and the
slope $\lambda$ is undefined the moment the accumulator equals $\pm$ the addend. Starting at a
constant $S$ and reading the answer relative to it keeps the accumulator on an ordinary point
throughout, and shifts the computed coset without touching the interference. The same device
appears as the offset $T$ in the windowed constructions of \S\ref{ec:sec:window}, and as
\code{P0} in the Classiq tutorial \cite{ec:classiq}.
\end{warnbox}

\section{Why it works}\label{ec:sec:why}

Section~\ref{ec:sec:orient} asserted \eqref{ec:eq:post}. This section derives it. The argument is
short, it is the only place in this part where the \emph{quantum} content lives --- everything after
\S\ref{ec:sec:grouplaw} is reversible arithmetic --- and it is checked numerically in
\code{bench/ec\_ablation.py}, which evaluates the amplitudes by brute force and compares them
against the closed form below.

\subsection{The state, in four lines}

Write $q=2^m$ for the size of each counting register and $\omega=e^{2\pi i/q}$. Start from a uniform
superposition and apply the oracle:
\[
\frac1q\sum_{u,v}\ket u\ket v\ket S
\;\longrightarrow\;
\frac1q\sum_{u,v}\ket u\ket v\,\big|\,S+[u]P+[v]Q\,\big\rangle
=\frac1q\sum_{u,v}\ket u\ket v\,\big|\,S+[u+kv]P\,\big\rangle ,
\]
the last equality being the only place $Q=[k]P$ is used. The third register therefore depends on
$(u,v)$ \emph{only} through
\[
w \;=\; u+kv \bmod r ,
\]
so group the terms by $w$:
\begin{equation}\label{ec:eq:grouped}
\frac1q\sum_{w\in\mathbb{Z}_r}\Bigg(\sum_{(u,v)\,:\,u+kv\equiv w}\ket u\ket v\Bigg)\big|\,S+[w]P\,\big\rangle .
\end{equation}
The third register is never measured and never needs to be: \eqref{ec:eq:grouped} is already a sum
of orthogonal branches, so the reduced state on $(u,v)$ is a mixture over $w$, and it is enough to
analyse one branch.

\begin{keybox}[title={Why the offset $S$ cannot matter}]
$S$ appears only inside the third register's label. It relabels which basis state each branch
attaches to; it does not touch the $(u,v)$ amplitudes at all, so it cannot change the measured
distribution. The same argument covers $w$: every branch of \eqref{ec:eq:grouped} gives the same
$|$amplitude$|$, which the benchmark confirms directly by computing two different branches and
comparing (result: \code{yes}). $S$ is there for the arithmetic --- to keep the
accumulator off the point at infinity --- and for nothing else.
\end{keybox}

\subsection{Two lattices}

Fix a branch $w$. Its $(u,v)$ set is a coset of
\[
L \;=\; \{(u,v)\in\mathbb{Z}^2 \;:\; u+kv\equiv 0 \ (\mathrm{mod}\ r)\}
\;=\; \big\langle\, (r,0),\ (-k,1) \,\big\rangle ,
\]
the \emph{period lattice}. The two generators say the two things that can be done without changing
the answer: take $r$ steps in $u$, or take one step in $v$ and $k$ steps back in $u$. Its
determinant is $r$, so $[\mathbb{Z}^2 : L]=r$ and the plane splits into exactly $r$ cosets --- one
per value of $w$, as it must.

\emph{The discrete logarithm is the slope of the second generator.} Finding $L$ is finding $k$;
that is the whole problem, restated.

A Fourier transform does not see $L$, it sees the \emph{dual}
\[
L^{*} \;=\; \{\,y\in\mathbb{Q}^2 \;:\; \langle y,x\rangle\in\mathbb{Z}\ \ \forall x\in L\,\}
\;=\; \big\langle\, \tfrac1r(1,k),\ (0,1) \,\big\rangle ,
\]
obtained as $(B^{-1})^{\mathsf T}$ from $L$'s basis $B$. Writing a point of $L^{*}$ as
$(a_1/r,\,a_2/r)$ and clearing denominators, membership says exactly
\[
a_2 \equiv k\,a_1 \ (\mathrm{mod}\ r) ,
\]
which is \eqref{ec:eq:post}. It is not a consequence of the dual lattice --- it \emph{is} the dual
lattice, written in coordinates.

\subsection{The amplitude}

Take $r\mid q$ for a moment and write $q=rM$. Then for each $v$ the admissible $u$ in $[0,q)$ are
$u_0(v)+tr$ with $u_0(v)=(w-kv)\bmod r$ and $t=0,\dots,M-1$ --- exactly $M$ of them. The amplitude
at $(j_1,j_2)$ after the transform is, up to normalisation,
\[
A(j_1,j_2)\;\propto\;\sum_{v=0}^{q-1}\omega^{v j_2}\,\omega^{u_0(v)j_1}
\underbrace{\sum_{t=0}^{M-1}\omega^{t r j_1}}_{\text{(i)}} .
\]
Sum (i) is geometric with ratio $\omega^{rj_1}$: it equals $M$ when $rj_1\equiv 0 \pmod q$, i.e.\
when $M \mid j_1$, and \emph{exactly zero} otherwise. So only $j_1=a_1M$ survives. For those,
$\omega^{u_0 j_1}=e^{2\pi i\,u_0 a_1/r}$ depends on $u_0$ only mod $r$ --- which is why reducing
$u_0$ into $[0,r)$ above was harmless --- and $u_0\equiv w-kv$, leaving
\[
A \;\propto\; e^{2\pi i\,w a_1/r}\sum_{v=0}^{q-1} e^{2\pi i\,v\,(j_2-k a_1 M)/q} .
\]
That second sum is $q$ when $j_2\equiv ka_1M \pmod q$ and zero otherwise. Writing $j_2=a_2M$ and
dividing through by $M$:
\[
\boxed{\;a_2\equiv k\,a_1 \pmod r\;}
\]
and every surviving outcome has $|A|^2=1/r$. The distribution is \emph{uniform on the $r$ points of
the dual lattice}, and nothing else has any amplitude at all.

\begin{center}\small
\begin{tabular}{@{}rrrrrcr@{}}
\toprule
$r$ & $k$ & $q$ & outcomes & each & closed form exact & usable $\varphi(r)/r$\\
\midrule
$4$ & $3$ & $16$ & $4$ & $1/4$ & \checkmark & $0.50$\\
$8$ & $3$ & $32$ & $8$ & $1/8$ & \checkmark & $0.50$\\
$4$ & $1$ & $64$ & $4$ & $1/4$ & \checkmark & $0.50$\\
\bottomrule
\end{tabular}
\end{center}

\subsection{Reading off $k$, and what it costs}

Given $(a_1,a_2)$ with $a_2\equiv ka_1$, recover $k=a_2a_1^{-1}\bmod r$ --- possible exactly when
$a_1$ is invertible. For prime $r$ that means $a_1\neq0$, so $r-1$ of the $r$ outcomes work and one
run succeeds with probability $(r-1)/r$: overwhelming, and it does not improve with more qubits
because the $a_1=0$ outcome is a fixed $1/r$ of the mass. Note $r$ is not secret --- it is the
curve's published group order --- so this is a division, not a search.

\begin{warnbox}[title={The clean case never actually happens}]
Everything above assumed $r\mid q$. For ECDLP $r$ is \emph{prime}, and $q$ is a power of two, so
$r\mid q$ only when $r=2$. The exact calculation is a device for seeing the structure; the real
situation is always the other one.

What goes wrong is only sum (i): with $r\nmid q$ it no longer vanishes, it becomes a Dirichlet
kernel. Amplitude leaks onto the integers either side of $a_1q/r$, which is why the measured peaks
of \S\ref{ec:sec:post} sit at $(6,5)$ and $(2,3)$ rather than at the exact dual-lattice points
$(1.6,3.2)$ and $(3.2,6.4)$, and why $a_i=\mathrm{round}(j_ir/q)$ has a rounding in it at all. The
leak shrinks as $q$ grows past $r$. Entries are $P(\text{usable and correct})$, i.e.\ the chance
of an outcome with $a_1$ invertible \emph{and} $a_2\equiv ka_1$; the last column is the irreducible
$a_1=0$:

\begin{center}\small
\begin{tabular}{@{}rrrrrr@{}}
\toprule
& \multicolumn{4}{c}{$m-\lceil\log_2 r\rceil$} & \\
\cmidrule(lr){2-5}
$r$ & $0$ & $1$ & $2$ & $3$ & $P(a_1{=}0)$\\
\midrule
$5$ & $60\%$ & $67\%$ & $74\%$ & $77\%$ & $20\%$\\
$7$ & $48\%$ & $67\%$ & $77\%$ & $81\%$ & $15\%$\\
$11$ & $69\%$ & $77\%$ & $84\%$ & $88\%$ & $9\%$\\
$13$ & $63\%$ & $75\%$ & $84\%$ & $88\%$ & $8\%$\\
\bottomrule
\end{tabular}
\end{center}

\noindent
At $m=\lceil\log_2 r\rceil$ the answer is barely better than a coin; three extra qubits per register
take it to $88\%$, with the residual $8$ being the irreducible $a_1=0$ outcome. This
is the whole content of \cite{ec:ekera}'s analysis for the discrete logarithm, and the reason the
post-processing of \S\ref{ec:sec:post} searches a small neighbourhood of the rounded value rather
than trusting it.
\end{warnbox}

\section{The group law, and why it resists reversibility}\label{ec:sec:grouplaw}

For $P_1=(x_1,y_1)$ and $P_2=(x_2,y_2)$ on $y^2=x^3+ax+b$ with $x_1\neq x_2$,
\begin{equation}\label{ec:eq:add}
\lambda=\frac{y_1-y_2}{x_1-x_2},\qquad
x_3=\lambda^2-x_1-x_2,\qquad
y_3=\lambda(x_2-x_3)-y_2 .
\end{equation}
Three subtractions, one squaring, one multiplication --- and one division. In Part~II every
arithmetic step was addition or multiplication by a compile-time constant. Here the divisor
$x_1-x_2$ is a quantum value, so the inverse must be computed \emph{by circuit}, on a
superposition, and then uncomputed. Section~\ref{ec:sec:kaliski} does that, and
\S\ref{ec:sec:cost} shows it dominating everything else by a factor of three to five.

\begin{center}\small\setlength{\tabcolsep}{5pt}
\begin{tabular}{@{}rllr@{}}
\toprule
Level & Object & Made of & In code\\
\midrule
5 & the two ladders & $2m$ rungs, one per counting qubit & \code{ec\_shor.py}\\
4 & $c$-$(+P_i)$ & 4 sub, 1 add, 2 div, 2 mul, 1 sqr, 1 neg & \code{ec\_pointadd.py}\\
3 & $x^{-1}\bmod p$ & $2n$ Kaliski rounds + correction & \code{ec\_kaliski.py}\\
3$'$ & $x\cdot y\bmod p$ & $n$ modular doublings + $n$ additions & \code{ec\_mult.py}\\
2 & $\pm\bmod p$, $2x\bmod p$ & plain add/sub plus one reduction & \code{ec\_modarith.py}\\
1 & $\ADD$ & CDKM chain, or Gidney's AND chain & \code{ec\_adders.py}\\
0 & AND & $4T$ to compute, $0T$ to uncompute & \code{ec\_gates.py}\\
\bottomrule
\end{tabular}
\end{center}

\noindent
Level~3 is new and it is the expensive one. Everything below it is Part~II's tower with the
modulus renamed.

\section{Level 2: modular arithmetic over $\mathbb{F}_p$}\label{ec:sec:modarith}

Two blocks carry everything above them, and both are from Roetteler, Naehrig, Svore and Lauter
\cite{ec:rnsl17}. Neither is hard; both are worth seeing, because \S\ref{ec:sec:approx} makes its
savings by attacking exactly the parts described here as ``the awkward bit''.

\paragraph{Modular doubling.}
Shift, subtract $p$, add it back if that went negative:
\[
2x \;\longrightarrow\; 2x-p \;\longrightarrow\; \begin{cases} 2x-p & 2x\ge p\\ 2x & \text{otherwise.}\end{cases}
\]
The shift is free --- it renames wires. The awkward bit is the borrow flag, which must return to
$\ket0$ or it is garbage in the sense of \S\ref{sec:garbage}. It clears for one CNOT: $2x$ is even
and $p$ is odd, so \emph{the answer is even exactly when no reduction happened}, and the answer's
own low bit is the flag.

\paragraph{Modular addition.} The same shape, but the flag-clearing trick is prettier.
After adding and subtracting $p$, the flag holds $[\,x+y<p\,]$, i.e.\ ``no reduction''. And
\begin{itemize}\setlength\itemsep{1pt}
\item no reduction means the answer is $x+y$, which is $\ge x$;
\item a reduction means the answer is $x+y-p$, which is $<x$ because $y<p$.
\end{itemize}
So $[\,\text{answer}<x\,]$ is exactly the complement of the flag, and one comparison against an
operand that is still lying there clears it. No extra ancilla, no second copy. In code:

\begin{lstlisting}
def cmodadd(m, ctrl, x, y, p):
    """y <- (y + x) mod p when ctrl (or unconditionally if ctrl is None).
    ..."""
    ctx, n = m.ctx, len(y)
    assert len(x) == n
    anc_x, anc_y = m.anc(1, "ax"), m.anc(1, "ay")
    xe, ye = x + anc_x, y + anc_y
    cp, a = m.anc(n + 1, "cp"), m.anc(n + 1, "a")

    if ctrl is None:
        A.add(ctx, xe, ye, a)
    else:
        A.cadd(ctx, ctrl, xe, ye, cp, a)
    A.sub_const(ctx, ye, p, cp, a)
    A.cadd_const(ctx, anc_y[0], y, p, cp[:n], a)
    if ctrl is None:
        A.lt_uint(ctx, y, x, anc_y[0], a)
    else:
        A.clt_uint(ctx, ctrl, y, x, anc_y[0], a, cp)
    ctx.x(anc_y[0])
    m.free(cp, a, anc_x, anc_y)
\end{lstlisting}

\noindent
The doubling is the same shape and its flag-clearing is one CNOT --- the last two lines below,
where \code{xe[0]} is the answer's own low bit:

\begin{lstlisting}
def moddbl(m, x, p):
    """x <- 2x mod p, in place: the register handle stays valid.
    ..."""
    ctx, n = m.ctx, len(x)
    z = m.anc(1, "z")
    xe = A.shift_up(x, z[0])                 # n+1 bits, holds 2x
    _moddbl_body(m, xe, p)
    # value bits now sit at xe[0..n-1]; xe[n] (the old MSB) is clean
    for i in range(n - 1, -1, -1):
        ctx.swap(xe[i], xe[i + 1])
    m.free(z)
\end{lstlisting}

\section{Level 3: inversion by Kaliski's algorithm}\label{ec:sec:kaliski}

The binary extended Euclidean algorithm branches four ways per step: $u$ even, $v$ even, both odd
with $u>v$, both odd with $u\le v$. A quantum circuit cannot branch --- every branch must run on
every input, on a superposition. Kaliski's formulation as reworked by H\"aner, Jaques, Naehrig,
Roetteler and Soeken \cite{ec:hjnrs20} makes all four cases the \emph{same} straight line:

\begin{center}\small
\begin{minipage}{0.86\textwidth}
\begin{tabbing}
xxx\=xxxx\=xxxxxxxxxxxxxxxxxxxxxxxxxxxxx\=\kill
\> \textbf{repeat} $2n$ times:\\
\> 1.\> \textbf{if} $v$ odd \textbf{and} ($u$ even \textbf{or} $u>v$): \> swap $(u,v)$, swap $(r,s)$; remember it\\
\> 2.\> \textbf{if} $u,v$ both odd: \> $v \mathrel{-}= u$; \quad $s \mathrel{+}= r$\\
\> 3.\> \> $v \mathrel{/}= 2$; \quad $r \mathrel{\times}= 2$ \ (the doubling only if the round fired)\\
\> 4.\> \textbf{if} it swapped: \> swap back
\end{tabbing}
\end{minipage}
\end{center}

\noindent
Started at $(u,v,r,s)=(p,x,0,1)$ this terminates with $u=1$, $v=0$, $s=p$ and
\begin{equation}\label{ec:eq:pseudo}
r \;=\; -x^{-1}\,2^{k} \bmod p,
\end{equation}
where $k$ counts the rounds that did anything. That is a \emph{pseudo}-inverse: right up to a sign
and a power of two, both of which the caller fixes --- one modular negation, and $k$ modular
halvings driven by a counter. Equation~\eqref{ec:eq:pseudo} is asserted directly in
\code{tests/test\_ec\_kaliski.py} for every invertible $x$ and every $p$ tested, alongside
$u=1$, $v=0$, $s\equiv0$ and the counter itself.

\begin{keybox}[title={Three bits per round, and why none can be dropped}]
The round is a permutation of $(u,v,r,s)$ chosen by its decision bits, so those bits must survive
or the round is not reversible. Three are kept:
\begin{itemize}\setlength\itemsep{1pt}
\item $f$ --- \emph{did this round do anything} ($v\neq0$ on entry). Not recoverable afterwards: an
inert round and the last active round both leave $v=0$.
\item $a$ --- \emph{did it swap}.
\item $b$ --- \emph{were both odd}, so the subtract-and-add fired. Not implied by $(f,a)$: the case
$f=1,a=0$ covers both ``$v$ was even'' and ``both odd with $u\le v$''.
\end{itemize}
The comparison $u>v$ is \emph{not} kept. It is consumed while building $a$ and then uncomputed by
running the comparator a second time --- before $u$ and $v$ change, which is the only window in
which that is possible.
\end{keybox}

\noindent
In code, with the decision bits built first and the comparison uncomputed before $u$ and $v$ move:

\begin{lstlisting}
def kaliski_round(m, u, v, r, s, p, rec, conditional=True):
    ctx, n = m.ctx, len(u)
    f, a, b = rec

    # --- f: is this round active? ---------------------------------------
    if conditional:
        MA.is_nonzero(m, v, f[0])
    else:
        ctx.x(f[0])                       # [106] Sec 3.3: never test v != 0

    # --- decision bits, from the pre-round u and v ------------------------
    ...

    # --- the straight-line body ------------------------------------------
    for i in range(n):
        ctx.cswap(a[0], u[i], v[i])
    for i in range(len(r)):
        ctx.cswap(a[0], r[i], s[i])

    cp, sc = m.anc(n, "cp"), m.anc(n, "sc")
    A.csub(ctx, b[0], u, v, cp, sc)       # v -= u  (no borrow: u <= v here)
    m.free(cp, sc)
    MA.cmodadd(m, b[0], r, s, p)          # s += r  mod p

    _shift_down(ctx, v)                   # v /= 2  (v is even by now)
    MA.cmoddbl(m, f[0], r, p)             # r *= 2  mod p

    for i in range(n):
        ctx.cswap(a[0], u[i], v[i])
    for i in range(len(r)):
        ctx.cswap(a[0], r[i], s[i])
\end{lstlisting}

\noindent
$r$ and $s$ are reduced mod $p$ every round. They have to be: $r$ doubles each round, and the tests
confirm it runs past $p$ within a few rounds if left alone. $u$ and $v$ are plain $n$-bit integers,
staying in $[0,p]$ throughout.

\paragraph{Division, and where the garbage goes.}
What a point addition actually wants is $\lambda = y/x$, and the inversion above leaves $6n$ qubits
of per-round records behind. Both are solved at once by the shape of \cite{ec:rnsl17}: invert,
multiply, copy the answer out, then \emph{run the whole thing backwards}. Only the copy survives;
records, counter, $u$, $v$ and $s$ all return to $\ket0$, which the simulator checks at the exact
instruction each is freed.

\begin{lstlisting}
def mod_div(m, ctrl, x, y, out, p):
    """out ^= ctrl * (y / x mod p).  [106] Fig. 7, the modular division.
    ..."""
    n = len(x)
    lam = m.anc(n, "lam")

    def build():
        inv = m.anc(n, "inv")
        recs, cnt = mod_inv(m, x, inv, p, ctrl=ctrl)
        import ec_mult as MU
        MU.modmul(m, y, inv, lam, p)
        return inv, recs, cnt

    # forward: lam = ctrl * y * x^{-1}
    start = len(m.qc.data)
    inv, recs, cnt = build()
    body = list(m.qc.data[start:])

    # lam is already 0 on the ctrl = 0 branch, so the copy needs no control
    for i in range(n):
        m.ctx.cx(lam[i], out[i])

    for ci in reversed(body):                         # undo everything else
        m.qc.append(ci.operation.inverse(), ci.qubits, ci.clbits)
    m.free(lam, inv)
    for rec in recs:
        m.free(*rec)
    if cnt is not None:
        m.free(cnt)
\end{lstlisting}

\begin{warnbox}[title={A control that has to reach all the way down --- here, though not everywhere}]
The division is controlled by the counting qubit $q$, and for the inversion built above it is not
enough to gate only the final copy. On the $q=0$ branch the $x$ register does not hold $x_1-x_2$;
by that point in the point addition it holds $x_1+2x_2$, which vanishes for perfectly ordinary
accumulators, so the circuit would be inverting zero --- and Kaliski on $0$ terminates with $u=p$
and $s=1$, not the $u=1$, $s\equiv0$ that the cleanup here relies on to return those registers to
$\ket0$.

The fix costs almost nothing. Load the inversion's \emph{setup} under the control ---
$u\leftarrow p$, $v\leftarrow x$, $s\leftarrow1$ --- and when $q=0$ every register starts at zero,
so $v=0$ makes all $2n$ rounds inert, $r$ stays $0$, and the correction does nothing. The rounds
themselves are never controlled; $n$ ANDs and a handful of CNOTs buy the whole thing, and it is
checked directly: with $q=0$ the output register is $0$ for every input.

\cite{ec:kim26} does \emph{not} do this, and does not need to. Its Fig.~7 runs the inversion and
the multiplication unconditionally and gates only the copy-out of $\lambda$, which is sound there
because the inversion is undone by its literal inverse --- whatever Inv did to the registers on a
degenerate input, Inv$^\dagger$ undoes. The gating above is the price of \emph{this} construction's
cheaper cleanup, not a correction to the paper.
\end{warnbox}

\section{Level 4: the eleven-step point addition}\label{ec:sec:padd}

The accumulator $(x_1,y_1)$ is quantum, the addend $(x_2,y_2)$ classical, and the control $q$ says
whether to add at all. In place, no garbage:

\begin{center}\small\setlength{\tabcolsep}{6pt}
\begin{tabular}{@{}rlll@{}}
\toprule
& operation & \multicolumn{2}{l}{after it, on the $q=1$ branch:}\\
\cmidrule(l){3-4}
& & $x$ register & $y$ register\\
\midrule
1 & $x \mathrel{-}= x_2$                      & $x_1-x_2$ & $y_1$\\
2 & $y \mathrel{-}= q\,y_2$                   & & $y_1-y_2$\\
3 & $\lambda \mathrel{\oplus}= q\cdot y/x$    & & \\
4 & $y \mathrel{\oplus}= x\lambda$            & & $\mathbf{0}$\\
5 & $x \mathrel{+}= 3x_2$                     & $x_1+2x_2$ & \\
6 & $x \mathrel{-}= \lambda^2$                & $x_2-x_3$ & \\
7 & $y \mathrel{+}= x\lambda$                 & & $\lambda(x_2-x_3)$\\
8 & $\lambda \mathrel{\oplus}= q\cdot y/x$    & & \\
9 & $x \mathrel{-}= (2-q)x_2$                 & $-x_3$ & \\
10 & $y \mathrel{-}= q\,y_2$                  & & $y_3$\\
11 & $x \leftarrow -x$ \ if $q$               & $x_3$ & \\
\bottomrule
\end{tabular}
\end{center}

\noindent
The circuit is the table, line for line:

\begin{lstlisting}
def point_add_ctrl(m, q, x1, y1, x2, y2, p):
    """(x1, y1) <- (x1, y1) + (x2, y2) when q; unchanged when q = 0.
    ..."""
    n = len(x1)
    lam = m.anc(n, "lam")

    MA.modsub_const(m, x1, x2, p)                 # 1  x <- x1 - x2
    MA.cmodsub_const(m, q, y1, y2, p)             # 2  y <- y1 - q y2
    K.mod_div(m, q, x1, y1, lam, p)               # 3  lambda <- q y/x
    MU.modmul_xor(m, x1, lam, y1, p)              # 4  y <- y XOR x lambda  (= 0)
    MA.modadd_const(m, x1, 3 * x2 % p, p)         # 5  x <- x + 3 x2
    MU.modsqr_sub(m, lam, x1, p)                  # 6  x <- x - lambda^2
    MU.modmul_add(m, x1, lam, y1, p)              # 7  y <- y + x lambda
    K.mod_div(m, q, x1, y1, lam, p)               # 8  lambda <- 0
    _csub_2x2_or_x2(m, q, x1, x2, p)              # 9  x <- -x3
    MA.cmodsub_const(m, q, y1, y2, p)             # 10 y <- y3
    MA.cmodneg(m, q, x1, p)                       # 11 x <- x3

    m.free(lam)
\end{lstlisting}

\noindent
Three things in that table are worth pausing on.

\emph{Step 4 is an XOR, not a subtraction.} At that moment $y$ holds $y_1-y_2$ and $x\lambda$ is by
definition $y_1-y_2$, so XOR-ing the product in zeroes the register, and the algebra guarantees it.
What is saved is the \emph{accumulation} --- CNOTs instead of a modular subtraction. The product
$x\lambda$ still has to be computed and uncomputed around it, which is two modular multiplications:
the step is cheaper than the alternative, not cheap.

\emph{Step 8 is the same circuit as step 3, and it \emph{clears} $\lambda$.} By then the registers
hold $x=x_2-x_3$ and $y=\lambda(x_2-x_3)$, so $y/x$ is again $\lambda$, and the XOR-accumulating
division undoes what step 3 wrote. The whole eleven-step arrangement exists to make that identity
true; without it $\lambda$ would be $n$ qubits of garbage per addition.

\emph{Step 9 is two of \cite{ec:hjnrs20}'s steps merged}, following \cite{ec:kim26}. The original
does $x\mathrel{-}=2x_2$ then a controlled $x\mathrel{+}=x_2$; the merged form subtracts
$(2-q)x_2$ once. Building that operand is free: $2x_2\bmod p$ and $x_2\bmod p$ are both classical,
so the register is X-gated to one of them and CNOTs from $q$ patch the bits where they differ.

\begin{warnbox}[title={Three exceptional cases, not one}]
The circuit is correct except on a set of density $O(1/p)$, and it does not detect the exceptions
--- a check would cost more than it saves. There are three, and only the first is the one usually
quoted:
\begin{itemize}\setlength\itemsep{1pt}
\item $x_1=x_2$ --- step 3 divides by $x_1-x_2$. Doubling, adding an inverse, or infinity.
\item $x_3=x_2$, i.e.\ $P_1=-[2]P_2$ --- step 8 divides by $x_2-x_3$. This one is invisible in
\eqref{ec:eq:add} and appears only once you follow what the \emph{registers} hold. The in-place
multiplier of \S\ref{ec:sec:dialog} reaches the same wall from the other direction, where it would
be asked to multiply by zero.
\item for the \emph{windowed} circuit of \S\ref{ec:sec:window} adding $\mathcal O$: $x_R=0$. This
one is not about the table above --- it belongs to \cite{ec:schrott26}'s Algorithm~1, whose own
step~6 divides by $x_R-x_{P_i}$, and that is $x_R$ when $P_i=\mathcal O$ is encoded as $(0,0)$.
Adding the identity looks like it could never fail.
\end{itemize}
\code{ec\_classical.point\_add\_exceptional} decides the first two; the tests exclude them
explicitly and report how many, rather than passing quietly on inputs they never exercise.
\end{warnbox}

\section{Making it fast}\label{ec:sec:fast}

Everything from here on is measured. Each subsection fixes a \emph{task} and varies one
implementation choice; comparing a modular doubling against a point addition would say nothing, so
nothing here does. Every variant is re-checked against the classical model \emph{inside} the
benchmark before its cost is recorded --- a cheaper circuit that is wrong is not a datapoint ---
and variants that are deliberately approximate record a measured failure rate instead. All of it
comes from \code{bench/ec\_ablation.py}, and the tables below are generated from its output rather
than typed.

\begin{keybox}[title={Which Toffoli count, and why it matters here more than in Part V}]
Both source papers count ``Toffolis'' as CCX, CCZ and AND gates together
\cite{ec:schrott26}. A measurement-based AND$^\dagger$ contains \emph{no} Toffoli and no $T$ gate,
so it is not counted. That is the convention used in every table below, and it is the one to
compare against published figures.

It is not the only defensible convention --- one could count every three-qubit operation performed,
uncomputes included --- and the two differ by nearly $2\times$ on Gidney-style circuits, which is
most of what this part builds. \code{ec\_cost} reports both (\code{toffoli\_paper} and
\code{toffoli\_equiv}) precisely so the choice is visible rather than buried. Quoting the wrong one
would make every optimisation here look roughly half as good as the literature says it is.
\end{keybox}

\subsection{Level 1: which adder, and how to control it}\label{ec:sec:fastadd}

Two adders, chosen per call site by how many ancillas happen to be free --- as \cite{ec:schrott26}
puts it, the ancilla-hungry one where ancillas are available and the ancilla-free one where they
are not.

\begin{center}\small
\begin{tabular}{@{}rrrrrrrr@{}}
\toprule
& \multicolumn{2}{c}{addition} & \multicolumn{3}{c}{controlled addition} & \multicolumn{2}{c}{qubits}\\
\cmidrule(lr){2-3}\cmidrule(lr){4-6}\cmidrule(lr){7-8}
& & & & \multicolumn{2}{c}{copy-then-add, with} & & \\
\cmidrule(lr){5-6}
$n$ & CDKM & Gidney & reconstruct & CDKM & Gidney & CDKM & Gidney\\
\midrule
$8$ & $16$ & $7$ & $32$ & $24$ & $15$ & $17$ & $23$\\
$16$ & $32$ & $15$ & $64$ & $48$ & $31$ & $33$ & $47$\\
$32$ & $64$ & $31$ & $128$ & $96$ & $63$ & $65$ & $95$\\
\bottomrule
\end{tabular}
\end{center}

\noindent
Plain addition is the $2n$-versus-$n$ of the textbook comparison, and lands exactly there:
$64=2n$ against $31=n-1$ at $n=32$. CDKM spends two Toffolis per bit, Gidney one AND
per bit, and the difference is that Gidney's targets start clean, so they uncompute for free --- the
$n-1$ rather than $n$ being the carry into the top bit, which no adder has to compute.

Controlled addition is where \cite{ec:kim26} makes its choice. Promoting the adder's CNOTs to
Toffolis costs $128$ at $n=32$; copying the operand into clean ancillas \emph{under the
control} and then adding unconditionally costs $63$, a factor $2.0$. The copies
all commute, so they also parallelise, which the reconstructed version cannot.

\begin{lstlisting}
def cadd(ctx, ctrl, x, y, copy, anc):
    """y += x if ctrl, via [106] Fig. 4(b): copy under control, then add.
    ..."""
    assert len(copy) == len(x)
    for i, q in enumerate(x):
        ctx.and_(ctrl, q, copy[i])
    add(ctx, copy, y, anc)
    for i, q in enumerate(x):
        ctx.and_dg(ctrl, q, copy[i])
\end{lstlisting}

\noindent
The middle column is there to keep that factor honest. Copy-then-add and the Gidney adder are two
separate decisions, and the $2.0\times$ bundles them; hold the adder fixed at CDKM and
the control style alone is worth $1.33\times$ ($96$ against $128$).
The rest is the adder.

What it costs is qubits. The last two columns are the plain adders' widths, $2n+1$ against
$3n-1$; for the controlled pair the gap is wider still, $66$ against $128$ at
$n=32$, since the copy needs a register of its own.

\subsection{Level 0: temporary ANDs}\label{ec:sec:fastand}

Part~V measured this for factoring; it applies unchanged, and matters more here because there is so
much more arithmetic to apply it to. Over one whole modular inversion:

\begin{center}\small
\begin{tabular}{@{}rrrrrr@{}}
\toprule
$n$ & ANDs & AND$^\dagger$s & $T$ (all Toffoli) & $T$ (temporary AND) & saving\\
\midrule
$5$ & $912$ & $892$ & $14\,938$ & $5\,958$ & $2.51\times$\\
$8$ & $2\,229$ & $2\,197$ & $36\,694$ & $14\,628$ & $2.51\times$\\
$12$ & $4\,785$ & $4\,737$ & $79\,254$ & $31\,740$ & $2.50\times$\\
\bottomrule
\end{tabular}
\end{center}

\noindent
A flat $2.50\times$, and the arithmetic says why: ANDs and AND$^\dagger$s arrive in
near-equal numbers, so the average is $(4+0)/2=2$ against $7$, i.e.\ $3.5\times$ on the ANDs alone,
diluted by the Fredkins and dirty-target Toffolis that cannot use the trick.

\subsection{Level 2: approximate and pseudo-Mersenne reduction}\label{ec:sec:approx}

Shor's circuit only has to work with large constant probability, and \cite{ec:schrott26} spends
that freedom on the reductions. Two ideas.

\emph{Compare only the top bits.} ``Is this at least $p$?'' is answered correctly unless the top
few bits tie, which happens with probability about $2^{-\mathrm{msbs}}$. That shrinks an $n$-bit
comparison to an $\mathrm{msbs}$-bit one and --- more importantly --- lets the flag be produced
\emph{before} the subtraction rather than recovered after it, collapsing the exact two-step
reduction of \S\ref{ec:sec:modarith} into a single controlled subtraction.

\emph{Exploit the shape of the prime.} secp256k1 uses $p=2^{u}-f$ with $f=2^{32}+977$, tiny against
$2^{256}$. Then ``at least $p$?'' is almost ``is bit $u$ set?'', and subtracting $p$ becomes ``clear
bit $u$, then add $f$'' --- and since $f$ is small, adding it need only touch the bottom few bits,
because the carry cannot travel far. An $n$-bit constant subtraction becomes a short one.

The pseudo-Mersenne doubling is the whole idea in eight lines --- note that nothing compares
anything, and that clearing the overflow bit \emph{is} the subtraction of $2^u$:

\begin{lstlisting}
def moddbl_pm(m, x, q, lsbs=None):
    pm = pseudo_mersenne(q)
    assert pm, f"q = {q} is not pseudo-Mersenne"
    u, f = pm
    ctx, n = m.ctx, len(x)
    assert u == n, "expected q just below 2^n"
    lsbs = lsbs or min(n, 2 * max(1, f.bit_length()) + 8)

    z = m.anc(1, "z")
    xe = A.shift_up(x, z[0])                      # n+1 bits; xe[n] is the overflow
    cp, sc = m.anc(lsbs, "cp"), m.anc(lsbs, "sc")
    A.cadd_const(ctx, xe[n], Reg(list(xe[:lsbs])), f, cp, sc)
    m.free(cp, sc)
    ctx.cx(xe[0], xe[n])                          # clears the bit == subtracts 2^u
    for i in range(n - 1, -1, -1):
        ctx.swap(xe[i], xe[i + 1])
    m.free(z)
\end{lstlisting}

\begin{center}\small
\begin{tabular}{@{}rlrrrr@{}}
\toprule
& & \multicolumn{2}{c}{doubling} & \multicolumn{2}{c}{controlled addition}\\
\cmidrule(lr){3-4}\cmidrule(lr){5-6}
$p$ & reduction & Toffoli & failure & Toffoli & failure\\
\midrule
$251$ & exact (Alg.~5 / Alg.~8) & $17$ & $0.0\%$ & $49$ & $0.0\%$\\
$251$ & top 4 bits only (Alg.~6 / Alg.~9) & $13$ & $2.4\%$ & $39$ & $7.0\%$\\
$251$ & pseudo-Mersenne (Alg.~7 / Alg.~10) & $5$ & $3.2\%$ & $33$ & $3.7\%$\\
$1021$ & exact (Alg.~5 / Alg.~8) & $21$ & $0.0\%$ & $61$ & $0.0\%$\\
$1021$ & top 4 bits only (Alg.~6 / Alg.~9) & $15$ & $3.0\%$ & $45$ & $9.7\%$\\
$1021$ & pseudo-Mersenne (Alg.~7 / Alg.~10) & $4$ & $3.1\%$ & $41$ & $1.3\%$\\
$4093$ & exact (Alg.~5 / Alg.~8) & $25$ & $0.0\%$ & $73$ & $0.0\%$\\
$4093$ & top 4 bits only (Alg.~6 / Alg.~9) & $17$ & $3.1\%$ & $51$ & $9.7\%$\\
$4093$ & pseudo-Mersenne (Alg.~7 / Alg.~10) & $4$ & $3.1\%$ & $49$ & $1.0\%$\\
\bottomrule
\end{tabular}
\end{center}

\noindent
Read the failure column as a rate, not a defect: these circuits are \emph{supposed} to be wrong
sometimes. The question is whether the rate tracks $2^{-\mathrm{msbs}}$, and it does --- at
$p=4093$, four MSBs give $3.1\%$ and eight give $0.17\%$, a factor $18$ for four
bits. The pseudo-Mersenne doubling is the headline: $4$ Toffolis against $25$
exact, a $6.2\times$ reduction, for a $3.1\%$ failure rate at this width.

\begin{warnbox}[title={Failure compounds, which is why the paper's budget is so large}]
A single operation failing a few per cent of the time sounds cheap. Chained through the
$\approx1.4n$ iterations of an in-place multiplication it is not. Measured end to end at $p=31$: a
two-bit comparison, which fails about $12\%$ per operation, turns a working multiplier into one that
fails $81\%$ of the time, while the full-width comparison stays exact. That is the calculation behind \cite{ec:schrott26}'s choice of $40$ to $50$
bits at $n=256$ --- per-operation failure near $2^{-40}$, so even several hundred iterations stay
negligible. A cheap approximation is only cheap once you count how often you use it. The measurement
is in \code{tests/test\_ec\_opt1128.py}, not an estimate.
\end{warnbox}

\subsection{Level 3: Montgomery multiplication over a carry-save tree}\label{ec:sec:mont}

\cite{ec:kim26} rewrites word-level Montgomery so that its two halves have the same shape. The
textbook reduction computes $M=c_w p'_w \bmod 2^w$ and then $c \mathrel{+}= Mp$; substituting $M$
gives $c \mathrel{+}= c_w\,(p'_w \bmod 2^w)\,p$, where the bracket depends only on the modulus. So
$d=(p'_w\bmod 2^w)\,p$ is a compile-time constant and the reduction becomes a second
``$c \mathrel{+}= c_w\cdot d$'', structurally identical to the multiplication step. One carry-save
tree then serves both. Its 3:2 compressor turns three addends into two with no carry propagation at
all,
\[
\begin{aligned}
S_i &= A_i \oplus B_i \oplus C_i && (\text{three CNOTs}),\\
M_i &= \MAJ(A_i,B_i,C_i) = A_i \oplus \big((A_i{\oplus}B_i)\wedge(A_i{\oplus}C_i)\big) && (\text{one AND}),
\end{aligned}
\]
with $A+B+C=S+2M$, so reducing $w$ addends to two costs $O(\log w)$ depth and only the final
addition has to propagate a carry.

\begin{lstlisting}
def csa_compress(m, A, B, C):
    """(A,B,C) -> (S, M) with A + B + C == S + M, inputs preserved.
    ..."""
    ctx, W = m.ctx, len(A)
    assert len(B) == len(C) == W
    S, M = m.anc(W, "csaS"), m.anc(W, "csaM")
    for i in range(W):                       # S = A ^ B ^ C
        ctx.cx(A[i], S[i])
        ctx.cx(B[i], S[i])
        ctx.cx(C[i], S[i])
    for i in range(W - 1):                   # M = MAJ, written one place up
        ctx.cx(A[i], B[i])
        ctx.cx(A[i], C[i])
        ctx.and_(B[i], C[i], M[i + 1])
        ctx.cx(A[i], M[i + 1])
        ctx.cx(A[i], B[i])
        ctx.cx(A[i], C[i])
    return S, M
\end{lstlisting}

\begin{center}\small
\begin{tabular}{@{}lrrrr@{}}
\toprule
multiplier & Toffoli & qubits & depth & depth vs.\ schoolbook\\
\midrule
\multicolumn{5}{@{}l}{\emph{$n=8$, $w=2$, $p=251$}}\\
\quad schoolbook double-and-add & $528$ & $46$ & $3\,211$ & ---\\
\quad Montgomery over a carry-save tree & $434$ & $88$ & $2\,600$ & $1.2\times$\\
\quad Montgomery over a QROM lookup & $312$ & $121$ & $1\,676$ & $1.9\times$\\
\multicolumn{5}{@{}l}{\emph{$n=12$, $w=4$, $p=4093$}}\\
\quad schoolbook double-and-add & $1\,176$ & $66$ & $7\,311$ & ---\\
\quad Montgomery over a carry-save tree & $1\,078$ & $268$ & $4\,994$ & $1.5\times$\\
\quad Montgomery over a QROM lookup & $699$ & $261$ & $3\,315$ & $2.2\times$\\
\multicolumn{5}{@{}l}{\emph{$n=16$, $w=4$, $p=65521$}}\\
\quad schoolbook double-and-add & $2\,080$ & $86$ & $13\,043$ & ---\\
\quad Montgomery over a carry-save tree & $1\,566$ & $324$ & $6\,669$ & $2.0\times$\\
\quad Montgomery over a QROM lookup & $1\,054$ & $347$ & $4\,780$ & $2.7\times$\\
\bottomrule
\end{tabular}
\end{center}

\begin{warnbox}[title={What this table does \emph{not} show, and why}]
The carry-save multiplier is a \emph{depth} optimisation, and depth is the column where it wins:
$2.0\times$ over schoolbook at $n=16$. It does not beat the QROM-lookup multiplier
of \cite{ec:hjnrs20} on gate count, and at these widths it should not be expected to --- a tree
over two or four words has almost nothing to flatten, and the rewrite costs real width. The two differ by
$2^w\lfloor c_wp'_w/2^w\rfloor\,p$, which with $c_w$ and $p'_w$ both below $2^w$ reaches about
$2^{2w}p$ --- so the reduction over-adds by that much and the accumulator settles at $n+2w$ bits
rather than the $n+2$ of textbook Montgomery, with a final reduction to mop up. The paper does not dwell
on this; the tests measure the width rather than trusting it.
\end{warnbox}

\subsection{Level 3: dropping the counter and postponing the reduction}\label{ec:sec:kalopt}

Two changes to \S\ref{ec:sec:kaliski}, both from \cite{ec:kim26}, and they compound.

\emph{Run every round unconditionally.} Drop the ``is $v$ still nonzero?'' test. An inert round does
nothing except double $r$, so instead of stopping at $-x^{-1}2^{k}$ the register runs on to
$-x^{-1}2^{2n}$ --- the same value for \emph{every} input, with no counter to keep and no $2n-k$
corrective halvings to perform. And $2^{2n}$ is exactly right: fed $x=X2^{n}$ in Montgomery form the
answer is $-X^{-1}2^{n}$, the Montgomery form of $-X^{-1}$, so the power of two has corrected itself
and only the sign is left --- one modular negation, against \cite{ec:hjnrs20}'s counter and its
$2n-k$ controlled halvings. It also removes one of the three recorded bits per round, since $f$ existed only to
say whether the round fired.

\emph{Postpone the modular reductions.} Do not reduce $r$ and $s$ mod $p$ each round; let them run
into a $2n$-bit register and reduce once at the end. This changes what a round \emph{is}: a modular
doubling (compare, conditionally subtract, clear the flag) becomes a plain left shift, which is a
renaming of wires and costs no Toffoli at all, and a modular addition becomes a plain addition.

\begin{center}\small
\begin{tabular}{@{}rrrrrrrrr@{}}
\toprule
& \multicolumn{3}{c}{Toffoli} & \multicolumn{3}{c}{depth} & \multicolumn{2}{c}{qubits}\\
\cmidrule(lr){2-4}\cmidrule(lr){5-7}\cmidrule(lr){8-9}
$n$ & \cite{ec:hjnrs20} & \cite{ec:kim26} & saved & \cite{ec:hjnrs20} & \cite{ec:kim26} & saved & \cite{ec:hjnrs20} & \cite{ec:kim26}\\
\midrule
$5$ & $1\,242$ & $817$ & $34\%$ & $4\,846$ & $2\,664$ & $45\%$ & $83$ & $84$\\
$7$ & $2\,328$ & $1\,593$ & $32\%$ & $9\,358$ & $5\,342$ & $43\%$ & $109$ & $116$\\
$8$ & $3\,045$ & $2\,077$ & $32\%$ & $12\,271$ & $7\,019$ & $43\%$ & $125$ & $132$\\
$10$ & $4\,647$ & $3\,237$ & $30\%$ & $19\,021$ & $11\,081$ & $42\%$ & $151$ & $164$\\
$12$ & $6\,585$ & $4\,653$ & $29\%$ & $27\,239$ & $16\,059$ & $41\%$ & $177$ & $196$\\
\bottomrule
\end{tabular}
\end{center}

\noindent
At $n=12$: $6\,585$ Toffolis become $4\,653$, a $29\%$ saving, and depth falls
by $41\%$, for about $11\%$ more qubits.

\cite{ec:kim26} reports $58$--$60\%$ for division at $n=192$--$521$, but on a different metric:
the \emph{product} of qubit count and $T$-depth, which is the trade-off measure that paper uses
throughout. The columns above are plain counts, which is why they are not directly comparable ---
a construction that spends qubits to buy depth looks better under a product metric than under
either factor alone. The direction agrees; the magnitudes should not be read against each other.

The $2n$-bit budget is not a comfortable margin, it is the exact requirement --- $r$ reaches $6$
bits for $p=7$ and $32$ bits for $p=65521$, both exactly $2n$. The test measures it rather than
assuming it.

\subsection{Level 3, rebuilt: the Euclidean dialog}\label{ec:sec:dialog}

Everything so far treats inversion and multiplication as two jobs. Schrottenloher
\cite{ec:schrott26}, following \cite{ec:khattar25}, observes that they are one, and this is the
single largest structural saving in this part.

Run the binary GCD on $(u,v)$ alone and \emph{do not} maintain the B\'ezout coefficients. Just
record the two decision bits per iteration --- call the record a \emph{dialog}. It is $1.5$ bits
per iteration on average, against the $2n$ qubits carrying $(r,s)$ would cost:

\begin{center}\small
\begin{minipage}{0.82\textwidth}
\begin{tabbing}
xxx\=xxxx\=xxxxxxxxxxxxxxxxxxxxxxxxxxxxxxxxxxxx\=\kill
\> \textbf{repeat} $\approx1.413n$ times:\\
\> 1.\> $b_0 := v \bmod 2$; \quad $b_1 := [\,u>v\,]$ \> record $b_0$ and $b_0b_1$\\
\> 2.\> \textbf{if} $b_0b_1$: \ swap $(u,v)$\\
\> 3.\> \textbf{if} $b_0$: \ $v \mathrel{-}= u$\\
\> 4.\> $v \mathrel{/}= 2$
\end{tabbing}
\end{minipage}
\end{center}

\noindent
The iteration count is not arbitrary. Halving $v$ always removes a bit; the subtraction removes a
second one half the time. \cite{ec:schrott26} averages the shrink factor of $uv$ over those two
cases, $\tfrac12\cdot\tfrac12+\tfrac12\cdot\tfrac14=\tfrac38$, and reads off
$2n/\log_2(8/3)\approx1.413n$ iterations. That step is heuristic --- it averages the factor rather
than its logarithm, which is not the same thing --- but the conclusion holds: measured on this
implementation's own dialog, the mean iteration count is $1.383n$ at $n=32$ and $1.409n$ at
$n=256$, climbing towards $1.413n$. The spread is $\approx0.6\sqrt n$, which is what the $+c\sqrt n$
safety margin is for.

\begin{lstlisting}
def dialog_round(m, u, v, rec):
    """One iteration of [1128] Algorithm 2.  `rec` = (b0, b0b1), kept.
    ..."""
    ctx, n = m.ctx, len(u)
    b0, b0b1 = rec
    ctx.cx(v[0], b0[0])                              # b0 = v mod 2

    t, sc = m.anc(1, "t"), m.anc(n, "sc")
    A.gt_uint(ctx, u, v, t[0], sc)                   # t = [u > v]
    ctx.and_(b0[0], t[0], b0b1[0])                   # b0b1 = b0 AND b1
    A.gt_uint(ctx, u, v, t[0], sc)
    m.free(t, sc)

    for i in range(n):
        ctx.cswap(b0b1[0], u[i], v[i])               # if b0b1: swap u, v

    cp, sc = m.anc(n, "cp"), m.anc(n, "sc")
    A.csub(ctx, b0[0], u, v, cp, sc)                 # if b0: v -= u
    m.free(cp, sc)

    for i in range(n - 1):                           # v /= 2 (v is even now)
        ctx.swap(v[i], v[i + 1])
\end{lstlisting}

\noindent
Now the second idea. The updates to $(r,s)$ that the full extended algorithm would have performed
are all \emph{linear} and depend only on those two bits, so they can be replayed afterwards, onto
any starting pair:
\[
s \leftarrow 2s \bmod p; \qquad \text{if } b_0:\ s \mathrel{+}= r; \qquad \text{if } b_0b_1:\ \text{swap}(r,s)
\]
Read \emph{forwards}, alongside the Euclidean algorithm as the textbook extended version would,
this maps $(1,0)\mapsto(0,\,x^{-1})$ --- an inversion --- and more generally
$(y,0)\mapsto(0,\,yx^{-1})$, which is an inversion and a multiplication for the price of one pass.
That is \cite{ec:schrott26}'s observation, and it is already the saving.

But the circuit reads the record \emph{backwards}, and gets a different map:
$(y,0)\mapsto(0,\,yx)$. Multiplication by $x$, not by $x^{-1}$. That is not a defect but the
reason the backwards reading is chosen: replayed in reverse the updates need only \emph{doubling}
and addition, where the forward reading needs modular \emph{halving}. And since the whole thing is
a circuit, running it in the other direction recovers the division. One component, both of the
point addition's multiplicative steps.

The replay is seven lines, and the whole saving is that none of them looks at $x$:

\begin{lstlisting}
def bezout_replay(m, r, s, recs, q, dbl=None, cadd=None):
    """(r, s) updated by [1128] Algorithm 3, reading the record in reverse.
    ..."""
    dbl = dbl or (lambda mm, reg: MA.moddbl(mm, reg, q))
    cadd = cadd or (lambda mm, c, a, b: MA.cmodadd(mm, c, a, b, q))
    for b0, b0b1 in reversed(recs):
        dbl(m, s)                                        # s <- 2s mod q
        cadd(m, b0[0], r, s)                             # if b0: s += r mod q
        for i in range(len(r)):
            m.ctx.cswap(b0b1[0], r[i], s[i])             # if b0b1: swap r, s
\end{lstlisting}

\noindent
Assembling the three phases --- dialog, replay, dialog backwards --- gives the in-place
multiplication, and the rebuild consumes the record on its way out, so nothing is left behind:

\begin{lstlisting}
def inplace_mul(m, x, y, q, iters=None, dbl=None, cadd=None):
    """|x, y> -> |x, y*x mod q>.  [1128] Algorithm 4.

    ..."""
    n = len(x)
    iters = iters or eea_iterations(n)

    recs, u, v = dialog(m, x, q, iters)                  # 1: x -> dialog
    clear_dialog_end(m, u, v)

    s = m.anc(n, "bz")                                   # 2: replay onto (y, 0)
    bezout_replay(m, y, s, recs, q, dbl, cadd)
    for a, b in zip(y, s):                               # r is 0, s is the answer
        m.ctx.swap(a, b)
    m.free(s)

    hi = m.anc(1, "vhi")                                 # 3: dialog -> x
    u2 = m.anc(n + 1, "u")
    v2 = Reg(list(x) + list(hi), "v")
    m.ctx.x(u2[0])                                       # the dialog ends at u = 1
    for rec in reversed(recs):
        m.emit_inverse(dialog_round, m, u2, v2, rec)
    A.encode_const(m.ctx, u2, q)                         # and starts at u = q
    m.free(u2, hi)
    for rec in recs:
        m.free(*rec)
\end{lstlisting}

\begin{keybox}[title={Direction matters, and it is easy to get wrong}]
Applying Algorithm~3's update rules in the record's original order is \emph{not} the inverse of
applying them in reverse, and it is not the transpose either --- it is simply a different linear
map. Writing ``$(1,0)$ gives $x^{-1}$'' next to ``$(y,0)$ gives $yx$'' is self-contradictory, since
the second is linear in $y$ and contains the first at $y=1$; an earlier draft of this section said
exactly that, and the code is what caught it.

The classical model in \code{ec\_classical.py} carries all three maps and the tests pin down which
is which: \code{bezout\_replay} (reversed) multiplies by $x$, \code{bezout\_replay\_inv} divides by
it, and \code{bezout\_inverse\_classical} --- an independent coefficient-tracking implementation
that needs modular halving where the replay needs only doubling --- yields $x^{-1}$ directly. The
third exists purely as a cross-check, which is why the direction claim is not taken on faith.
\end{keybox}

\noindent
Compared against what it replaces --- Kaliski inversion plus an out-of-place multiplication, which
is what a point addition otherwise needs:

\begin{center}\small
\begin{tabular}{@{}rrrrrrr@{}}
\toprule
& \multicolumn{2}{c}{invert-then-multiply} & \multicolumn{2}{c}{dialog + replay} & \multicolumn{2}{c}{saving}\\
\cmidrule(lr){2-3}\cmidrule(lr){4-5}\cmidrule(lr){6-7}
$p$ & Toffoli & qubits & Toffoli & qubits & Toffoli & qubits\\
\midrule
$7$ & $1\,095$ & $64$ & $612$ & $39$ & $44\%$ & $39\%$\\
$11$ & $1\,910$ & $83$ & $957$ & $48$ & $50\%$ & $42\%$\\
$31$ & $2\,889$ & $99$ & $1\,378$ & $57$ & $52\%$ & $42\%$\\
$61$ & $4\,068$ & $115$ & $1\,875$ & $66$ & $54\%$ & $43\%$\\
\bottomrule
\end{tabular}
\end{center}

\noindent
At $p=61$ that is $54\%$ of the gates and $43\%$ of the qubits, and the
margin widens with $n$. The record itself is
squeezed further by a trick worth its own paragraph: each pair $(b_0, b_0b_1)$ can only be
$(0,0)$, $(1,0)$ or $(1,1)$ --- never $(0,1)$, since $b_0b_1$ cannot be set unless $b_0$ is --- so
three pairs carry $3^3=27<32$ states and fit in five bits instead of six. One qubit freed per three
iterations, taking the record from $2.83n$ to $2.355n$ asymptotically, and (unlike a shift-based
encoding) leaving it a \emph{fixed} size, which is what lets the rest of the circuit be laid out
statically.

\subsection{Level 4, rebuilt: projective coordinates}\label{ec:sec:proj}

The other way to make an inversion cheap is not to do it. In Jacobian coordinates
$(X:Y:Z)\sim(x,y)=(X/Z^2,\,Y/Z^3)$ the addition law is inversion-free, and with the addend in
affine form ($Z_2=1$, which it is, being a classical point) every term involving $Z_2$ drops out:
\[
\begin{aligned}
U_2 &= x_2Z_1^2, & S_2 &= y_2Z_1^3, & H &= U_2-X_1, & R &= S_2-Y_1,\\
X_3 &= R^2-H^3-2X_1H^2, & Y_3 &= R(X_1H^2-X_3)-Y_1H^3, & Z_3 &= Z_1H. &&
\end{aligned}
\]
Eleven multiplications, no division --- \cite{ec:kim26} reports this as the fewest of any quantum
projective addition, against a previous best of twelve.

The catch is that projective coordinates are not a unique representation. $(P+Q)-Q$ returns a
\emph{different triple} for the same point, so the input cannot be erased and every addition leaves
its accumulator behind. The fifteen intermediates \emph{can} all be cleaned --- each by reversing
the step that made it, with that step's inputs still live --- and this implementation cleans them,
but the input point is $3n$ qubits of irreducible garbage per addition.

The eighteen steps, and then the unwinding --- note the last two lines, which are the whole
subtlety:

\begin{lstlisting}
def jacobian_add(m, X1, Y1, Z1, x2, y2, X3, Y3, Z3, p):
    """(X3:Y3:Z3) <- (X1:Y1:Z1) + (x2, y2).  [106] Algorithm 4.
    ..."""
    n = len(X1)
    A = m.anc

    ...

    Z1sq = A(n, "Z1sq"); _, b1 = m.step(sqr, Z1, Z1sq)                    # 1
    U2 = A(n, "U2");     _, b2 = m.step(MU.modmul_const, m, Z1sq, x2, U2, p)   # 2
    Z1cu = A(n, "Z1cu"); _, b3 = m.step(mul, Z1sq, Z1, Z1cu)              # 3
    H = A(n, "H");       _, b4 = m.step(sub_into, U2, X1, H)              # 4
    ...
    def mk2V():
        for q, r in zip(V, twoV):
            m.ctx.cx(q, r)
        MA.modadd(m, V, twoV, p)
    _, b11 = m.step(mk2V)                                                 # 11

    T0 = A(n, "T0");     _, b12 = m.step(sub_into, Rsq, Hcu, T0)          # 12
    _, b13 = m.step(sub_into, T0, twoV, X3)                        # 13  OUTPUT
    T1 = A(n, "T1");     _, b14 = m.step(sub_into, V, X3, T1)             # 14
    T2 = A(n, "T2");     _, b15 = m.step(mul, R, T1, T2)                  # 15
    T3 = A(n, "T3");     _, b16 = m.step(mul, Y1, Hcu, T3)                # 16
    _, b17 = m.step(sub_into, T2, T3, Y3)                          # 17  OUTPUT
    _, b18 = m.step(mul, Z1, H, Z3)                                # 18  OUTPUT

    # unwind everything but the three outputs.  Note T0 goes back via step 12,
    # not step 13 -- step 13 wrote X3, which stays.
    for body in (b16, b15, b14, b12, b11, b10, b9, b8, b7, b6, b5, b4, b3, b2, b1):
        m.undo(body)
    m.free(T3, T2, T1, T0, twoV, V, Rsq, R, S2, Hcu, Hsq, H, Z1cu, U2, Z1sq)
\end{lstlisting}

\begin{warnbox}[title={The order of unwinding is not simply ``backwards''}]
$X_3$ is an output and must survive, so the intermediate $T_0$ cannot be released by reversing the
step that consumed it. It is released by reversing the step that \emph{produced} it, from $R^2$ and
$H^3$, which are still live at that moment. Getting this wrong destroys the answer;
\code{Machine.undo} exists precisely so that a recorded step can be reversed out of order, and the
tests sweep the projective representatives $Z$ of every input point --- $Z\in\{1,2\}$ by default
and all of $[1,p)$ under \code{SHOR\_EC\_FULL=1} --- a sweep that catches representation bugs an
affine-only test cannot see.
\end{warnbox}

\subsection{Level 5: windowing, zig-zag, and the semiclassical transform}\label{ec:sec:window}

Three counting arguments, none of which changes a single arithmetic circuit.

\emph{Windowing} \cite{ec:gidney-ekera} groups $w$ control qubits into one address and looks up
$P_i$ from a table of $2^w$ precomputed points, so $2n+2$ additions become $2\lceil (n{+}1)/w\rceil$.
A lookup costs $2^w$ Toffolis, an order of magnitude below an addition, so $w$ is chosen large ---
both \cite{ec:schrott26} and \cite{ec:babbush26} use $w=16$.

\emph{Zig-zag} \cite{ec:grassl16} attacks the projective garbage. Write $\rho$ for the number of
accumulator registers --- not $m$, which above is the counting-register width. Additions are
arranged in runs of $\rho, \rho-1, \dots, 1$: fill all $\rho$ registers, then walk back
un-computing all but the last of the run --- legal, because each state's predecessor is still live
--- returning $\rho-1$ registers, and repeat. Total forward additions $\rho(\rho+1)/2$, so $\rho$
registers serve them all.

\emph{The semiclassical transform} \cite{ec:griffiths-niu} is the same trick as
\S\ref{sec:onectrl}: every controlled rotation in the inverse QFT is controlled by a qubit about to
be measured, so measure it first and make the control classical. Both scalar registers collapse to
\emph{one recycled qubit}. This is the assumption behind both papers reporting only the cost of the
point arithmetic.

The windowed addition itself, with the three lookups and the two in-place multiplications that
replace \S\ref{ec:sec:padd}'s two divisions:

\begin{lstlisting}
def windowed_point_add(m, addr, x2, y2, points, p, iters=None):
    """[1128] Algorithm 1.  `points[i]` is the window entry for address i.
    ..."""
    ctx, n = m.ctx, len(x2)
    w = len(addr)
    assert len(points) == 1 << w, "the window must be full"
    assert points[0].inf, "entry 0 must be the point at infinity"

    ...

    one = m.anc(1, "one")
    ctx.x(one[0])
    c = m.anc(1, "c")
    MA.is_nonzero(m, addr, c[0])                       # 1

    x1, y1 = m.anc(n, "x1"), m.anc(n, "y1")
    _lookup(m, one[0], addr, x1, xs)                   # 2
    _lookup(m, one[0], addr, y1, ys)
    MA.modsub(m, x1, x2, p)                            # 3
    MA.modsub(m, y1, y2, p)                            # 4
    _lookup(m, one[0], addr, x1, xs)                   # 5
    _lookup(m, one[0], addr, y1, ys)

    E.inplace_div(m, x2, y2, p, iters)                 # 6  y2 <- y2 / x2

    xx = m.anc(n, "3x")
    _lookup(m, one[0], addr, xx, x3s)                  # 7
    MA.modadd(m, xx, x2, p)                            # 8
    _lookup(m, one[0], addr, xx, x3s)                  # 9
    m.free(xx)

    _csub_square(m, c[0], y2, x2, p)                   # 10 if c: x2 -= y2^2
    E.inplace_mul(m, x2, y2, p, iters)                 # 11 y2 <- y2 * x2
    MA.cmodneg(m, c[0], x2, p)                         # 12 if c: x2 <- -x2
    ...
\end{lstlisting}

\begin{center}\small
\begin{tabular}{@{}rrrrrrrrr@{}}
\toprule
& \multicolumn{3}{c}{point additions} & \multicolumn{3}{c}{zig-zag registers} & \multicolumn{2}{c}{garbage qubits}\\
\cmidrule(lr){2-4}\cmidrule(lr){5-7}\cmidrule(lr){8-9}
$n$ & none & $w{=}11$ & $w{=}16$ & none & $w{=}11$ & $w{=}16$ & none & $w{=}16$\\
\midrule
$192$ & $386$ & $36$ & $26$ & $28$ & $8$ & $7$ & $16\,128$ & $4\,032$\\
$256$ & $514$ & $48$ & $34$ & $32$ & $10$ & $8$ & $24\,576$ & $6\,144$\\
$384$ & $770$ & $70$ & $50$ & $39$ & $12$ & $10$ & $44\,928$ & $11\,520$\\
$521$ & $1\,044$ & $96$ & $66$ & $46$ & $14$ & $11$ & $71\,898$ & $17\,193$\\
\bottomrule
\end{tabular}
\end{center}

\noindent
At $n=256$: $514$ additions become $34$, a $15\times$ reduction,
and the registers holding projective garbage fall from $32$ to $8$. The two
compose --- windowing cuts how many additions there are, zig-zag cuts how many of their outputs are
live at once --- and together they take the garbage from $24\,576$ qubits to $6\,144$.

\section{Everything at once}\label{ec:sec:cost}

The subsections above each hold one thing fixed. This one does not: it stacks them, on the operation
that actually matters --- one point addition --- and then on the whole circuit.

\begin{center}\small
\begin{tabular}{@{}llrrrr@{}}
\toprule
curve & construction & Toffoli & qubits & depth & vs.\ affine\\
\midrule
\multicolumn{6}{@{}l}{\footnotesize Alg.~3, Alg.~4 are \cite{ec:kim26}'s; Alg.~1 is \cite{ec:schrott26}'s.}\\
$p=7$ & affine, in place (Alg.~3) & $2\,756$ & $64$ & $11\,890$ & ---\\
 & Jacobian, out of place (Alg.~4) & $1\,544$ & $75$ & $8\,511$ & $44\%$\\
$p=11$ & affine, in place (Alg.~3) & $4\,769$ & $83$ & $21\,024$ & ---\\
 & Jacobian, out of place (Alg.~4) & $2\,642$ & $98$ & $15\,126$ & $45\%$\\
$p=7$ & windowed $w{=}2$, dialog mult.\ (Alg.~1) & $1\,523$ & $49$ & $6\,893$ & $45\%$\\
\bottomrule
\end{tabular}
\end{center}

\noindent
Three constructions of the same map, and they stack differently. The affine addition is in place and
leaves nothing behind, and pays for that with an inversion. The Jacobian addition removes the
inversion entirely and is $45\%$ cheaper, but leaves $3n$ qubits of input point as
garbage --- which is only affordable once windowing and zig-zag have cut how many such inputs are
live at once. The windowed addition is cheapest of the three at $45\%$ below affine,
and it gets there differently again: not by changing coordinates, but by replacing
invert-then-multiply with a single Euclidean dialog, and by amortising three table lookups over a
whole window of additions.

\begin{keybox}[title={How the savings compose --- and where they do not}]
Read down the stack for one point addition at $n=256$, taking each factor from the measurements
above:
\begin{itemize}\setlength\itemsep{2pt}
\item \emph{Gate-level, and only partly independent.} The Gidney adder ($\approx2\times$ on the
inner additions) and the temporary AND ($2.50\times$ on $T$) are \emph{not} two savings:
they are one saving counted at two levels. The adder wins by choosing a gate whose target starts
clean, and the temporary-AND accounting is exactly what makes that gate's uncompute free.
Multiplying the two would double-count. Copy-then-add controls are a smaller further factor than
the table suggests: the $2.0\times$ there bundles the control style together with the
switch to Gidney's adder inside it. Hold the adder fixed and the control style alone is worth
$1.33\times$ ($96$ against $128$ at $n=32$, CDKM on both sides) --- it
buys parallelism and depth more than gates. Pseudo-Mersenne reduction
($6.2\times$ on doublings) is the one that cleanly composes, because it acts on a
different part of the circuit again: the reductions rather than the additions.
\item \emph{Structural, and they replace rather than compose.} The dialog
($54\%$ off the multiplicative step) and projective coordinates
($45\%$ off the addition) are two different answers to the same problem --- the
inversion --- so you take one, not both. Every published construction picks a side.
\item \emph{Count-level, and these do multiply the whole thing.} Windowing divides the number of
point additions by $15$ at $n=256$; zig-zag divides the live projective garbage by
$4$; the semiclassical transform deletes both counting registers. None of them touches an
arithmetic circuit, which is precisely why they compose cleanly with everything above.
\end{itemize}
The last group dominates. An optimisation that halves a point addition is worth half as much as one
that runs fifteen times fewer of them, which is why both papers spend their first section on
windowing and their last on gate counts.
\end{keybox}

\noindent
And the whole circuit, on the toy curve $y^2=x^3+5x+4 \bmod 7$ with $\mathrm{ord}(G)=5$:

\begin{center}\small
\begin{tabular}{@{}lrrrr@{}}
\toprule
ECDLP circuit & qubits & Toffoli & gates & depth\\
\midrule
real arithmetic, two counting registers & $69$ & $16\,488$ & $105\,339$ & $70\,991$\\
permutation oracle, two counting registers & $12$ & --- & $603$ & $282$\\
permutation oracle, semiclassical QFT & $7$ & --- & $619$ & $287$\\
\bottomrule
\end{tabular}
\end{center}

\noindent
$69$ qubits and $105\,339$ gates to compute a discrete logarithm in a group of order
five. That ratio is the honest summary of this part: the arithmetic is enormous, and the
counting-level tricks are the only reason cryptographic sizes are discussed at all. Note the last
two rows --- the semiclassical transform takes $12$ qubits to $7$ by recycling one
control qubit, which is precisely why both papers report only the point arithmetic and treat the
scalar registers as free.

\subsection{Space at cryptographic size}\label{ec:sec:space}

The one place where a projection is worth making, because it is what the papers' headline qubit
counts are built from. The dialog record dominates the space of \cite{ec:schrott26}'s construction:

\begin{center}\small
\begin{tabular}{@{}rrrrr@{}}
\toprule
$n$ & record, raw & with Fig.~1 packing & per $n$ & with register sharing\\
\midrule
$256$ & $802$ & $669$ & $2.613n$ & $579$\\
$384$ & $1\,180$ & $984$ & $2.562n$ & $859$\\
$521$ & $1\,582$ & $1\,319$ & $2.532n$ & $1\,157$\\
\bottomrule
\end{tabular}
\end{center}

\noindent
The middle column converges to the paper's $2.355n$ from above --- the gap is the $O(\sqrt n)$
safety margin on the iteration count, which is a constant fraction at these sizes and vanishes
asymptotically. The last column is \emph{not} implemented here: it is Sec 3.1's register sharing, in
which $u$ and $v$ shrink as the algorithm proceeds and the freed qubits absorb the record. It is
priced here but not built --- \S\ref{ec:sec:gaps} says why.

The paper's own closed form for the total --- its asymptotic $2.355n$ record plus $2n$ for the pair
of modular integers the reconstruction holds, i.e.\ $4.355n+O(\sqrt n)$ --- gives
$1151$ qubits at $n=256$ against the paper's $1192$, and $1407$ against $1446$
for the gate-optimised variant. That is the paper's formula evaluated, \emph{not} the table above
summed: the table's middle column carries its $O(\sqrt n)$ margin explicitly and would give a
larger figure. Neither is a check on the other.

The full-algorithm count of $2^{26.11}$ Toffolis is likewise the paper's equation~(1)
evaluated at the paper's own per-addition figure --- restated for scale, not independently
confirmed. Nothing in this package measures a $256$-bit point addition.

\section{Verifying it}\label{ec:sec:verify}

A wrong ECDLP circuit does not return a wrong answer. It returns a \emph{flat histogram} --- and
since the classical post-processing verifies its own candidates against the curve, it will happily
print the correct $k$ from pure noise. Testing the output is therefore worth almost nothing; three
other things are worth a great deal.

\paragraph{Exact simulation at full width.}
Every arithmetic circuit here is a permutation of basis states, so pushing one basis state through
in $O(\text{gates})$ time is exact --- and scales to hundreds of qubits where a statevector dies at
thirty. Ancillas are checked back to $\ket0$ \emph{at the instruction where they are freed}, not
merely at the end, which is what catches a circuit that computes the right value and leaves a dirty
flag behind. This is sound only for circuits whose phases carry no information; the sole phase and
Hadamard gates outside the AND gadget are in the state preparation and the Fourier transform, which
are run on Aer instead.

\begin{lstlisting}
def simulate(qc, init=None, checks=(), strict=True):
    """Push one basis state through `qc` and return the final bit assignment.

    ..."""
    idx = {q: i for i, q in enumerate(qc.qubits)}
    bits = [0] * qc.num_qubits
    if init:
        for k, v in init.items():
            bits[idx[k] if k in idx else k] = int(v) & 1

    ...
        name = op.name
        w = [idx[q] for q in qargs]

        if name in _SKIP or name in _DIAGONAL:
            return
        if name == "x":
            bits[w[0]] ^= 1
            return
        if name == "swap":
            bits[w[0]], bits[w[1]] = bits[w[1]], bits[w[0]]
            return
        if name == "ecand":
            if strict and bits[w[2]] != 0:
                raise SimError("AND target was not |0>: temporary AND invalid")
            bits[w[2]] = bits[w[0]] & bits[w[1]]
            return
        if name == "ecand_dg":
            if strict and bits[w[2]] != (bits[w[0]] & bits[w[1]]):
                raise SimError("AND-dagger target did not hold a AND b")
            bits[w[2]] = 0
            ...
\end{lstlisting}

\paragraph{Statevector cross-checks.}
The fast simulator only sees basis states, and a bug in it would hide from every test that uses it.
So the small circuits are also evolved on genuine superpositions and checked for exact fidelity
\emph{and} for the absence of leaked amplitude --- the statement that ancillas are unentangled,
which basis-state testing cannot see. A stray diagonal gate would appear here as a phase, and phases
are exactly what Shor depends on.

\paragraph{The listings, too.}
Every listing in this part is extracted from \code{shor\_qiskit/ec\_*.py} by AST at build time, not
copied by hand, and the whole part is regenerated and diffed against the document by
\code{bench/ec\_check\_tex.py}. So a listing cannot drift from the function it claims to show:
change the code and the document either follows or the check fails. Where a listing is shortened,
the elision is visible as a bare \code{...} --- nothing is silently trimmed.

\paragraph{Property-based tests over random curves.}
Random prime, random non-singular $(a,b)$, random points and scalars. That is where curve-shape
assumptions hide: $a=0$, $b=0$, curves with a point of order two, primes where $2$ is or is not a
square. The harness shrinks failures --- staying inside each field's declared domain, and
re-checking its own result before reporting it --- and it has been validated against injected bugs,
where it isolates the exact offending modulus.

\begin{center}\small
\begin{tabular}{@{}ll@{}}
\toprule
what & result\\
\midrule
classical models (Kaliski, Montgomery, dialog, Jacobian) & exact\\
adders, comparators, modular arithmetic & exhaustive $p\le61$, random to $n=128$\\
Kaliski inversion and division & exhaustive; records and counter clean\\
affine point addition & every point pair, both control branches\\
projective point addition & every point pair; every $Z$ under \code{SHOR\_EC\_FULL}\\
in-place multiplication (dialog) & every $(x,y)$ for $p\le31$ ($p\le61$ full)\\
ECDLP end to end & every logarithm on two toy curves\\
\bottomrule
\end{tabular}
\end{center}

\subsection{The distribution, and the control that makes it mean something}\label{ec:sec:post}

On $y^2=x^3+5x+4\bmod7$ with $G=(0,5)$ of order $5$, and $Q=[2]G$, the measured peaks sit at
$(j_1,j_2)\in\{(0,0),(6,5),(2,3),(3,6),(5,2)\}$ carrying $20\%$, $10\%$, $10\%$, $10\%$, $10\%$ of
the shots. Rounding by $a_i=\mathrm{round}(j_i\cdot5/8)$ turns those into $(a_1,a_2) =
(1,2),(2,4),(3,1),(4,3)$ --- and every one of them satisfies $a_2=2a_1 \bmod 5$, which is
\eqref{ec:eq:post} with $k=2$.

\begin{lstlisting}
def ecdlp_postprocess(counts, order, bits, search=1):
    """Recover k from measured (j1, j2) pairs of the two ECDLP registers.

    `counts` maps (j1, j2) -> number of shots.  Returns [(k, weight)] sorted by
    weight, heaviest first.
    ..."""
    if isinstance(counts, dict):
        items = list(counts.items())
    else:
        items = [(pair, 1) for pair in counts]

    votes, q = {}, 1 << bits
    for (j1, j2), c in items:
        for da in range(-search, search + 1):
            for db in range(-search, search + 1):
                a1 = (round(j1 * order / q) + da) % order
                a2 = (round(j2 * order / q) + db) % order
                if gcd(a1, order) != 1:
                    continue
                k = a2 * pow(a1, -1, order) % order
                w = c / (1 + abs(da) + abs(db)) ** 2
                votes[k] = votes.get(k, 0.0) + w
    return sorted(votes.items(), key=lambda kv: -kv[1])
\end{lstlisting}

\begin{warnbox}[title={The control that catches a broken circuit}]
The correct $k$ must be the \emph{top} candidate, not merely \emph{a} candidate that verifies. On
both toy curves it is, for every logarithm, holding at least $30\%$ of the vote against a uniform
$20\%$ and $7.7\%$ respectively. And feeding the same post-processing a flat, random histogram
returns a top share of $21.2\%$ against a uniform $20\%$ --- no peak, no signal.

That control exists because an early version of this post-processing did not have it. Eker\aa{}'s
limited search \cite{ec:ekera} widens the rounding when $2^m$ is not a multiple of $r$, and with the
radius set too generously it enumerated the entire group: every $k$ ``verified'', on every input,
including noise. The bug was invisible in the pass/fail of an end-to-end test and obvious the
moment the vote shares were printed next to the uniform baseline.
\end{warnbox}

\section{What is not here}\label{ec:sec:gaps}

Named rather than glossed:

\begin{itemize}\setlength\itemsep{2pt}
\item \textbf{Register sharing} (\cite{ec:schrott26} Sec 3.1), which takes the dialog space from
$2.83n$ to $2.12n$ by shrinking $u$ and $v$ as the algorithm proceeds. It is left out not for being
probabilistic --- \S\ref{ec:sec:approx} builds and measures several circuits that are --- but
because its failure mode is a \emph{register overflow} rather than a wrong answer, so the
exact-simulation harness of \S\ref{ec:sec:verify} would be checking a circuit whose register layout
is itself only probably right. Priced in \S\ref{ec:sec:space}, not built.
\item \textbf{Algorithm 11} of \cite{ec:schrott26}, the pseudo-Mersenne controlled addition that
also handles $x+y=p$. That case cannot arise from random inputs, only in the first few
B\'ezout-replay iterations, and using the exact adder there is the same fix at negligible cost.
\item \textbf{Gidney's dirty-ancilla constant adder} \cite{ec:gidney25}. Constant addition here uses
clean ancillas: correct, and more qubits than the papers' accounting.
\item \textbf{The Ed25519 extension} of \cite{ec:kim26} Sec 4.4, in extended Edwards coordinates.
The curve model here is short Weierstrass only.
\item \textbf{Physical-level estimates} --- surface codes, QLDPC, wall-clock runtime. This part
stops at the logical layer, as Parts II--VI do.
\end{itemize}

\noindent
And one thing implemented in only one of its two forms: \cite{ec:kim26} distinguishes a
\emph{depth-optimised} Montgomery multiplier, which allocates fresh output qubits per two-operand
addition, from a \emph{qubit-optimised} one, which uncomputes and reuses them. The implementation
here always accumulates in place --- the qubit-optimised shape --- which is why the depth win from
the carry-save tree shows up smaller than the paper's depth-optimised figures.

\begin{keybox}[title={Where this part sits}]
The affine construction is \cite{ec:rnsl17} as improved by \cite{ec:hjnrs20}, with
\cite{ec:kim26}'s modifications: 2017--2020 with a 2026 finish. The projective addition and the
carry-save Montgomery multiplier are \cite{ec:kim26}. The Euclidean dialog, the approximate and
pseudo-Mersenne arithmetic, and the windowed addition are \cite{ec:schrott26}. What is deliberately
absent is the line that leaves this architecture entirely --- the residue-number-system
constructions that cut qubit counts by changing what a number \emph{is}, rather than how it is
added.

The honest limit of what runs on a laptop: the circuit with real arithmetic is verified on every
basis state, which for a permutation is the whole truth, but at $69$ qubits for a
seven-element field it cannot be run in superposition. The end-to-end demonstration uses a
permutation oracle of identical circuit shape. Both halves are tested; the join between them is an
argument, not a simulation.
\end{keybox}

\begin{thebibliography}{99}\small

\bibitem{shor}
P.~W. Shor,
\emph{Polynomial-time algorithms for prime factorization and discrete logarithms on a quantum
computer}, SIAM Review \textbf{41}(2):303--332, 1999 (SIAM J.\ Comput.\ \textbf{26}(5):1484--1509,
1997). \href{https://arxiv.org/abs/quant-ph/9508027}{arXiv:quant-ph/9508027}.

\bibitem{kitaev}
A.~Yu. Kitaev,
\emph{Quantum measurements and the Abelian stabilizer problem},
1995. \href{https://arxiv.org/abs/quant-ph/9511026}{arXiv:quant-ph/9511026}.

\bibitem{cemm}
R.~Cleve, A.~Ekert, C.~Macchiavello, and M.~Mosca,
\emph{Quantum algorithms revisited},
Proc.\ R.\ Soc.\ Lond.\ A \textbf{454}:339--354, 1998.
\href{https://arxiv.org/abs/quant-ph/9708016}{arXiv:quant-ph/9708016}.

\bibitem{vbe}
V.~Vedral, A.~Barenco, and A.~Ekert,
\emph{Quantum networks for elementary arithmetic operations},
Phys.\ Rev.\ A \textbf{54}:147--153, 1996.
\href{https://arxiv.org/abs/quant-ph/9511018}{arXiv:quant-ph/9511018}.

\bibitem{bcdp}
D.~Beckman, A.~N. Chari, S.~Devabhaktuni, and J.~Preskill,
\emph{Efficient networks for quantum factoring},
Phys.\ Rev.\ A \textbf{54}:1034--1063, 1996.
\href{https://arxiv.org/abs/quant-ph/9602016}{arXiv:quant-ph/9602016}.

\bibitem{zalka}
C.~Zalka,
\emph{Fast versions of Shor's quantum factoring algorithm},
1998. \href{https://arxiv.org/abs/quant-ph/9806084}{arXiv:quant-ph/9806084}.

\bibitem{draper}
T.~G. Draper,
\emph{Addition on a quantum computer},
2000. \href{https://arxiv.org/abs/quant-ph/0008033}{arXiv:quant-ph/0008033}.

\bibitem{beauregard}
S.~Beauregard,
\emph{Circuit for Shor's algorithm using $2n+3$ qubits},
Quantum Inf.\ Comput.\ \textbf{3}(2):175--185, 2003.
\href{https://arxiv.org/abs/quant-ph/0205095}{arXiv:quant-ph/0205095}.

\bibitem{cuccaro}
S.~A. Cuccaro, T.~G. Draper, S.~A. Kutin, and D.~P. Moulton,
\emph{A new quantum ripple-carry addition circuit},
2004. \href{https://arxiv.org/abs/quant-ph/0410184}{arXiv:quant-ph/0410184}.

\bibitem{pavlidis}
A.~Pavlidis and D.~Gizopoulos,
\emph{Fast quantum modular exponentiation architecture for Shor's factoring algorithm},
Quantum Inf.\ Comput.\ \textbf{14}(7--8):649--682, 2014.
\href{https://arxiv.org/abs/1207.0511}{arXiv:1207.0511}.

\bibitem{barenco-aqft}
A.~Barenco, A.~Ekert, K.-A. Suominen, and P.~T\"orm\"a,
\emph{Approximate quantum Fourier transform and decoherence},
Phys.\ Rev.\ A \textbf{54}:139--146, 1996.
\href{https://arxiv.org/abs/quant-ph/9601018}{arXiv:quant-ph/9601018}.

\bibitem{cheung}
D.~Cheung,
\emph{Improved bounds for the approximate QFT},
2004. \href{https://arxiv.org/abs/quant-ph/0403071}{arXiv:quant-ph/0403071}.

\bibitem{griffiths-niu}
R.~B. Griffiths and C.-S. Niu,
\emph{Semiclassical Fourier transform for quantum computation},
Phys.\ Rev.\ Lett.\ \textbf{76}:3228--3231, 1996.
\href{https://arxiv.org/abs/quant-ph/9511007}{arXiv:quant-ph/9511007}.

\bibitem{mosca-ekert}
M.~Mosca and A.~Ekert,
\emph{The hidden subgroup problem and eigenvalue estimation on a quantum computer},
Proc.\ QCQC'98, LNCS 1509, 1999.
\href{https://arxiv.org/abs/quant-ph/9903071}{arXiv:quant-ph/9903071}.

\bibitem{parker-plenio}
S.~Parker and M.~B. Plenio,
\emph{Efficient factorization with a single pure qubit and $\log N$ mixed qubits},
Phys.\ Rev.\ Lett.\ \textbf{85}:3049--3052, 2000.
\href{https://arxiv.org/abs/quant-ph/0001066}{arXiv:quant-ph/0001066}.

\bibitem{gidney-and}
C.~Gidney,
\emph{Halving the cost of quantum addition},
Quantum \textbf{2}:74, 2018.
\href{https://arxiv.org/abs/1709.06648}{arXiv:1709.06648}.

\bibitem{babbush} R.~Babbush, C.~Gidney, D.~W. Berry, N.~Wiebe, J.~McClean, A.~Paler,
A.~Fowler, and H.~Neven, \emph{Encoding electronic spectra in quantum circuits with linear
$T$ complexity}, Phys.\ Rev.\ X \textbf{8}:041015, 2018.
\href{https://arxiv.org/abs/1805.03662}{arXiv:1805.03662}.
\bibitem{zalka-coset} C.~Zalka, \emph{Shor's algorithm with fewer (pure) qubits}, 2006.
\href{https://arxiv.org/abs/quant-ph/0601097}{arXiv:quant-ph/0601097}. --- the coset
representation of modular integers.
\bibitem{gidney-ekera} C.~Gidney and M.~Eker\aa, \emph{How to factor 2048 bit RSA integers
in 8 hours using 20 million noisy qubits}, Quantum \textbf{5}:433, 2021.
\href{https://arxiv.org/abs/1905.09749}{arXiv:1905.09749}.
\bibitem{gidney-window}
C.~Gidney,
\emph{Windowed quantum arithmetic},
2019. \href{https://arxiv.org/abs/1905.07682}{arXiv:1905.07682}.

\bibitem{nielsen-chuang}
M.~A. Nielsen and I.~L. Chuang,
\emph{Quantum Computation and Quantum Information},
10th Anniversary Ed., Cambridge University Press, 2010.

\bibitem{vandersypen}
L.~M.~K. Vandersypen, M.~Steffen, G.~Breyta, C.~S. Yannoni, M.~H. Sherwood, and I.~L. Chuang,
\emph{Experimental realization of Shor's quantum factoring algorithm using nuclear magnetic
resonance},
Nature \textbf{414}:883--887, 2001.

\bibitem{notebook15}
\emph{Building Shor's algorithm from scratch: factoring 15},
Jupyter notebook (Qiskit), after \cite{draper,beauregard}, Nielsen--Chuang
\cite{nielsen-chuang}, and the Qiskit community tutorials.

\bibitem{repo}
B.~Assel,
\emph{qiskit-shor: full implementation of Shor's factoring algorithm with the Qiskit SDK},
GitHub repository, 2025--2026.
\url{https://github.com/benjamin-assel/qiskit-shor}.

\bibitem{blog}
B.~Assel,
\emph{Where are we with Shor's algorithm?},
Towards Data Science, 2025.
\url{https://towardsdatascience.com/where-are-we-with-shors-algorithm/}.


\bibitem{ec:proos-zalka}
J.~Proos and C.~Zalka,
\emph{Shor's discrete logarithm quantum algorithm for elliptic curves},
Quantum Inf.\ Comput.\ \textbf{3}(4):317--344, 2003.

\bibitem{ec:rnsl17}
M.~Roetteler, M.~Naehrig, K.~M. Svore, and K.~Lauter,
\emph{Quantum resource estimates for computing elliptic curve discrete logarithms},
ASIACRYPT 2017, LNCS 10625, pp.~241--270.

\bibitem{ec:hjnrs20}
T.~H\"aner, S.~Jaques, M.~Naehrig, M.~Roetteler, and M.~Soeken,
\emph{Improved quantum circuits for elliptic curve discrete logarithms},
PQCrypto 2020, LNCS 12100, pp.~425--444.

\bibitem{ec:kim26}
H.~Kim, K.~Jang, S.~Wang, V.~Srivastava, A.~Baksi, G.~Song, H.~Seo, and A.~Chattopadhyay,
\emph{New quantum circuits for ECDLP: breaking prime elliptic curve cryptography},
IACR ePrint 2026/106.
\url{https://eprint.iacr.org/2026/106}.

\bibitem{ec:schrott26}
A.~Schrottenloher,
\emph{Optimized point addition circuits for elliptic curve discrete logarithms},
IACR ePrint 2026/1128.
\url{https://eprint.iacr.org/2026/1128}.

\bibitem{ec:babbush26}
R.~Babbush, A.~Zalcman, C.~Gidney, M.~Broughton, T.~Khattar, H.~Neven, T.~Bergamaschi, J.~Drake,
and D.~Boneh,
\emph{Securing elliptic curve cryptocurrencies against quantum vulnerabilities: resource estimates
and mitigations},
arXiv:2603.28846, 2026.

\bibitem{ec:khattar25}
T.~Khattar, N.~Shutty, C.~Gidney, A.~Zalcman, N.~Yosri, D.~Maslov, R.~Babbush, and S.~P. Jordan,
\emph{Verifiable quantum advantage via optimized DQI circuits},
arXiv:2510.10967, 2025.

\bibitem{ec:ekera}
M.~Eker\aa{},
\emph{Revisiting Shor's quantum algorithm for computing general discrete logarithms},
arXiv:1905.09084, 2019.

\bibitem{ec:gidney25}
C.~Gidney,
\emph{A classical-quantum adder with constant workspace and linear gates},
arXiv:2507.23079, 2025.

\bibitem{ec:grassl16}
M.~Grassl, B.~Langenberg, M.~Roetteler, and R.~Steinwandt,
\emph{Applying Grover's algorithm to AES: quantum resource estimates},
PQCrypto 2016, LNCS 9606, pp.~29--43.

\bibitem{ec:classiq}
Classiq,
\emph{Elliptic curve discrete logarithm} (tutorial notebook), 2025--2026.
\url{https://docs.classiq.io/explore/algorithms/number_theory_and_cryptography/elliptic_curves/elliptic_curve_discrete_log}.

\end{thebibliography}


\end{document}
